% Môn: Vật lý
% Lớp: 12
% Chương: Chương III. Từ trường
% Bài: L12C3 Bài 18. Ứng dụng hiện tượng cảm ứng điện từ
% Số câu: 324
% App đọc file này trực tiếp — sửa rồi Commit trên GitHub, bấm Đồng bộ GitHub .tex

% L12C3 Bài 18. Ứng dụng hiện tượng cảm ứng điện từ

% ===== Câu 1 =====
% ID: L12C3B18-01-DS
% Mức: TH
\dangbt{Bài toán tổng hợp về máy biến áp, truyền tải điện và an toàn điện}
\begin{ex}
% ID: L12C3B18-01-DS
% Mức: TH
Xét Bài toán tổng hợp về máy biến áp, truyền tải điện và an toàn điện (Tình huống 1). Hãy xác định tính đúng, sai của bốn phát biểu sau.
\choiceTF
{\True An toàn điện yêu cầu cách điện, nối đất phù hợp và không chạm phần dẫn điện đang có điện áp nguy hiểm.}
{Có thể bỏ qua nối đất và thiết bị bảo vệ khi dùng điện áp cao.}
{\True Khi giải cần xác định đúng dữ kiện, phương chiều, điều kiện áp dụng và đơn vị.}
{Kết quả không phụ thuộc dữ kiện và có thể bỏ qua điều kiện.}
\loigiai{
\begin{itemchoice}
\item (Đúng) An toàn điện yêu cầu cách điện, nối đất phù hợp và không chạm phần dẫn điện đang có điện áp nguy hiểm. (Vì): An toàn điện yêu cầu cách điện, nối đất phù hợp và không chạm phần dẫn điện đang có điện áp nguy hiểm. b) Sai: Có thể bỏ qua nối đất và thiết bị bảo vệ khi dùng điện áp cao. c) Đúng: Khi giải cần xác định đúng dữ kiện, phương chiều, điều kiện áp dụng và đơn vị. d) Sai: Kết quả không phụ thuộc dữ kiện và có thể bỏ qua điều kiện. Kiến thức cốt lõi: An toàn điện yêu cầu cách điện, nối đất phù hợp và không chạm phần dẫn điện đang có điện áp nguy hiểm
\item (Sai) Có thể bỏ qua nối đất và thiết bị bảo vệ khi dùng điện áp cao.
\item (Đúng) Khi giải cần xác định đúng dữ kiện, phương chiều, điều kiện áp dụng và đơn vị.
\item (Sai) Kết quả không phụ thuộc dữ kiện và có thể bỏ qua điều kiện.
\end{itemchoice}
}
\end{ex}

% ===== Câu 2 =====
% ID: L12C3B18-01-TL
% Mức: TH
\begin{bt}
% ID: L12C3B18-01-TL
% Mức: TH
Trình bày kiến thức cốt lõi của Bài toán tổng hợp về máy biến áp, truyền tải điện và an toàn điện và nêu điều kiện áp dụng.
\loigiai{
Cần nêu được: An toàn điện yêu cầu cách điện, nối đất phù hợp và không chạm phần dẫn điện đang có điện áp nguy hiểm. Khi vận dụng phải xác định đúng dữ kiện, phương chiều nếu có, điều kiện áp dụng, hệ đơn vị và kiểm tra tính hợp lí. Nhận định “Có thể bỏ qua nối đất và thiết bị bảo vệ khi dùng điện áp cao.” sai.
}
\end{bt}

% ===== Câu 3 =====
% ID: L12C3B18-01-TLN
% Mức: TH
\begin{ex}
% ID: L12C3B18-01-TLN
% Mức: TH
Trong bốn phát biểu về Bài toán tổng hợp về máy biến áp, truyền tải điện và an toàn điện, có bao nhiêu phát biểu đúng? (Chỉ ghi một số nguyên.) (1) An toàn điện yêu cầu cách điện, nối đất phù hợp và không chạm phần dẫn điện đang có điện áp nguy hiểm. (2) Có thể bỏ qua nối đất và thiết bị bảo vệ khi dùng điện áp cao. (3) Cần kiểm tra điều kiện áp dụng. (4) Có thể bỏ qua đơn vị.
\shortans{2}
\loigiai{
Đối chiếu với kiến thức: An toàn điện yêu cầu cách điện, nối đất phù hợp và không chạm phần dẫn điện đang có điện áp nguy hiểm. Có 2 phát biểu đúng.
}
\end{ex}

% ===== Câu 4 =====
% ID: L12C3B18-01-TN
% Mức: NB
\begin{ex}
% ID: L12C3B18-01-TN
% Mức: NB
Phát biểu nào sau đây đúng về Bài toán tổng hợp về máy biến áp, truyền tải điện và an toàn điện?
\choice
{Có thể bỏ qua nối đất và thiết bị bảo vệ khi dùng điện áp cao.}
{Có thể bỏ qua phương, chiều và đơn vị trong mọi trường hợp.}
{Kết quả không phụ thuộc dữ kiện và điều kiện áp dụng.}
{\True An toàn điện yêu cầu cách điện, nối đất phù hợp và không chạm phần dẫn điện đang có điện áp nguy hiểm.}
\loigiai{
An toàn điện yêu cầu cách điện, nối đất phù hợp và không chạm phần dẫn điện đang có điện áp nguy hiểm. Vì vậy phương án $D$ đúng; các phương án còn lại sai.
}
\end{ex}

% ===== Câu 5 =====
% ID: L12C3B18-02-DS
% Mức: TH
\begin{ex}
% ID: L12C3B18-02-DS
% Mức: TH
Xét Bài toán tổng hợp về máy biến áp, truyền tải điện và an toàn điện (Tình huống 2). Hãy xác định tính đúng, sai của bốn phát biểu sau.
\choiceTF
{Có thể bỏ qua nối đất và thiết bị bảo vệ khi dùng điện áp cao.}
{\True An toàn điện yêu cầu cách điện, nối đất phù hợp và không chạm phần dẫn điện đang có điện áp nguy hiểm.}
{Kết quả không phụ thuộc dữ kiện và có thể bỏ qua điều kiện.}
{\True Khi giải cần xác định đúng dữ kiện, phương chiều, điều kiện áp dụng và đơn vị.}
\loigiai{
\begin{itemchoice}
\item (Sai) Có thể bỏ qua nối đất và thiết bị bảo vệ khi dùng điện áp cao. (Vì): Có thể bỏ qua nối đất và thiết bị bảo vệ khi dùng điện áp cao. b) Đúng: An toàn điện yêu cầu cách điện, nối đất phù hợp và không chạm phần dẫn điện đang có điện áp nguy hiểm. c) Sai: Kết quả không phụ thuộc dữ kiện và có thể bỏ qua điều kiện. d) Đúng: Khi giải cần xác định đúng dữ kiện, phương chiều, điều kiện áp dụng và đơn vị. Kiến thức cốt lõi: An toàn điện yêu cầu cách điện, nối đất phù hợp và không chạm phần dẫn điện đang có điện áp nguy hiểm
\item (Đúng) An toàn điện yêu cầu cách điện, nối đất phù hợp và không chạm phần dẫn điện đang có điện áp nguy hiểm.
\item (Sai) Kết quả không phụ thuộc dữ kiện và có thể bỏ qua điều kiện.
\item (Đúng) Khi giải cần xác định đúng dữ kiện, phương chiều, điều kiện áp dụng và đơn vị.
\end{itemchoice}
}
\end{ex}

% ===== Câu 6 =====
% ID: L12C3B18-02-TL
% Mức: TH
\begin{bt}
% ID: L12C3B18-02-TL
% Mức: TH
Giải thích vì sao nhận định sau đúng: “An toàn điện yêu cầu cách điện, nối đất phù hợp và không chạm phần dẫn điện đang có điện áp nguy hiểm.”
\loigiai{
Cần nêu được: An toàn điện yêu cầu cách điện, nối đất phù hợp và không chạm phần dẫn điện đang có điện áp nguy hiểm. Khi vận dụng phải xác định đúng dữ kiện, phương chiều nếu có, điều kiện áp dụng, hệ đơn vị và kiểm tra tính hợp lí. Nhận định “Có thể bỏ qua nối đất và thiết bị bảo vệ khi dùng điện áp cao.” sai.
}
\end{bt}

% ===== Câu 7 =====
% ID: L12C3B18-02-TLN
% Mức: TH
\begin{ex}
% ID: L12C3B18-02-TLN
% Mức: TH
Trong bốn phát biểu về Bài toán tổng hợp về máy biến áp, truyền tải điện và an toàn điện, có bao nhiêu phát biểu đúng? (Chỉ ghi một số nguyên.) (1) Có thể bỏ qua nối đất và thiết bị bảo vệ khi dùng điện áp cao. (2) An toàn điện yêu cầu cách điện, nối đất phù hợp và không chạm phần dẫn điện đang có điện áp nguy hiểm. (3) Kết quả phải phù hợp ý nghĩa vật lí. (4) Mọi đại lượng đều là vô hướng.
\shortans{1}
\loigiai{
Đối chiếu với kiến thức: An toàn điện yêu cầu cách điện, nối đất phù hợp và không chạm phần dẫn điện đang có điện áp nguy hiểm. Có 1 phát biểu đúng.
}
\end{ex}

% ===== Câu 8 =====
% ID: L12C3B18-02-TN
% Mức: NB
\begin{ex}
% ID: L12C3B18-02-TN
% Mức: NB
Trong nội dung Bài toán tổng hợp về máy biến áp, truyền tải điện và an toàn điện?
\choice
{\True An toàn điện yêu cầu cách điện, nối đất phù hợp và không chạm phần dẫn điện đang có điện áp nguy hiểm.}
{Có thể bỏ qua nối đất và thiết bị bảo vệ khi dùng điện áp cao.}
{Có thể bỏ qua phương, chiều và đơn vị trong mọi trường hợp.}
{Kết quả không phụ thuộc dữ kiện và điều kiện áp dụng.}
\loigiai{
An toàn điện yêu cầu cách điện, nối đất phù hợp và không chạm phần dẫn điện đang có điện áp nguy hiểm. Vì vậy phương án $A$ đúng; các phương án còn lại sai.
}
\end{ex}

% ===== Câu 9 =====
% ID: L12C3B18-03-DS
% Mức: VD
\begin{ex}
% ID: L12C3B18-03-DS
% Mức: VD
Xét Bài toán tổng hợp về máy biến áp, truyền tải điện và an toàn điện (Tình huống 3). Hãy xác định tính đúng, sai của bốn phát biểu sau.
\choiceTF
{\True Khi giải cần xác định đúng dữ kiện, phương chiều, điều kiện áp dụng và đơn vị.}
{Kết quả không phụ thuộc dữ kiện và có thể bỏ qua điều kiện.}
{\True An toàn điện yêu cầu cách điện, nối đất phù hợp và không chạm phần dẫn điện đang có điện áp nguy hiểm.}
{Có thể bỏ qua nối đất và thiết bị bảo vệ khi dùng điện áp cao.}
\loigiai{
\begin{itemchoice}
\item (Đúng) Khi giải cần xác định đúng dữ kiện, phương chiều, điều kiện áp dụng và đơn vị. (Vì): Khi giải cần xác định đúng dữ kiện, phương chiều, điều kiện áp dụng và đơn vị. b) Sai: Kết quả không phụ thuộc dữ kiện và có thể bỏ qua điều kiện. c) Đúng: An toàn điện yêu cầu cách điện, nối đất phù hợp và không chạm phần dẫn điện đang có điện áp nguy hiểm. d) Sai: Có thể bỏ qua nối đất và thiết bị bảo vệ khi dùng điện áp cao. Kiến thức cốt lõi: An toàn điện yêu cầu cách điện, nối đất phù hợp và không chạm phần dẫn điện đang có điện áp nguy hiểm
\item (Sai) Kết quả không phụ thuộc dữ kiện và có thể bỏ qua điều kiện.
\item (Đúng) An toàn điện yêu cầu cách điện, nối đất phù hợp và không chạm phần dẫn điện đang có điện áp nguy hiểm.
\item (Sai) Có thể bỏ qua nối đất và thiết bị bảo vệ khi dùng điện áp cao.
\end{itemchoice}
}
\end{ex}

% ===== Câu 10 =====
% ID: L12C3B18-03-TL
% Mức: VD
\begin{bt}
% ID: L12C3B18-03-TL
% Mức: VD
Phân tích sai lầm trong nhận định: “Có thể bỏ qua nối đất và thiết bị bảo vệ khi dùng điện áp cao.”
\loigiai{
Cần nêu được: An toàn điện yêu cầu cách điện, nối đất phù hợp và không chạm phần dẫn điện đang có điện áp nguy hiểm. Khi vận dụng phải xác định đúng dữ kiện, phương chiều nếu có, điều kiện áp dụng, hệ đơn vị và kiểm tra tính hợp lí. Nhận định “Có thể bỏ qua nối đất và thiết bị bảo vệ khi dùng điện áp cao.” sai.
}
\end{bt}

% ===== Câu 11 =====
% ID: L12C3B18-03-TLN
% Mức: TH
\begin{ex}
% ID: L12C3B18-03-TLN
% Mức: TH
Trong bốn phát biểu về Bài toán tổng hợp về máy biến áp, truyền tải điện và an toàn điện, có bao nhiêu phát biểu đúng? (Chỉ ghi một số nguyên.) (1) An toàn điện yêu cầu cách điện, nối đất phù hợp và không chạm phần dẫn điện đang có điện áp nguy hiểm. (2) Có thể bỏ qua nối đất và thiết bị bảo vệ khi dùng điện áp cao. (3) Cần kiểm tra điều kiện áp dụng. (4) Có thể bỏ qua đơn vị.
\shortans{2}
\loigiai{
Đối chiếu với kiến thức: An toàn điện yêu cầu cách điện, nối đất phù hợp và không chạm phần dẫn điện đang có điện áp nguy hiểm. Có 2 phát biểu đúng.
}
\end{ex}

% ===== Câu 12 =====
% ID: L12C3B18-03-TN
% Mức: NB
\begin{ex}
% ID: L12C3B18-03-TN
% Mức: NB
Tình huống 3: chọn kết luận chính xác về Bài toán tổng hợp về máy biến áp, truyền tải điện và an toàn điện?
\choice
{Kết quả không phụ thuộc dữ kiện và điều kiện áp dụng.}
{\True An toàn điện yêu cầu cách điện, nối đất phù hợp và không chạm phần dẫn điện đang có điện áp nguy hiểm.}
{Có thể bỏ qua nối đất và thiết bị bảo vệ khi dùng điện áp cao.}
{Có thể bỏ qua phương, chiều và đơn vị trong mọi trường hợp.}
\loigiai{
An toàn điện yêu cầu cách điện, nối đất phù hợp và không chạm phần dẫn điện đang có điện áp nguy hiểm. Vì vậy phương án $B$ đúng; các phương án còn lại sai.
}
\end{ex}

% ===== Câu 13 =====
% ID: L12C3B18-04-DS
% Mức: VD
\begin{ex}
% ID: L12C3B18-04-DS
% Mức: VD
Xét Bài toán tổng hợp về máy biến áp, truyền tải điện và an toàn điện (Tình huống 4). Hãy xác định tính đúng, sai của bốn phát biểu sau.
\choiceTF
{Kết quả không phụ thuộc dữ kiện và có thể bỏ qua điều kiện.}
{\True Khi giải cần xác định đúng dữ kiện, phương chiều, điều kiện áp dụng và đơn vị.}
{Có thể bỏ qua nối đất và thiết bị bảo vệ khi dùng điện áp cao.}
{\True An toàn điện yêu cầu cách điện, nối đất phù hợp và không chạm phần dẫn điện đang có điện áp nguy hiểm.}
\loigiai{
\begin{itemchoice}
\item (Sai) Kết quả không phụ thuộc dữ kiện và có thể bỏ qua điều kiện. (Vì): Kết quả không phụ thuộc dữ kiện và có thể bỏ qua điều kiện. b) Đúng: Khi giải cần xác định đúng dữ kiện, phương chiều, điều kiện áp dụng và đơn vị. c) Sai: Có thể bỏ qua nối đất và thiết bị bảo vệ khi dùng điện áp cao. d) Đúng: An toàn điện yêu cầu cách điện, nối đất phù hợp và không chạm phần dẫn điện đang có điện áp nguy hiểm. Kiến thức cốt lõi: An toàn điện yêu cầu cách điện, nối đất phù hợp và không chạm phần dẫn điện đang có điện áp nguy hiểm
\item (Đúng) Khi giải cần xác định đúng dữ kiện, phương chiều, điều kiện áp dụng và đơn vị.
\item (Sai) Có thể bỏ qua nối đất và thiết bị bảo vệ khi dùng điện áp cao.
\item (Đúng) An toàn điện yêu cầu cách điện, nối đất phù hợp và không chạm phần dẫn điện đang có điện áp nguy hiểm.
\end{itemchoice}
}
\end{ex}

% ===== Câu 14 =====
% ID: L12C3B18-04-TL
% Mức: VD
\begin{bt}
% ID: L12C3B18-04-TL
% Mức: VD
Đề xuất các bước giải một bài toán thuộc dạng Bài toán tổng hợp về máy biến áp, truyền tải điện và an toàn điện.
\loigiai{
Cần nêu được: An toàn điện yêu cầu cách điện, nối đất phù hợp và không chạm phần dẫn điện đang có điện áp nguy hiểm. Khi vận dụng phải xác định đúng dữ kiện, phương chiều nếu có, điều kiện áp dụng, hệ đơn vị và kiểm tra tính hợp lí. Nhận định “Có thể bỏ qua nối đất và thiết bị bảo vệ khi dùng điện áp cao.” sai.
}
\end{bt}

% ===== Câu 15 =====
% ID: L12C3B18-04-TLN
% Mức: VD
\begin{ex}
% ID: L12C3B18-04-TLN
% Mức: VD
Trong bốn phát biểu về Bài toán tổng hợp về máy biến áp, truyền tải điện và an toàn điện, có bao nhiêu phát biểu đúng? (Chỉ ghi một số nguyên.) (1) Có thể bỏ qua nối đất và thiết bị bảo vệ khi dùng điện áp cao. (2) An toàn điện yêu cầu cách điện, nối đất phù hợp và không chạm phần dẫn điện đang có điện áp nguy hiểm. (3) Kết quả phải phù hợp ý nghĩa vật lí. (4) Mọi đại lượng đều là vô hướng.
\shortans{2}
\loigiai{
Đối chiếu với kiến thức: An toàn điện yêu cầu cách điện, nối đất phù hợp và không chạm phần dẫn điện đang có điện áp nguy hiểm. Có 2 phát biểu đúng.
}
\end{ex}

% ===== Câu 16 =====
% ID: L12C3B18-04-TN
% Mức: TH
\begin{ex}
% ID: L12C3B18-04-TN
% Mức: TH
Nội dung nào mô tả đúng nhất Bài toán tổng hợp về máy biến áp, truyền tải điện và an toàn điện?
\choice
{Có thể bỏ qua phương, chiều và đơn vị trong mọi trường hợp.}
{Kết quả không phụ thuộc dữ kiện và điều kiện áp dụng.}
{\True An toàn điện yêu cầu cách điện, nối đất phù hợp và không chạm phần dẫn điện đang có điện áp nguy hiểm.}
{Có thể bỏ qua nối đất và thiết bị bảo vệ khi dùng điện áp cao.}
\loigiai{
An toàn điện yêu cầu cách điện, nối đất phù hợp và không chạm phần dẫn điện đang có điện áp nguy hiểm. Vì vậy phương án $C$ đúng; các phương án còn lại sai.
}
\end{ex}

% ===== Câu 17 =====
% ID: L12C3B18-05-DS
% Mức: VD
\begin{ex}
% ID: L12C3B18-05-DS
% Mức: VD
Xét Bài toán tổng hợp về máy biến áp, truyền tải điện và an toàn điện (Tình huống 5). Hãy xác định tính đúng, sai của bốn phát biểu sau.
\choiceTF
{\True An toàn điện yêu cầu cách điện, nối đất phù hợp và không chạm phần dẫn điện đang có điện áp nguy hiểm.}
{Kết quả không phụ thuộc dữ kiện và có thể bỏ qua điều kiện.}
{Có thể bỏ qua nối đất và thiết bị bảo vệ khi dùng điện áp cao.}
{\True Khi giải cần xác định đúng dữ kiện, phương chiều, điều kiện áp dụng và đơn vị.}
\loigiai{
\begin{itemchoice}
\item (Đúng) An toàn điện yêu cầu cách điện, nối đất phù hợp và không chạm phần dẫn điện đang có điện áp nguy hiểm. (Vì): An toàn điện yêu cầu cách điện, nối đất phù hợp và không chạm phần dẫn điện đang có điện áp nguy hiểm. b) Sai: Kết quả không phụ thuộc dữ kiện và có thể bỏ qua điều kiện. c) Sai: Có thể bỏ qua nối đất và thiết bị bảo vệ khi dùng điện áp cao. d) Đúng: Khi giải cần xác định đúng dữ kiện, phương chiều, điều kiện áp dụng và đơn vị. Kiến thức cốt lõi: An toàn điện yêu cầu cách điện, nối đất phù hợp và không chạm phần dẫn điện đang có điện áp nguy hiểm
\item (Sai) Kết quả không phụ thuộc dữ kiện và có thể bỏ qua điều kiện.
\item (Sai) Có thể bỏ qua nối đất và thiết bị bảo vệ khi dùng điện áp cao.
\item (Đúng) Khi giải cần xác định đúng dữ kiện, phương chiều, điều kiện áp dụng và đơn vị.
\end{itemchoice}
}
\end{ex}

% ===== Câu 18 =====
% ID: L12C3B18-05-TL
% Mức: VD
\begin{bt}
% ID: L12C3B18-05-TL
% Mức: VD
Nêu một tình huống thực tiễn liên quan đến Bài toán tổng hợp về máy biến áp, truyền tải điện và an toàn điện và giải thích bằng kiến thức Vật lí.
\loigiai{
Cần nêu được: An toàn điện yêu cầu cách điện, nối đất phù hợp và không chạm phần dẫn điện đang có điện áp nguy hiểm. Khi vận dụng phải xác định đúng dữ kiện, phương chiều nếu có, điều kiện áp dụng, hệ đơn vị và kiểm tra tính hợp lí. Nhận định “Có thể bỏ qua nối đất và thiết bị bảo vệ khi dùng điện áp cao.” sai.
}
\end{bt}

% ===== Câu 19 =====
% ID: L12C3B18-05-TLN
% Mức: VD
\begin{ex}
% ID: L12C3B18-05-TLN
% Mức: VD
Trong bốn phát biểu về Bài toán tổng hợp về máy biến áp, truyền tải điện và an toàn điện, có bao nhiêu phát biểu đúng? (Chỉ ghi một số nguyên.) (1) An toàn điện yêu cầu cách điện, nối đất phù hợp và không chạm phần dẫn điện đang có điện áp nguy hiểm. (2) Có thể bỏ qua nối đất và thiết bị bảo vệ khi dùng điện áp cao. (3) Cần kiểm tra điều kiện áp dụng. (4) Có thể bỏ qua đơn vị.
\shortans{2}
\loigiai{
Đối chiếu với kiến thức: An toàn điện yêu cầu cách điện, nối đất phù hợp và không chạm phần dẫn điện đang có điện áp nguy hiểm. Có 2 phát biểu đúng.
}
\end{ex}

% ===== Câu 20 =====
% ID: L12C3B18-05-TN
% Mức: TH
\begin{ex}
% ID: L12C3B18-05-TN
% Mức: TH
Khi phân tích Bài toán tổng hợp về máy biến áp, truyền tải điện và an toàn điện?
\choice
{Có thể bỏ qua nối đất và thiết bị bảo vệ khi dùng điện áp cao.}
{Có thể bỏ qua phương, chiều và đơn vị trong mọi trường hợp.}
{Kết quả không phụ thuộc dữ kiện và điều kiện áp dụng.}
{\True An toàn điện yêu cầu cách điện, nối đất phù hợp và không chạm phần dẫn điện đang có điện áp nguy hiểm.}
\loigiai{
An toàn điện yêu cầu cách điện, nối đất phù hợp và không chạm phần dẫn điện đang có điện áp nguy hiểm. Vì vậy phương án $D$ đúng; các phương án còn lại sai.
}
\end{ex}

% ===== Câu 21 =====
% ID: L12C3B18-06-DS
% Mức: TH
\dangbt{Giải thích nguyên lí bếp từ, phanh điện từ và gia nhiệt cảm ứng}
\begin{ex}
% ID: L12C3B18-06-DS
% Mức: TH
Xét Giải thích nguyên lí bếp từ, phanh điện từ và gia nhiệt cảm ứng (Tình huống 1). Hãy xác định tính đúng, sai của bốn phát biểu sau.
\choiceTF
{\True Bếp từ làm nóng đáy nồi dẫn điện nhờ dòng điện Foucault và sự tỏa nhiệt Joule.}
{Bếp từ làm nóng trực tiếp mọi vật liệu bằng điện trường tĩnh.}
{\True Khi giải cần xác định đúng dữ kiện, phương chiều, điều kiện áp dụng và đơn vị.}
{Kết quả không phụ thuộc dữ kiện và có thể bỏ qua điều kiện.}
\loigiai{
\begin{itemchoice}
\item (Đúng) Bếp từ làm nóng đáy nồi dẫn điện nhờ dòng điện Foucault và sự tỏa nhiệt Joule. (Vì): Bếp từ làm nóng đáy nồi dẫn điện nhờ dòng điện Foucault và sự tỏa nhiệt Joule. b) Sai: Bếp từ làm nóng trực tiếp mọi vật liệu bằng điện trường tĩnh. c) Đúng: Khi giải cần xác định đúng dữ kiện, phương chiều, điều kiện áp dụng và đơn vị. d) Sai: Kết quả không phụ thuộc dữ kiện và có thể bỏ qua điều kiện. Kiến thức cốt lõi: Bếp từ làm nóng đáy nồi dẫn điện nhờ dòng điện Foucault và sự tỏa nhiệt Joule
\item (Sai) Bếp từ làm nóng trực tiếp mọi vật liệu bằng điện trường tĩnh.
\item (Đúng) Khi giải cần xác định đúng dữ kiện, phương chiều, điều kiện áp dụng và đơn vị.
\item (Sai) Kết quả không phụ thuộc dữ kiện và có thể bỏ qua điều kiện.
\end{itemchoice}
}
\end{ex}

% ===== Câu 22 =====
% ID: L12C3B18-06-TL
% Mức: VDC
\dangbt{Bài toán tổng hợp về máy biến áp, truyền tải điện và an toàn điện}
\begin{bt}
% ID: L12C3B18-06-TL
% Mức: VDC
So sánh nhận định đúng “An toàn điện yêu cầu cách điện, nối đất phù hợp và không chạm phần dẫn điện đang có điện áp nguy hiểm.” với nhận định sai “Có thể bỏ qua nối đất và thiết bị bảo vệ khi dùng điện áp cao.”; chỉ rõ căn cứ Vật lí.
\loigiai{
Cần nêu được: An toàn điện yêu cầu cách điện, nối đất phù hợp và không chạm phần dẫn điện đang có điện áp nguy hiểm. Khi vận dụng phải xác định đúng dữ kiện, phương chiều nếu có, điều kiện áp dụng, hệ đơn vị và kiểm tra tính hợp lí. Nhận định “Có thể bỏ qua nối đất và thiết bị bảo vệ khi dùng điện áp cao.” sai.
}
\end{bt}

% ===== Câu 23 =====
% ID: L12C3B18-06-TLN
% Mức: VDC
\begin{ex}
% ID: L12C3B18-06-TLN
% Mức: VDC
Trong bốn phát biểu về Bài toán tổng hợp về máy biến áp, truyền tải điện và an toàn điện, có bao nhiêu phát biểu đúng? (Chỉ ghi một số nguyên.) (1) Khi giải cần xác định đúng dữ kiện, phương chiều, điều kiện áp dụng và đơn vị. (2) Có thể bỏ qua nối đất và thiết bị bảo vệ khi dùng điện áp cao. (3) Phải dùng đúng định luật cho tình huống. (4) An toàn điện yêu cầu cách điện, nối đất phù hợp và không chạm phần dẫn điện đang có điện áp nguy hiểm.
\shortans{3}
\loigiai{
Đối chiếu với kiến thức: An toàn điện yêu cầu cách điện, nối đất phù hợp và không chạm phần dẫn điện đang có điện áp nguy hiểm. Có 3 phát biểu đúng.
}
\end{ex}

% ===== Câu 24 =====
% ID: L12C3B18-06-TN
% Mức: TH
\begin{ex}
% ID: L12C3B18-06-TN
% Mức: TH
Một học sinh ôn tập Bài toán tổng hợp về máy biến áp, truyền tải điện và an toàn điện?
\choice
{\True An toàn điện yêu cầu cách điện, nối đất phù hợp và không chạm phần dẫn điện đang có điện áp nguy hiểm.}
{Có thể bỏ qua nối đất và thiết bị bảo vệ khi dùng điện áp cao.}
{Có thể bỏ qua phương, chiều và đơn vị trong mọi trường hợp.}
{Kết quả không phụ thuộc dữ kiện và điều kiện áp dụng.}
\loigiai{
An toàn điện yêu cầu cách điện, nối đất phù hợp và không chạm phần dẫn điện đang có điện áp nguy hiểm. Vì vậy phương án $A$ đúng; các phương án còn lại sai.
}
\end{ex}

% ===== Câu 25 =====
% ID: L12C3B18-07-DS
% Mức: TH
\dangbt{Giải thích nguyên lí bếp từ, phanh điện từ và gia nhiệt cảm ứng}
\begin{ex}
% ID: L12C3B18-07-DS
% Mức: TH
Xét Giải thích nguyên lí bếp từ, phanh điện từ và gia nhiệt cảm ứng (Tình huống 2). Hãy xác định tính đúng, sai của bốn phát biểu sau.
\choiceTF
{Bếp từ làm nóng trực tiếp mọi vật liệu bằng điện trường tĩnh.}
{\True Bếp từ làm nóng đáy nồi dẫn điện nhờ dòng điện Foucault và sự tỏa nhiệt Joule.}
{Kết quả không phụ thuộc dữ kiện và có thể bỏ qua điều kiện.}
{\True Khi giải cần xác định đúng dữ kiện, phương chiều, điều kiện áp dụng và đơn vị.}
\loigiai{
\begin{itemchoice}
\item (Sai) Bếp từ làm nóng trực tiếp mọi vật liệu bằng điện trường tĩnh. (Vì): Bếp từ làm nóng trực tiếp mọi vật liệu bằng điện trường tĩnh. b) Đúng: Bếp từ làm nóng đáy nồi dẫn điện nhờ dòng điện Foucault và sự tỏa nhiệt Joule. c) Sai: Kết quả không phụ thuộc dữ kiện và có thể bỏ qua điều kiện. d) Đúng: Khi giải cần xác định đúng dữ kiện, phương chiều, điều kiện áp dụng và đơn vị. Kiến thức cốt lõi: Bếp từ làm nóng đáy nồi dẫn điện nhờ dòng điện Foucault và sự tỏa nhiệt Joule
\item (Đúng) Bếp từ làm nóng đáy nồi dẫn điện nhờ dòng điện Foucault và sự tỏa nhiệt Joule.
\item (Sai) Kết quả không phụ thuộc dữ kiện và có thể bỏ qua điều kiện.
\item (Đúng) Khi giải cần xác định đúng dữ kiện, phương chiều, điều kiện áp dụng và đơn vị.
\end{itemchoice}
}
\end{ex}

% ===== Câu 26 =====
% ID: L12C3B18-07-TL
% Mức: TH
\begin{bt}
% ID: L12C3B18-07-TL
% Mức: TH
Trình bày kiến thức cốt lõi của Giải thích nguyên lí bếp từ, phanh điện từ và gia nhiệt cảm ứng và nêu điều kiện áp dụng.
\loigiai{
Cần nêu được: Bếp từ làm nóng đáy nồi dẫn điện nhờ dòng điện Foucault và sự tỏa nhiệt Joule. Khi vận dụng phải xác định đúng dữ kiện, phương chiều nếu có, điều kiện áp dụng, hệ đơn vị và kiểm tra tính hợp lí. Nhận định “Bếp từ làm nóng trực tiếp mọi vật liệu bằng điện trường tĩnh.” sai.
}
\end{bt}

% ===== Câu 27 =====
% ID: L12C3B18-07-TLN
% Mức: TH
\begin{ex}
% ID: L12C3B18-07-TLN
% Mức: TH
Trong bốn phát biểu về Giải thích nguyên lí bếp từ, phanh điện từ và gia nhiệt cảm ứng, có bao nhiêu phát biểu đúng? (Chỉ ghi một số nguyên.) (1) Bếp từ làm nóng đáy nồi dẫn điện nhờ dòng điện Foucault và sự tỏa nhiệt Joule. (2) Bếp từ làm nóng trực tiếp mọi vật liệu bằng điện trường tĩnh. (3) Cần kiểm tra điều kiện áp dụng. (4) Có thể bỏ qua đơn vị.
\shortans{2}
\loigiai{
Đối chiếu với kiến thức: Bếp từ làm nóng đáy nồi dẫn điện nhờ dòng điện Foucault và sự tỏa nhiệt Joule. Có 2 phát biểu đúng.
}
\end{ex}

% ===== Câu 28 =====
% ID: L12C3B18-07-TN
% Mức: TH
\dangbt{Bài toán tổng hợp về máy biến áp, truyền tải điện và an toàn điện}
\begin{ex}
% ID: L12C3B18-07-TN
% Mức: TH
Chọn phát biểu không trái với kiến thức về Bài toán tổng hợp về máy biến áp, truyền tải điện và an toàn điện?
\choice
{Kết quả không phụ thuộc dữ kiện và điều kiện áp dụng.}
{\True An toàn điện yêu cầu cách điện, nối đất phù hợp và không chạm phần dẫn điện đang có điện áp nguy hiểm.}
{Có thể bỏ qua nối đất và thiết bị bảo vệ khi dùng điện áp cao.}
{Có thể bỏ qua phương, chiều và đơn vị trong mọi trường hợp.}
\loigiai{
An toàn điện yêu cầu cách điện, nối đất phù hợp và không chạm phần dẫn điện đang có điện áp nguy hiểm. Vì vậy phương án $B$ đúng; các phương án còn lại sai.
}
\end{ex}

% ===== Câu 29 =====
% ID: L12C3B18-08-DS
% Mức: VD
\dangbt{Giải thích nguyên lí bếp từ, phanh điện từ và gia nhiệt cảm ứng}
\begin{ex}
% ID: L12C3B18-08-DS
% Mức: VD
Xét Giải thích nguyên lí bếp từ, phanh điện từ và gia nhiệt cảm ứng (Tình huống 3). Hãy xác định tính đúng, sai của bốn phát biểu sau.
\choiceTF
{\True Khi giải cần xác định đúng dữ kiện, phương chiều, điều kiện áp dụng và đơn vị.}
{Kết quả không phụ thuộc dữ kiện và có thể bỏ qua điều kiện.}
{\True Bếp từ làm nóng đáy nồi dẫn điện nhờ dòng điện Foucault và sự tỏa nhiệt Joule.}
{Bếp từ làm nóng trực tiếp mọi vật liệu bằng điện trường tĩnh.}
\loigiai{
\begin{itemchoice}
\item (Đúng) Khi giải cần xác định đúng dữ kiện, phương chiều, điều kiện áp dụng và đơn vị. (Vì): Khi giải cần xác định đúng dữ kiện, phương chiều, điều kiện áp dụng và đơn vị. b) Sai: Kết quả không phụ thuộc dữ kiện và có thể bỏ qua điều kiện. c) Đúng: Bếp từ làm nóng đáy nồi dẫn điện nhờ dòng điện Foucault và sự tỏa nhiệt Joule. d) Sai: Bếp từ làm nóng trực tiếp mọi vật liệu bằng điện trường tĩnh. Kiến thức cốt lõi: Bếp từ làm nóng đáy nồi dẫn điện nhờ dòng điện Foucault và sự tỏa nhiệt Joule
\item (Sai) Kết quả không phụ thuộc dữ kiện và có thể bỏ qua điều kiện.
\item (Đúng) Bếp từ làm nóng đáy nồi dẫn điện nhờ dòng điện Foucault và sự tỏa nhiệt Joule.
\item (Sai) Bếp từ làm nóng trực tiếp mọi vật liệu bằng điện trường tĩnh.
\end{itemchoice}
}
\end{ex}

% ===== Câu 30 =====
% ID: L12C3B18-08-TL
% Mức: TH
\begin{bt}
% ID: L12C3B18-08-TL
% Mức: TH
Giải thích vì sao nhận định sau đúng: “Bếp từ làm nóng đáy nồi dẫn điện nhờ dòng điện Foucault và sự tỏa nhiệt Joule.”
\loigiai{
Cần nêu được: Bếp từ làm nóng đáy nồi dẫn điện nhờ dòng điện Foucault và sự tỏa nhiệt Joule. Khi vận dụng phải xác định đúng dữ kiện, phương chiều nếu có, điều kiện áp dụng, hệ đơn vị và kiểm tra tính hợp lí. Nhận định “Bếp từ làm nóng trực tiếp mọi vật liệu bằng điện trường tĩnh.” sai.
}
\end{bt}

% ===== Câu 31 =====
% ID: L12C3B18-08-TLN
% Mức: TH
\begin{ex}
% ID: L12C3B18-08-TLN
% Mức: TH
Trong bốn phát biểu về Giải thích nguyên lí bếp từ, phanh điện từ và gia nhiệt cảm ứng, có bao nhiêu phát biểu đúng? (Chỉ ghi một số nguyên.) (1) Bếp từ làm nóng trực tiếp mọi vật liệu bằng điện trường tĩnh. (2) Bếp từ làm nóng đáy nồi dẫn điện nhờ dòng điện Foucault và sự tỏa nhiệt Joule. (3) Kết quả phải phù hợp ý nghĩa vật lí. (4) Mọi đại lượng đều là vô hướng.
\shortans{1}
\loigiai{
Đối chiếu với kiến thức: Bếp từ làm nóng đáy nồi dẫn điện nhờ dòng điện Foucault và sự tỏa nhiệt Joule. Có 1 phát biểu đúng.
}
\end{ex}

% ===== Câu 32 =====
% ID: L12C3B18-08-TN
% Mức: VD
\dangbt{Bài toán tổng hợp về máy biến áp, truyền tải điện và an toàn điện}
\begin{ex}
% ID: L12C3B18-08-TN
% Mức: VD
Cơ sở Vật lí đúng dùng để giải Bài toán tổng hợp về máy biến áp, truyền tải điện và an toàn điện?
\choice
{Có thể bỏ qua phương, chiều và đơn vị trong mọi trường hợp.}
{Kết quả không phụ thuộc dữ kiện và điều kiện áp dụng.}
{\True An toàn điện yêu cầu cách điện, nối đất phù hợp và không chạm phần dẫn điện đang có điện áp nguy hiểm.}
{Có thể bỏ qua nối đất và thiết bị bảo vệ khi dùng điện áp cao.}
\loigiai{
An toàn điện yêu cầu cách điện, nối đất phù hợp và không chạm phần dẫn điện đang có điện áp nguy hiểm. Vì vậy phương án $C$ đúng; các phương án còn lại sai.
}
\end{ex}

% ===== Câu 33 =====
% ID: L12C3B18-09-DS
% Mức: VD
\dangbt{Giải thích nguyên lí bếp từ, phanh điện từ và gia nhiệt cảm ứng}
\begin{ex}
% ID: L12C3B18-09-DS
% Mức: VD
Xét Giải thích nguyên lí bếp từ, phanh điện từ và gia nhiệt cảm ứng (Tình huống 4). Hãy xác định tính đúng, sai của bốn phát biểu sau.
\choiceTF
{Kết quả không phụ thuộc dữ kiện và có thể bỏ qua điều kiện.}
{\True Khi giải cần xác định đúng dữ kiện, phương chiều, điều kiện áp dụng và đơn vị.}
{Bếp từ làm nóng trực tiếp mọi vật liệu bằng điện trường tĩnh.}
{\True Bếp từ làm nóng đáy nồi dẫn điện nhờ dòng điện Foucault và sự tỏa nhiệt Joule.}
\loigiai{
\begin{itemchoice}
\item (Sai) Kết quả không phụ thuộc dữ kiện và có thể bỏ qua điều kiện. (Vì): Kết quả không phụ thuộc dữ kiện và có thể bỏ qua điều kiện. b) Đúng: Khi giải cần xác định đúng dữ kiện, phương chiều, điều kiện áp dụng và đơn vị. c) Sai: Bếp từ làm nóng trực tiếp mọi vật liệu bằng điện trường tĩnh. d) Đúng: Bếp từ làm nóng đáy nồi dẫn điện nhờ dòng điện Foucault và sự tỏa nhiệt Joule. Kiến thức cốt lõi: Bếp từ làm nóng đáy nồi dẫn điện nhờ dòng điện Foucault và sự tỏa nhiệt Joule
\item (Đúng) Khi giải cần xác định đúng dữ kiện, phương chiều, điều kiện áp dụng và đơn vị.
\item (Sai) Bếp từ làm nóng trực tiếp mọi vật liệu bằng điện trường tĩnh.
\item (Đúng) Bếp từ làm nóng đáy nồi dẫn điện nhờ dòng điện Foucault và sự tỏa nhiệt Joule.
\end{itemchoice}
}
\end{ex}

% ===== Câu 34 =====
% ID: L12C3B18-09-TL
% Mức: VD
\begin{bt}
% ID: L12C3B18-09-TL
% Mức: VD
Phân tích sai lầm trong nhận định: “Bếp từ làm nóng trực tiếp mọi vật liệu bằng điện trường tĩnh.”
\loigiai{
Cần nêu được: Bếp từ làm nóng đáy nồi dẫn điện nhờ dòng điện Foucault và sự tỏa nhiệt Joule. Khi vận dụng phải xác định đúng dữ kiện, phương chiều nếu có, điều kiện áp dụng, hệ đơn vị và kiểm tra tính hợp lí. Nhận định “Bếp từ làm nóng trực tiếp mọi vật liệu bằng điện trường tĩnh.” sai.
}
\end{bt}

% ===== Câu 35 =====
% ID: L12C3B18-09-TLN
% Mức: TH
\begin{ex}
% ID: L12C3B18-09-TLN
% Mức: TH
Trong bốn phát biểu về Giải thích nguyên lí bếp từ, phanh điện từ và gia nhiệt cảm ứng, có bao nhiêu phát biểu đúng? (Chỉ ghi một số nguyên.) (1) Bếp từ làm nóng đáy nồi dẫn điện nhờ dòng điện Foucault và sự tỏa nhiệt Joule. (2) Bếp từ làm nóng trực tiếp mọi vật liệu bằng điện trường tĩnh. (3) Cần kiểm tra điều kiện áp dụng. (4) Có thể bỏ qua đơn vị.
\shortans{2}
\loigiai{
Đối chiếu với kiến thức: Bếp từ làm nóng đáy nồi dẫn điện nhờ dòng điện Foucault và sự tỏa nhiệt Joule. Có 2 phát biểu đúng.
}
\end{ex}

% ===== Câu 36 =====
% ID: L12C3B18-09-TN
% Mức: VD
\dangbt{Bài toán tổng hợp về máy biến áp, truyền tải điện và an toàn điện}
\begin{ex}
% ID: L12C3B18-09-TN
% Mức: VD
Trong một bài thực hành về Bài toán tổng hợp về máy biến áp, truyền tải điện và an toàn điện?
\choice
{Có thể bỏ qua nối đất và thiết bị bảo vệ khi dùng điện áp cao.}
{Có thể bỏ qua phương, chiều và đơn vị trong mọi trường hợp.}
{Kết quả không phụ thuộc dữ kiện và điều kiện áp dụng.}
{\True An toàn điện yêu cầu cách điện, nối đất phù hợp và không chạm phần dẫn điện đang có điện áp nguy hiểm.}
\loigiai{
An toàn điện yêu cầu cách điện, nối đất phù hợp và không chạm phần dẫn điện đang có điện áp nguy hiểm. Vì vậy phương án $D$ đúng; các phương án còn lại sai.
}
\end{ex}

% ===== Câu 37 =====
% ID: L12C3B18-10-DS
% Mức: VD
\dangbt{Giải thích nguyên lí bếp từ, phanh điện từ và gia nhiệt cảm ứng}
\begin{ex}
% ID: L12C3B18-10-DS
% Mức: VD
Xét Giải thích nguyên lí bếp từ, phanh điện từ và gia nhiệt cảm ứng (Tình huống 5). Hãy xác định tính đúng, sai của bốn phát biểu sau.
\choiceTF
{\True Bếp từ làm nóng đáy nồi dẫn điện nhờ dòng điện Foucault và sự tỏa nhiệt Joule.}
{Kết quả không phụ thuộc dữ kiện và có thể bỏ qua điều kiện.}
{Bếp từ làm nóng trực tiếp mọi vật liệu bằng điện trường tĩnh.}
{\True Khi giải cần xác định đúng dữ kiện, phương chiều, điều kiện áp dụng và đơn vị.}
\loigiai{
\begin{itemchoice}
\item (Đúng) Bếp từ làm nóng đáy nồi dẫn điện nhờ dòng điện Foucault và sự tỏa nhiệt Joule. (Vì): Bếp từ làm nóng đáy nồi dẫn điện nhờ dòng điện Foucault và sự tỏa nhiệt Joule. b) Sai: Kết quả không phụ thuộc dữ kiện và có thể bỏ qua điều kiện. c) Sai: Bếp từ làm nóng trực tiếp mọi vật liệu bằng điện trường tĩnh. d) Đúng: Khi giải cần xác định đúng dữ kiện, phương chiều, điều kiện áp dụng và đơn vị. Kiến thức cốt lõi: Bếp từ làm nóng đáy nồi dẫn điện nhờ dòng điện Foucault và sự tỏa nhiệt Joule
\item (Sai) Kết quả không phụ thuộc dữ kiện và có thể bỏ qua điều kiện.
\item (Sai) Bếp từ làm nóng trực tiếp mọi vật liệu bằng điện trường tĩnh.
\item (Đúng) Khi giải cần xác định đúng dữ kiện, phương chiều, điều kiện áp dụng và đơn vị.
\end{itemchoice}
}
\end{ex}

% ===== Câu 38 =====
% ID: L12C3B18-10-TL
% Mức: VD
\begin{bt}
% ID: L12C3B18-10-TL
% Mức: VD
Đề xuất các bước giải một bài toán thuộc dạng Giải thích nguyên lí bếp từ, phanh điện từ và gia nhiệt cảm ứng.
\loigiai{
Cần nêu được: Bếp từ làm nóng đáy nồi dẫn điện nhờ dòng điện Foucault và sự tỏa nhiệt Joule. Khi vận dụng phải xác định đúng dữ kiện, phương chiều nếu có, điều kiện áp dụng, hệ đơn vị và kiểm tra tính hợp lí. Nhận định “Bếp từ làm nóng trực tiếp mọi vật liệu bằng điện trường tĩnh.” sai.
}
\end{bt}

% ===== Câu 39 =====
% ID: L12C3B18-10-TLN
% Mức: VD
\begin{ex}
% ID: L12C3B18-10-TLN
% Mức: VD
Trong bốn phát biểu về Giải thích nguyên lí bếp từ, phanh điện từ và gia nhiệt cảm ứng, có bao nhiêu phát biểu đúng? (Chỉ ghi một số nguyên.) (1) Bếp từ làm nóng trực tiếp mọi vật liệu bằng điện trường tĩnh. (2) Bếp từ làm nóng đáy nồi dẫn điện nhờ dòng điện Foucault và sự tỏa nhiệt Joule. (3) Kết quả phải phù hợp ý nghĩa vật lí. (4) Mọi đại lượng đều là vô hướng.
\shortans{2}
\loigiai{
Đối chiếu với kiến thức: Bếp từ làm nóng đáy nồi dẫn điện nhờ dòng điện Foucault và sự tỏa nhiệt Joule. Có 2 phát biểu đúng.
}
\end{ex}

% ===== Câu 40 =====
% ID: L12C3B18-10-TN
% Mức: VD
\dangbt{Bài toán tổng hợp về máy biến áp, truyền tải điện và an toàn điện}
\begin{ex}
% ID: L12C3B18-10-TN
% Mức: VD
Kết luận đúng khi vận dụng Bài toán tổng hợp về máy biến áp, truyền tải điện và an toàn điện?
\choice
{\True An toàn điện yêu cầu cách điện, nối đất phù hợp và không chạm phần dẫn điện đang có điện áp nguy hiểm.}
{Có thể bỏ qua nối đất và thiết bị bảo vệ khi dùng điện áp cao.}
{Có thể bỏ qua phương, chiều và đơn vị trong mọi trường hợp.}
{Kết quả không phụ thuộc dữ kiện và điều kiện áp dụng.}
\loigiai{
An toàn điện yêu cầu cách điện, nối đất phù hợp và không chạm phần dẫn điện đang có điện áp nguy hiểm. Vì vậy phương án $A$ đúng; các phương án còn lại sai.
}
\end{ex}

% ===== Câu 41 =====
% ID: L12C3B18-100-TN
% Mức: VD
\dangbt{Tính điện áp, số vòng dây và cường độ dòng điện của máy biến áp lí tưởng}
\begin{ex}
% ID: L12C3B18-100-TN
% Mức: VD
Kết luận đúng khi vận dụng Tính điện áp, số vòng dây và cường độ dòng điện của máy biến áp lí tưởng?
\choice
{\True Máy biến áp lí tưởng thỏa U2/U1 = N2/N1 và U1I1 = U2I2.}
{Máy biến áp lí tưởng luôn có U2 = U1 bất kể số vòng dây.}
{Có thể bỏ qua phương, chiều và đơn vị trong mọi trường hợp.}
{Kết quả không phụ thuộc dữ kiện và điều kiện áp dụng.}
\loigiai{
Máy biến áp lí tưởng thỏa U2/U1 = N2/N1 và U1I1 = U2I2. Vì vậy phương án $A$ đúng; các phương án còn lại sai.
}
\end{ex}

% ===== Câu 42 =====
% ID: L12C3B18-101-TN
% Mức: NB
\dangbt{Truyền tải điện năng đi xa và giảm hao phí trên đường dây}
\begin{ex}
% ID: L12C3B18-101-TN
% Mức: NB
Phát biểu nào sau đây đúng về Truyền tải điện năng đi xa và giảm hao phí trên đường dây?
\choice
{Muốn giảm hao phí phải giảm điện áp truyền tải.}
{Có thể bỏ qua phương, chiều và đơn vị trong mọi trường hợp.}
{Kết quả không phụ thuộc dữ kiện và điều kiện áp dụng.}
{\True Với công suất truyền không đổi, tăng điện áp truyền tải làm giảm dòng điện và giảm hao phí I^2R.}
\loigiai{
Với công suất truyền không đổi, tăng điện áp truyền tải làm giảm dòng điện và giảm hao phí I^2R. Vì vậy phương án $D$ đúng; các phương án còn lại sai.
}
\end{ex}

% ===== Câu 43 =====
% ID: L12C3B18-102-TN
% Mức: NB
\begin{ex}
% ID: L12C3B18-102-TN
% Mức: NB
Trong nội dung Truyền tải điện năng đi xa và giảm hao phí trên đường dây?
\choice
{\True Với công suất truyền không đổi, tăng điện áp truyền tải làm giảm dòng điện và giảm hao phí I^2R.}
{Muốn giảm hao phí phải giảm điện áp truyền tải.}
{Có thể bỏ qua phương, chiều và đơn vị trong mọi trường hợp.}
{Kết quả không phụ thuộc dữ kiện và điều kiện áp dụng.}
\loigiai{
Với công suất truyền không đổi, tăng điện áp truyền tải làm giảm dòng điện và giảm hao phí I^2R. Vì vậy phương án $A$ đúng; các phương án còn lại sai.
}
\end{ex}

% ===== Câu 44 =====
% ID: L12C3B18-103-TN
% Mức: NB
\begin{ex}
% ID: L12C3B18-103-TN
% Mức: NB
Tình huống 3: chọn kết luận chính xác về Truyền tải điện năng đi xa và giảm hao phí trên đường dây?
\choice
{Kết quả không phụ thuộc dữ kiện và điều kiện áp dụng.}
{\True Với công suất truyền không đổi, tăng điện áp truyền tải làm giảm dòng điện và giảm hao phí I^2R.}
{Muốn giảm hao phí phải giảm điện áp truyền tải.}
{Có thể bỏ qua phương, chiều và đơn vị trong mọi trường hợp.}
\loigiai{
Với công suất truyền không đổi, tăng điện áp truyền tải làm giảm dòng điện và giảm hao phí I^2R. Vì vậy phương án $B$ đúng; các phương án còn lại sai.
}
\end{ex}

% ===== Câu 45 =====
% ID: L12C3B18-104-TN
% Mức: TH
\begin{ex}
% ID: L12C3B18-104-TN
% Mức: TH
Nội dung nào mô tả đúng nhất Truyền tải điện năng đi xa và giảm hao phí trên đường dây?
\choice
{Có thể bỏ qua phương, chiều và đơn vị trong mọi trường hợp.}
{Kết quả không phụ thuộc dữ kiện và điều kiện áp dụng.}
{\True Với công suất truyền không đổi, tăng điện áp truyền tải làm giảm dòng điện và giảm hao phí I^2R.}
{Muốn giảm hao phí phải giảm điện áp truyền tải.}
\loigiai{
Với công suất truyền không đổi, tăng điện áp truyền tải làm giảm dòng điện và giảm hao phí I^2R. Vì vậy phương án $C$ đúng; các phương án còn lại sai.
}
\end{ex}

% ===== Câu 46 =====
% ID: L12C3B18-105-TN
% Mức: TH
\begin{ex}
% ID: L12C3B18-105-TN
% Mức: TH
Khi phân tích Truyền tải điện năng đi xa và giảm hao phí trên đường dây?
\choice
{Muốn giảm hao phí phải giảm điện áp truyền tải.}
{Có thể bỏ qua phương, chiều và đơn vị trong mọi trường hợp.}
{Kết quả không phụ thuộc dữ kiện và điều kiện áp dụng.}
{\True Với công suất truyền không đổi, tăng điện áp truyền tải làm giảm dòng điện và giảm hao phí I^2R.}
\loigiai{
Với công suất truyền không đổi, tăng điện áp truyền tải làm giảm dòng điện và giảm hao phí I^2R. Vì vậy phương án $D$ đúng; các phương án còn lại sai.
}
\end{ex}

% ===== Câu 47 =====
% ID: L12C3B18-106-TN
% Mức: TH
\begin{ex}
% ID: L12C3B18-106-TN
% Mức: TH
Một học sinh ôn tập Truyền tải điện năng đi xa và giảm hao phí trên đường dây?
\choice
{\True Với công suất truyền không đổi, tăng điện áp truyền tải làm giảm dòng điện và giảm hao phí I^2R.}
{Muốn giảm hao phí phải giảm điện áp truyền tải.}
{Có thể bỏ qua phương, chiều và đơn vị trong mọi trường hợp.}
{Kết quả không phụ thuộc dữ kiện và điều kiện áp dụng.}
\loigiai{
Với công suất truyền không đổi, tăng điện áp truyền tải làm giảm dòng điện và giảm hao phí I^2R. Vì vậy phương án $A$ đúng; các phương án còn lại sai.
}
\end{ex}

% ===== Câu 48 =====
% ID: L12C3B18-107-TN
% Mức: TH
\begin{ex}
% ID: L12C3B18-107-TN
% Mức: TH
Chọn phát biểu không trái với kiến thức về Truyền tải điện năng đi xa và giảm hao phí trên đường dây?
\choice
{Kết quả không phụ thuộc dữ kiện và điều kiện áp dụng.}
{\True Với công suất truyền không đổi, tăng điện áp truyền tải làm giảm dòng điện và giảm hao phí I^2R.}
{Muốn giảm hao phí phải giảm điện áp truyền tải.}
{Có thể bỏ qua phương, chiều và đơn vị trong mọi trường hợp.}
\loigiai{
Với công suất truyền không đổi, tăng điện áp truyền tải làm giảm dòng điện và giảm hao phí I^2R. Vì vậy phương án $B$ đúng; các phương án còn lại sai.
}
\end{ex}

% ===== Câu 49 =====
% ID: L12C3B18-108-TN
% Mức: VD
\begin{ex}
% ID: L12C3B18-108-TN
% Mức: VD
Cơ sở Vật lí đúng dùng để giải Truyền tải điện năng đi xa và giảm hao phí trên đường dây?
\choice
{Có thể bỏ qua phương, chiều và đơn vị trong mọi trường hợp.}
{Kết quả không phụ thuộc dữ kiện và điều kiện áp dụng.}
{\True Với công suất truyền không đổi, tăng điện áp truyền tải làm giảm dòng điện và giảm hao phí I^2R.}
{Muốn giảm hao phí phải giảm điện áp truyền tải.}
\loigiai{
Với công suất truyền không đổi, tăng điện áp truyền tải làm giảm dòng điện và giảm hao phí I^2R. Vì vậy phương án $C$ đúng; các phương án còn lại sai.
}
\end{ex}

% ===== Câu 50 =====
% ID: L12C3B18-109-TN
% Mức: VD
\begin{ex}
% ID: L12C3B18-109-TN
% Mức: VD
Trong một bài thực hành về Truyền tải điện năng đi xa và giảm hao phí trên đường dây?
\choice
{Muốn giảm hao phí phải giảm điện áp truyền tải.}
{Có thể bỏ qua phương, chiều và đơn vị trong mọi trường hợp.}
{Kết quả không phụ thuộc dữ kiện và điều kiện áp dụng.}
{\True Với công suất truyền không đổi, tăng điện áp truyền tải làm giảm dòng điện và giảm hao phí I^2R.}
\loigiai{
Với công suất truyền không đổi, tăng điện áp truyền tải làm giảm dòng điện và giảm hao phí I^2R. Vì vậy phương án $D$ đúng; các phương án còn lại sai.
}
\end{ex}

% ===== Câu 51 =====
% ID: L12C3B18-11-DS
% Mức: TH
\dangbt{Nhận biết cấu tạo, nguyên lí hoạt động của máy biến áp}
\begin{ex}
% ID: L12C3B18-11-DS
% Mức: TH
Xét Nhận biết cấu tạo, nguyên lí hoạt động của máy biến áp (Tình huống 1). Hãy xác định tính đúng, sai của bốn phát biểu sau.
\choiceTF
{\True Máy biến áp gồm lõi sắt và hai cuộn dây, hoạt động dựa trên cảm ứng điện từ với dòng điện xoay chiều.}
{Máy biến áp lí tưởng hoạt động được với nguồn điện một chiều không đổi.}
{\True Khi giải cần xác định đúng dữ kiện, phương chiều, điều kiện áp dụng và đơn vị.}
{Kết quả không phụ thuộc dữ kiện và có thể bỏ qua điều kiện.}
\loigiai{
\begin{itemchoice}
\item (Đúng) Máy biến áp gồm lõi sắt và hai cuộn dây, hoạt động dựa trên cảm ứng điện từ với dòng điện xoay chiều. (Vì): Máy biến áp gồm lõi sắt và hai cuộn dây, hoạt động dựa trên cảm ứng điện từ với dòng điện xoay chiều. b) Sai: Máy biến áp lí tưởng hoạt động được với nguồn điện một chiều không đổi. c) Đúng: Khi giải cần xác định đúng dữ kiện, phương chiều, điều kiện áp dụng và đơn vị. d) Sai: Kết quả không phụ thuộc dữ kiện và có thể bỏ qua điều kiện. Kiến thức cốt lõi: Máy biến áp gồm lõi sắt và hai cuộn dây, hoạt động dựa trên cảm ứng điện từ với dòng điện xoay chiều
\item (Sai) Máy biến áp lí tưởng hoạt động được với nguồn điện một chiều không đổi.
\item (Đúng) Khi giải cần xác định đúng dữ kiện, phương chiều, điều kiện áp dụng và đơn vị.
\item (Sai) Kết quả không phụ thuộc dữ kiện và có thể bỏ qua điều kiện.
\end{itemchoice}
}
\end{ex}

% ===== Câu 52 =====
% ID: L12C3B18-11-TL
% Mức: VD
\dangbt{Giải thích nguyên lí bếp từ, phanh điện từ và gia nhiệt cảm ứng}
\begin{bt}
% ID: L12C3B18-11-TL
% Mức: VD
Nêu một tình huống thực tiễn liên quan đến Giải thích nguyên lí bếp từ, phanh điện từ và gia nhiệt cảm ứng và giải thích bằng kiến thức Vật lí.
\loigiai{
Cần nêu được: Bếp từ làm nóng đáy nồi dẫn điện nhờ dòng điện Foucault và sự tỏa nhiệt Joule. Khi vận dụng phải xác định đúng dữ kiện, phương chiều nếu có, điều kiện áp dụng, hệ đơn vị và kiểm tra tính hợp lí. Nhận định “Bếp từ làm nóng trực tiếp mọi vật liệu bằng điện trường tĩnh.” sai.
}
\end{bt}

% ===== Câu 53 =====
% ID: L12C3B18-11-TLN
% Mức: VD
\begin{ex}
% ID: L12C3B18-11-TLN
% Mức: VD
Trong bốn phát biểu về Giải thích nguyên lí bếp từ, phanh điện từ và gia nhiệt cảm ứng, có bao nhiêu phát biểu đúng? (Chỉ ghi một số nguyên.) (1) Bếp từ làm nóng đáy nồi dẫn điện nhờ dòng điện Foucault và sự tỏa nhiệt Joule. (2) Bếp từ làm nóng trực tiếp mọi vật liệu bằng điện trường tĩnh. (3) Cần kiểm tra điều kiện áp dụng. (4) Có thể bỏ qua đơn vị.
\shortans{2}
\loigiai{
Đối chiếu với kiến thức: Bếp từ làm nóng đáy nồi dẫn điện nhờ dòng điện Foucault và sự tỏa nhiệt Joule. Có 2 phát biểu đúng.
}
\end{ex}

% ===== Câu 54 =====
% ID: L12C3B18-11-TN
% Mức: NB
\begin{ex}
% ID: L12C3B18-11-TN
% Mức: NB
Phát biểu nào sau đây đúng về Giải thích nguyên lí bếp từ, phanh điện từ và gia nhiệt cảm ứng?
\choice
{Bếp từ làm nóng trực tiếp mọi vật liệu bằng điện trường tĩnh.}
{Có thể bỏ qua phương, chiều và đơn vị trong mọi trường hợp.}
{Kết quả không phụ thuộc dữ kiện và điều kiện áp dụng.}
{\True Bếp từ làm nóng đáy nồi dẫn điện nhờ dòng điện Foucault và sự tỏa nhiệt Joule.}
\loigiai{
Bếp từ làm nóng đáy nồi dẫn điện nhờ dòng điện Foucault và sự tỏa nhiệt Joule. Vì vậy phương án $D$ đúng; các phương án còn lại sai.
}
\end{ex}

% ===== Câu 55 =====
% ID: L12C3B18-110-TN
% Mức: VD
\dangbt{Truyền tải điện năng đi xa và giảm hao phí trên đường dây}
\begin{ex}
% ID: L12C3B18-110-TN
% Mức: VD
Kết luận đúng khi vận dụng Truyền tải điện năng đi xa và giảm hao phí trên đường dây?
\choice
{\True Với công suất truyền không đổi, tăng điện áp truyền tải làm giảm dòng điện và giảm hao phí I^2R.}
{Muốn giảm hao phí phải giảm điện áp truyền tải.}
{Có thể bỏ qua phương, chiều và đơn vị trong mọi trường hợp.}
{Kết quả không phụ thuộc dữ kiện và điều kiện áp dụng.}
\loigiai{
Với công suất truyền không đổi, tăng điện áp truyền tải làm giảm dòng điện và giảm hao phí I^2R. Vì vậy phương án $A$ đúng; các phương án còn lại sai.
}
\end{ex}

% ===== Câu 56 =====
% ID: L12C3B18-111-TN
% Mức: NB
\dangbt{Bài toán tổng hợp về máy biến áp, truyền tải điện và an toàn điện}
\begin{ex}
% ID: L12C3B18-111-TN
% Mức: NB
Phát biểu nào sau đây đúng về Bài toán tổng hợp về máy biến áp, truyền tải điện và an toàn điện?
\choice
{Có thể bỏ qua nối đất và thiết bị bảo vệ khi dùng điện áp cao.}
{Có thể bỏ qua phương, chiều và đơn vị trong mọi trường hợp.}
{Kết quả không phụ thuộc dữ kiện và điều kiện áp dụng.}
{\True An toàn điện yêu cầu cách điện, nối đất phù hợp và không chạm phần dẫn điện đang có điện áp nguy hiểm.}
\loigiai{
An toàn điện yêu cầu cách điện, nối đất phù hợp và không chạm phần dẫn điện đang có điện áp nguy hiểm. Vì vậy phương án $D$ đúng; các phương án còn lại sai.
}
\end{ex}

% ===== Câu 57 =====
% ID: L12C3B18-112-TN
% Mức: NB
\begin{ex}
% ID: L12C3B18-112-TN
% Mức: NB
Trong nội dung Bài toán tổng hợp về máy biến áp, truyền tải điện và an toàn điện?
\choice
{\True An toàn điện yêu cầu cách điện, nối đất phù hợp và không chạm phần dẫn điện đang có điện áp nguy hiểm.}
{Có thể bỏ qua nối đất và thiết bị bảo vệ khi dùng điện áp cao.}
{Có thể bỏ qua phương, chiều và đơn vị trong mọi trường hợp.}
{Kết quả không phụ thuộc dữ kiện và điều kiện áp dụng.}
\loigiai{
An toàn điện yêu cầu cách điện, nối đất phù hợp và không chạm phần dẫn điện đang có điện áp nguy hiểm. Vì vậy phương án $A$ đúng; các phương án còn lại sai.
}
\end{ex}

% ===== Câu 58 =====
% ID: L12C3B18-113-TN
% Mức: NB
\begin{ex}
% ID: L12C3B18-113-TN
% Mức: NB
Tình huống 3: chọn kết luận chính xác về Bài toán tổng hợp về máy biến áp, truyền tải điện và an toàn điện?
\choice
{Kết quả không phụ thuộc dữ kiện và điều kiện áp dụng.}
{\True An toàn điện yêu cầu cách điện, nối đất phù hợp và không chạm phần dẫn điện đang có điện áp nguy hiểm.}
{Có thể bỏ qua nối đất và thiết bị bảo vệ khi dùng điện áp cao.}
{Có thể bỏ qua phương, chiều và đơn vị trong mọi trường hợp.}
\loigiai{
An toàn điện yêu cầu cách điện, nối đất phù hợp và không chạm phần dẫn điện đang có điện áp nguy hiểm. Vì vậy phương án $B$ đúng; các phương án còn lại sai.
}
\end{ex}

% ===== Câu 59 =====
% ID: L12C3B18-114-TN
% Mức: TH
\begin{ex}
% ID: L12C3B18-114-TN
% Mức: TH
Nội dung nào mô tả đúng nhất Bài toán tổng hợp về máy biến áp, truyền tải điện và an toàn điện?
\choice
{Có thể bỏ qua phương, chiều và đơn vị trong mọi trường hợp.}
{Kết quả không phụ thuộc dữ kiện và điều kiện áp dụng.}
{\True An toàn điện yêu cầu cách điện, nối đất phù hợp và không chạm phần dẫn điện đang có điện áp nguy hiểm.}
{Có thể bỏ qua nối đất và thiết bị bảo vệ khi dùng điện áp cao.}
\loigiai{
An toàn điện yêu cầu cách điện, nối đất phù hợp và không chạm phần dẫn điện đang có điện áp nguy hiểm. Vì vậy phương án $C$ đúng; các phương án còn lại sai.
}
\end{ex}

% ===== Câu 60 =====
% ID: L12C3B18-115-TN
% Mức: TH
\begin{ex}
% ID: L12C3B18-115-TN
% Mức: TH
Khi phân tích Bài toán tổng hợp về máy biến áp, truyền tải điện và an toàn điện?
\choice
{Có thể bỏ qua nối đất và thiết bị bảo vệ khi dùng điện áp cao.}
{Có thể bỏ qua phương, chiều và đơn vị trong mọi trường hợp.}
{Kết quả không phụ thuộc dữ kiện và điều kiện áp dụng.}
{\True An toàn điện yêu cầu cách điện, nối đất phù hợp và không chạm phần dẫn điện đang có điện áp nguy hiểm.}
\loigiai{
An toàn điện yêu cầu cách điện, nối đất phù hợp và không chạm phần dẫn điện đang có điện áp nguy hiểm. Vì vậy phương án $D$ đúng; các phương án còn lại sai.
}
\end{ex}

% ===== Câu 61 =====
% ID: L12C3B18-116-TN
% Mức: TH
\begin{ex}
% ID: L12C3B18-116-TN
% Mức: TH
Một học sinh ôn tập Bài toán tổng hợp về máy biến áp, truyền tải điện và an toàn điện?
\choice
{\True An toàn điện yêu cầu cách điện, nối đất phù hợp và không chạm phần dẫn điện đang có điện áp nguy hiểm.}
{Có thể bỏ qua nối đất và thiết bị bảo vệ khi dùng điện áp cao.}
{Có thể bỏ qua phương, chiều và đơn vị trong mọi trường hợp.}
{Kết quả không phụ thuộc dữ kiện và điều kiện áp dụng.}
\loigiai{
An toàn điện yêu cầu cách điện, nối đất phù hợp và không chạm phần dẫn điện đang có điện áp nguy hiểm. Vì vậy phương án $A$ đúng; các phương án còn lại sai.
}
\end{ex}

% ===== Câu 62 =====
% ID: L12C3B18-117-TN
% Mức: TH
\begin{ex}
% ID: L12C3B18-117-TN
% Mức: TH
Chọn phát biểu không trái với kiến thức về Bài toán tổng hợp về máy biến áp, truyền tải điện và an toàn điện?
\choice
{Kết quả không phụ thuộc dữ kiện và điều kiện áp dụng.}
{\True An toàn điện yêu cầu cách điện, nối đất phù hợp và không chạm phần dẫn điện đang có điện áp nguy hiểm.}
{Có thể bỏ qua nối đất và thiết bị bảo vệ khi dùng điện áp cao.}
{Có thể bỏ qua phương, chiều và đơn vị trong mọi trường hợp.}
\loigiai{
An toàn điện yêu cầu cách điện, nối đất phù hợp và không chạm phần dẫn điện đang có điện áp nguy hiểm. Vì vậy phương án $B$ đúng; các phương án còn lại sai.
}
\end{ex}

% ===== Câu 63 =====
% ID: L12C3B18-118-TN
% Mức: VD
\begin{ex}
% ID: L12C3B18-118-TN
% Mức: VD
Cơ sở Vật lí đúng dùng để giải Bài toán tổng hợp về máy biến áp, truyền tải điện và an toàn điện?
\choice
{Có thể bỏ qua phương, chiều và đơn vị trong mọi trường hợp.}
{Kết quả không phụ thuộc dữ kiện và điều kiện áp dụng.}
{\True An toàn điện yêu cầu cách điện, nối đất phù hợp và không chạm phần dẫn điện đang có điện áp nguy hiểm.}
{Có thể bỏ qua nối đất và thiết bị bảo vệ khi dùng điện áp cao.}
\loigiai{
An toàn điện yêu cầu cách điện, nối đất phù hợp và không chạm phần dẫn điện đang có điện áp nguy hiểm. Vì vậy phương án $C$ đúng; các phương án còn lại sai.
}
\end{ex}

% ===== Câu 64 =====
% ID: L12C3B18-119-TN
% Mức: VD
\begin{ex}
% ID: L12C3B18-119-TN
% Mức: VD
Trong một bài thực hành về Bài toán tổng hợp về máy biến áp, truyền tải điện và an toàn điện?
\choice
{Có thể bỏ qua nối đất và thiết bị bảo vệ khi dùng điện áp cao.}
{Có thể bỏ qua phương, chiều và đơn vị trong mọi trường hợp.}
{Kết quả không phụ thuộc dữ kiện và điều kiện áp dụng.}
{\True An toàn điện yêu cầu cách điện, nối đất phù hợp và không chạm phần dẫn điện đang có điện áp nguy hiểm.}
\loigiai{
An toàn điện yêu cầu cách điện, nối đất phù hợp và không chạm phần dẫn điện đang có điện áp nguy hiểm. Vì vậy phương án $D$ đúng; các phương án còn lại sai.
}
\end{ex}

% ===== Câu 65 =====
% ID: L12C3B18-12-DS
% Mức: TH
\dangbt{Nhận biết cấu tạo, nguyên lí hoạt động của máy biến áp}
\begin{ex}
% ID: L12C3B18-12-DS
% Mức: TH
Xét Nhận biết cấu tạo, nguyên lí hoạt động của máy biến áp (Tình huống 2). Hãy xác định tính đúng, sai của bốn phát biểu sau.
\choiceTF
{Máy biến áp lí tưởng hoạt động được với nguồn điện một chiều không đổi.}
{\True Máy biến áp gồm lõi sắt và hai cuộn dây, hoạt động dựa trên cảm ứng điện từ với dòng điện xoay chiều.}
{Kết quả không phụ thuộc dữ kiện và có thể bỏ qua điều kiện.}
{\True Khi giải cần xác định đúng dữ kiện, phương chiều, điều kiện áp dụng và đơn vị.}
\loigiai{
\begin{itemchoice}
\item (Sai) Máy biến áp lí tưởng hoạt động được với nguồn điện một chiều không đổi. (Vì): Máy biến áp lí tưởng hoạt động được với nguồn điện một chiều không đổi. b) Đúng: Máy biến áp gồm lõi sắt và hai cuộn dây, hoạt động dựa trên cảm ứng điện từ với dòng điện xoay chiều. c) Sai: Kết quả không phụ thuộc dữ kiện và có thể bỏ qua điều kiện. d) Đúng: Khi giải cần xác định đúng dữ kiện, phương chiều, điều kiện áp dụng và đơn vị. Kiến thức cốt lõi: Máy biến áp gồm lõi sắt và hai cuộn dây, hoạt động dựa trên cảm ứng điện từ với dòng điện xoay chiều
\item (Đúng) Máy biến áp gồm lõi sắt và hai cuộn dây, hoạt động dựa trên cảm ứng điện từ với dòng điện xoay chiều.
\item (Sai) Kết quả không phụ thuộc dữ kiện và có thể bỏ qua điều kiện.
\item (Đúng) Khi giải cần xác định đúng dữ kiện, phương chiều, điều kiện áp dụng và đơn vị.
\end{itemchoice}
}
\end{ex}

% ===== Câu 66 =====
% ID: L12C3B18-12-TL
% Mức: VDC
\dangbt{Giải thích nguyên lí bếp từ, phanh điện từ và gia nhiệt cảm ứng}
\begin{bt}
% ID: L12C3B18-12-TL
% Mức: VDC
So sánh nhận định đúng “Bếp từ làm nóng đáy nồi dẫn điện nhờ dòng điện Foucault và sự tỏa nhiệt Joule.” với nhận định sai “Bếp từ làm nóng trực tiếp mọi vật liệu bằng điện trường tĩnh.”; chỉ rõ căn cứ Vật lí.
\loigiai{
Cần nêu được: Bếp từ làm nóng đáy nồi dẫn điện nhờ dòng điện Foucault và sự tỏa nhiệt Joule. Khi vận dụng phải xác định đúng dữ kiện, phương chiều nếu có, điều kiện áp dụng, hệ đơn vị và kiểm tra tính hợp lí. Nhận định “Bếp từ làm nóng trực tiếp mọi vật liệu bằng điện trường tĩnh.” sai.
}
\end{bt}

% ===== Câu 67 =====
% ID: L12C3B18-12-TLN
% Mức: VDC
\begin{ex}
% ID: L12C3B18-12-TLN
% Mức: VDC
Trong bốn phát biểu về Giải thích nguyên lí bếp từ, phanh điện từ và gia nhiệt cảm ứng, có bao nhiêu phát biểu đúng? (Chỉ ghi một số nguyên.) (1) Khi giải cần xác định đúng dữ kiện, phương chiều, điều kiện áp dụng và đơn vị. (2) Bếp từ làm nóng trực tiếp mọi vật liệu bằng điện trường tĩnh. (3) Phải dùng đúng định luật cho tình huống. (4) Bếp từ làm nóng đáy nồi dẫn điện nhờ dòng điện Foucault và sự tỏa nhiệt Joule.
\shortans{3}
\loigiai{
Đối chiếu với kiến thức: Bếp từ làm nóng đáy nồi dẫn điện nhờ dòng điện Foucault và sự tỏa nhiệt Joule. Có 3 phát biểu đúng.
}
\end{ex}

% ===== Câu 68 =====
% ID: L12C3B18-12-TN
% Mức: NB
\begin{ex}
% ID: L12C3B18-12-TN
% Mức: NB
Trong nội dung Giải thích nguyên lí bếp từ, phanh điện từ và gia nhiệt cảm ứng?
\choice
{\True Bếp từ làm nóng đáy nồi dẫn điện nhờ dòng điện Foucault và sự tỏa nhiệt Joule.}
{Bếp từ làm nóng trực tiếp mọi vật liệu bằng điện trường tĩnh.}
{Có thể bỏ qua phương, chiều và đơn vị trong mọi trường hợp.}
{Kết quả không phụ thuộc dữ kiện và điều kiện áp dụng.}
\loigiai{
Bếp từ làm nóng đáy nồi dẫn điện nhờ dòng điện Foucault và sự tỏa nhiệt Joule. Vì vậy phương án $A$ đúng; các phương án còn lại sai.
}
\end{ex}

% ===== Câu 69 =====
% ID: L12C3B18-120-TN
% Mức: VD
\dangbt{Bài toán tổng hợp về máy biến áp, truyền tải điện và an toàn điện}
\begin{ex}
% ID: L12C3B18-120-TN
% Mức: VD
Kết luận đúng khi vận dụng Bài toán tổng hợp về máy biến áp, truyền tải điện và an toàn điện?
\choice
{\True An toàn điện yêu cầu cách điện, nối đất phù hợp và không chạm phần dẫn điện đang có điện áp nguy hiểm.}
{Có thể bỏ qua nối đất và thiết bị bảo vệ khi dùng điện áp cao.}
{Có thể bỏ qua phương, chiều và đơn vị trong mọi trường hợp.}
{Kết quả không phụ thuộc dữ kiện và điều kiện áp dụng.}
\loigiai{
An toàn điện yêu cầu cách điện, nối đất phù hợp và không chạm phần dẫn điện đang có điện áp nguy hiểm. Vì vậy phương án $A$ đúng; các phương án còn lại sai.
}
\end{ex}

% ===== Câu 70 =====
% ID: L12C3B18-13-DS
% Mức: VD
\dangbt{Nhận biết cấu tạo, nguyên lí hoạt động của máy biến áp}
\begin{ex}
% ID: L12C3B18-13-DS
% Mức: VD
Xét Nhận biết cấu tạo, nguyên lí hoạt động của máy biến áp (Tình huống 3). Hãy xác định tính đúng, sai của bốn phát biểu sau.
\choiceTF
{\True Khi giải cần xác định đúng dữ kiện, phương chiều, điều kiện áp dụng và đơn vị.}
{Kết quả không phụ thuộc dữ kiện và có thể bỏ qua điều kiện.}
{\True Máy biến áp gồm lõi sắt và hai cuộn dây, hoạt động dựa trên cảm ứng điện từ với dòng điện xoay chiều.}
{Máy biến áp lí tưởng hoạt động được với nguồn điện một chiều không đổi.}
\loigiai{
\begin{itemchoice}
\item (Đúng) Khi giải cần xác định đúng dữ kiện, phương chiều, điều kiện áp dụng và đơn vị. (Vì): Khi giải cần xác định đúng dữ kiện, phương chiều, điều kiện áp dụng và đơn vị. b) Sai: Kết quả không phụ thuộc dữ kiện và có thể bỏ qua điều kiện. c) Đúng: Máy biến áp gồm lõi sắt và hai cuộn dây, hoạt động dựa trên cảm ứng điện từ với dòng điện xoay chiều. d) Sai: Máy biến áp lí tưởng hoạt động được với nguồn điện một chiều không đổi. Kiến thức cốt lõi: Máy biến áp gồm lõi sắt và hai cuộn dây, hoạt động dựa trên cảm ứng điện từ với dòng điện xoay chiều
\item (Sai) Kết quả không phụ thuộc dữ kiện và có thể bỏ qua điều kiện.
\item (Đúng) Máy biến áp gồm lõi sắt và hai cuộn dây, hoạt động dựa trên cảm ứng điện từ với dòng điện xoay chiều.
\item (Sai) Máy biến áp lí tưởng hoạt động được với nguồn điện một chiều không đổi.
\end{itemchoice}
}
\end{ex}

% ===== Câu 71 =====
% ID: L12C3B18-13-TL
% Mức: TH
\begin{bt}
% ID: L12C3B18-13-TL
% Mức: TH
Trình bày kiến thức cốt lõi của Nhận biết cấu tạo, nguyên lí hoạt động của máy biến áp và nêu điều kiện áp dụng.
\loigiai{
Cần nêu được: Máy biến áp gồm lõi sắt và hai cuộn dây, hoạt động dựa trên cảm ứng điện từ với dòng điện xoay chiều. Khi vận dụng phải xác định đúng dữ kiện, phương chiều nếu có, điều kiện áp dụng, hệ đơn vị và kiểm tra tính hợp lí. Nhận định “Máy biến áp lí tưởng hoạt động được với nguồn điện một chiều không đổi.” sai.
}
\end{bt}

% ===== Câu 72 =====
% ID: L12C3B18-13-TLN
% Mức: TH
\begin{ex}
% ID: L12C3B18-13-TLN
% Mức: TH
Trong bốn phát biểu về Nhận biết cấu tạo, nguyên lí hoạt động của máy biến áp, có bao nhiêu phát biểu đúng? (Chỉ ghi một số nguyên.) (1) Máy biến áp gồm lõi sắt và hai cuộn dây, hoạt động dựa trên cảm ứng điện từ với dòng điện xoay chiều. (2) Máy biến áp lí tưởng hoạt động được với nguồn điện một chiều không đổi. (3) Cần kiểm tra điều kiện áp dụng. (4) Có thể bỏ qua đơn vị.
\shortans{2}
\loigiai{
Đối chiếu với kiến thức: Máy biến áp gồm lõi sắt và hai cuộn dây, hoạt động dựa trên cảm ứng điện từ với dòng điện xoay chiều. Có 2 phát biểu đúng.
}
\end{ex}

% ===== Câu 73 =====
% ID: L12C3B18-13-TN
% Mức: NB
\dangbt{Giải thích nguyên lí bếp từ, phanh điện từ và gia nhiệt cảm ứng}
\begin{ex}
% ID: L12C3B18-13-TN
% Mức: NB
Tình huống 3: chọn kết luận chính xác về Giải thích nguyên lí bếp từ, phanh điện từ và gia nhiệt cảm ứng?
\choice
{Kết quả không phụ thuộc dữ kiện và điều kiện áp dụng.}
{\True Bếp từ làm nóng đáy nồi dẫn điện nhờ dòng điện Foucault và sự tỏa nhiệt Joule.}
{Bếp từ làm nóng trực tiếp mọi vật liệu bằng điện trường tĩnh.}
{Có thể bỏ qua phương, chiều và đơn vị trong mọi trường hợp.}
\loigiai{
Bếp từ làm nóng đáy nồi dẫn điện nhờ dòng điện Foucault và sự tỏa nhiệt Joule. Vì vậy phương án $B$ đúng; các phương án còn lại sai.
}
\end{ex}

% ===== Câu 74 =====
% ID: L12C3B18-14-DS
% Mức: VD
\dangbt{Nhận biết cấu tạo, nguyên lí hoạt động của máy biến áp}
\begin{ex}
% ID: L12C3B18-14-DS
% Mức: VD
Xét Nhận biết cấu tạo, nguyên lí hoạt động của máy biến áp (Tình huống 4). Hãy xác định tính đúng, sai của bốn phát biểu sau.
\choiceTF
{Kết quả không phụ thuộc dữ kiện và có thể bỏ qua điều kiện.}
{\True Khi giải cần xác định đúng dữ kiện, phương chiều, điều kiện áp dụng và đơn vị.}
{Máy biến áp lí tưởng hoạt động được với nguồn điện một chiều không đổi.}
{\True Máy biến áp gồm lõi sắt và hai cuộn dây, hoạt động dựa trên cảm ứng điện từ với dòng điện xoay chiều.}
\loigiai{
\begin{itemchoice}
\item (Sai) Kết quả không phụ thuộc dữ kiện và có thể bỏ qua điều kiện. (Vì): Kết quả không phụ thuộc dữ kiện và có thể bỏ qua điều kiện. b) Đúng: Khi giải cần xác định đúng dữ kiện, phương chiều, điều kiện áp dụng và đơn vị. c) Sai: Máy biến áp lí tưởng hoạt động được với nguồn điện một chiều không đổi. d) Đúng: Máy biến áp gồm lõi sắt và hai cuộn dây, hoạt động dựa trên cảm ứng điện từ với dòng điện xoay chiều. Kiến thức cốt lõi: Máy biến áp gồm lõi sắt và hai cuộn dây, hoạt động dựa trên cảm ứng điện từ với dòng điện xoay chiều
\item (Đúng) Khi giải cần xác định đúng dữ kiện, phương chiều, điều kiện áp dụng và đơn vị.
\item (Sai) Máy biến áp lí tưởng hoạt động được với nguồn điện một chiều không đổi.
\item (Đúng) Máy biến áp gồm lõi sắt và hai cuộn dây, hoạt động dựa trên cảm ứng điện từ với dòng điện xoay chiều.
\end{itemchoice}
}
\end{ex}

% ===== Câu 75 =====
% ID: L12C3B18-14-TL
% Mức: TH
\begin{bt}
% ID: L12C3B18-14-TL
% Mức: TH
Giải thích vì sao nhận định sau đúng: “Máy biến áp gồm lõi sắt và hai cuộn dây, hoạt động dựa trên cảm ứng điện từ với dòng điện xoay chiều.”
\loigiai{
Cần nêu được: Máy biến áp gồm lõi sắt và hai cuộn dây, hoạt động dựa trên cảm ứng điện từ với dòng điện xoay chiều. Khi vận dụng phải xác định đúng dữ kiện, phương chiều nếu có, điều kiện áp dụng, hệ đơn vị và kiểm tra tính hợp lí. Nhận định “Máy biến áp lí tưởng hoạt động được với nguồn điện một chiều không đổi.” sai.
}
\end{bt}

% ===== Câu 76 =====
% ID: L12C3B18-14-TLN
% Mức: TH
\begin{ex}
% ID: L12C3B18-14-TLN
% Mức: TH
Trong bốn phát biểu về Nhận biết cấu tạo, nguyên lí hoạt động của máy biến áp, có bao nhiêu phát biểu đúng? (Chỉ ghi một số nguyên.) (1) Máy biến áp lí tưởng hoạt động được với nguồn điện một chiều không đổi. (2) Máy biến áp gồm lõi sắt và hai cuộn dây, hoạt động dựa trên cảm ứng điện từ với dòng điện xoay chiều. (3) Kết quả phải phù hợp ý nghĩa vật lí. (4) Mọi đại lượng đều là vô hướng.
\shortans{1}
\loigiai{
Đối chiếu với kiến thức: Máy biến áp gồm lõi sắt và hai cuộn dây, hoạt động dựa trên cảm ứng điện từ với dòng điện xoay chiều. Có 1 phát biểu đúng.
}
\end{ex}

% ===== Câu 77 =====
% ID: L12C3B18-14-TN
% Mức: TH
\dangbt{Giải thích nguyên lí bếp từ, phanh điện từ và gia nhiệt cảm ứng}
\begin{ex}
% ID: L12C3B18-14-TN
% Mức: TH
Nội dung nào mô tả đúng nhất Giải thích nguyên lí bếp từ, phanh điện từ và gia nhiệt cảm ứng?
\choice
{Có thể bỏ qua phương, chiều và đơn vị trong mọi trường hợp.}
{Kết quả không phụ thuộc dữ kiện và điều kiện áp dụng.}
{\True Bếp từ làm nóng đáy nồi dẫn điện nhờ dòng điện Foucault và sự tỏa nhiệt Joule.}
{Bếp từ làm nóng trực tiếp mọi vật liệu bằng điện trường tĩnh.}
\loigiai{
Bếp từ làm nóng đáy nồi dẫn điện nhờ dòng điện Foucault và sự tỏa nhiệt Joule. Vì vậy phương án $C$ đúng; các phương án còn lại sai.
}
\end{ex}

% ===== Câu 78 =====
% ID: L12C3B18-15-DS
% Mức: VD
\dangbt{Nhận biết cấu tạo, nguyên lí hoạt động của máy biến áp}
\begin{ex}
% ID: L12C3B18-15-DS
% Mức: VD
Xét Nhận biết cấu tạo, nguyên lí hoạt động của máy biến áp (Tình huống 5). Hãy xác định tính đúng, sai của bốn phát biểu sau.
\choiceTF
{\True Máy biến áp gồm lõi sắt và hai cuộn dây, hoạt động dựa trên cảm ứng điện từ với dòng điện xoay chiều.}
{Kết quả không phụ thuộc dữ kiện và có thể bỏ qua điều kiện.}
{Máy biến áp lí tưởng hoạt động được với nguồn điện một chiều không đổi.}
{\True Khi giải cần xác định đúng dữ kiện, phương chiều, điều kiện áp dụng và đơn vị.}
\loigiai{
\begin{itemchoice}
\item (Đúng) Máy biến áp gồm lõi sắt và hai cuộn dây, hoạt động dựa trên cảm ứng điện từ với dòng điện xoay chiều. (Vì): Máy biến áp gồm lõi sắt và hai cuộn dây, hoạt động dựa trên cảm ứng điện từ với dòng điện xoay chiều. b) Sai: Kết quả không phụ thuộc dữ kiện và có thể bỏ qua điều kiện. c) Sai: Máy biến áp lí tưởng hoạt động được với nguồn điện một chiều không đổi. d) Đúng: Khi giải cần xác định đúng dữ kiện, phương chiều, điều kiện áp dụng và đơn vị. Kiến thức cốt lõi: Máy biến áp gồm lõi sắt và hai cuộn dây, hoạt động dựa trên cảm ứng điện từ với dòng điện xoay chiều
\item (Sai) Kết quả không phụ thuộc dữ kiện và có thể bỏ qua điều kiện.
\item (Sai) Máy biến áp lí tưởng hoạt động được với nguồn điện một chiều không đổi.
\item (Đúng) Khi giải cần xác định đúng dữ kiện, phương chiều, điều kiện áp dụng và đơn vị.
\end{itemchoice}
}
\end{ex}

% ===== Câu 79 =====
% ID: L12C3B18-15-TL
% Mức: VD
\begin{bt}
% ID: L12C3B18-15-TL
% Mức: VD
Phân tích sai lầm trong nhận định: “Máy biến áp lí tưởng hoạt động được với nguồn điện một chiều không đổi.”
\loigiai{
Cần nêu được: Máy biến áp gồm lõi sắt và hai cuộn dây, hoạt động dựa trên cảm ứng điện từ với dòng điện xoay chiều. Khi vận dụng phải xác định đúng dữ kiện, phương chiều nếu có, điều kiện áp dụng, hệ đơn vị và kiểm tra tính hợp lí. Nhận định “Máy biến áp lí tưởng hoạt động được với nguồn điện một chiều không đổi.” sai.
}
\end{bt}

% ===== Câu 80 =====
% ID: L12C3B18-15-TLN
% Mức: TH
\begin{ex}
% ID: L12C3B18-15-TLN
% Mức: TH
Trong bốn phát biểu về Nhận biết cấu tạo, nguyên lí hoạt động của máy biến áp, có bao nhiêu phát biểu đúng? (Chỉ ghi một số nguyên.) (1) Máy biến áp gồm lõi sắt và hai cuộn dây, hoạt động dựa trên cảm ứng điện từ với dòng điện xoay chiều. (2) Máy biến áp lí tưởng hoạt động được với nguồn điện một chiều không đổi. (3) Cần kiểm tra điều kiện áp dụng. (4) Có thể bỏ qua đơn vị.
\shortans{2}
\loigiai{
Đối chiếu với kiến thức: Máy biến áp gồm lõi sắt và hai cuộn dây, hoạt động dựa trên cảm ứng điện từ với dòng điện xoay chiều. Có 2 phát biểu đúng.
}
\end{ex}

% ===== Câu 81 =====
% ID: L12C3B18-15-TN
% Mức: TH
\dangbt{Giải thích nguyên lí bếp từ, phanh điện từ và gia nhiệt cảm ứng}
\begin{ex}
% ID: L12C3B18-15-TN
% Mức: TH
Khi phân tích Giải thích nguyên lí bếp từ, phanh điện từ và gia nhiệt cảm ứng?
\choice
{Bếp từ làm nóng trực tiếp mọi vật liệu bằng điện trường tĩnh.}
{Có thể bỏ qua phương, chiều và đơn vị trong mọi trường hợp.}
{Kết quả không phụ thuộc dữ kiện và điều kiện áp dụng.}
{\True Bếp từ làm nóng đáy nồi dẫn điện nhờ dòng điện Foucault và sự tỏa nhiệt Joule.}
\loigiai{
Bếp từ làm nóng đáy nồi dẫn điện nhờ dòng điện Foucault và sự tỏa nhiệt Joule. Vì vậy phương án $D$ đúng; các phương án còn lại sai.
}
\end{ex}

% ===== Câu 82 =====
% ID: L12C3B18-16-DS
% Mức: TH
\dangbt{Nhận biết dòng điện Foucault và các ứng dụng của cảm ứng điện từ}
\begin{ex}
% ID: L12C3B18-16-DS
% Mức: TH
Xét Nhận biết dòng điện Foucault và các ứng dụng của cảm ứng điện từ (Tình huống 1). Hãy xác định tính đúng, sai của bốn phát biểu sau.
\choiceTF
{\True Dòng điện Foucault là dòng điện cảm ứng xuất hiện trong khối vật dẫn khi từ thông qua nó biến thiên.}
{Dòng điện Foucault chỉ xuất hiện trong chất cách điện.}
{\True Khi giải cần xác định đúng dữ kiện, phương chiều, điều kiện áp dụng và đơn vị.}
{Kết quả không phụ thuộc dữ kiện và có thể bỏ qua điều kiện.}
\loigiai{
\begin{itemchoice}
\item (Đúng) Dòng điện Foucault là dòng điện cảm ứng xuất hiện trong khối vật dẫn khi từ thông qua nó biến thiên. (Vì): òng điện Foucault là dòng điện cảm ứng xuất hiện trong khối vật dẫn khi từ thông qua nó biến thiên. b) Sai: Dòng điện Foucault chỉ xuất hiện trong chất cách điện. c) Đúng: Khi giải cần xác định đúng dữ kiện, phương chiều, điều kiện áp dụng và đơn vị. d) Sai: Kết quả không phụ thuộc dữ kiện và có thể bỏ qua điều kiện. Kiến thức cốt lõi: Dòng điện Foucault là dòng điện cảm ứng xuất hiện trong khối vật dẫn khi từ thông qua nó biến thiên
\item (Sai) Dòng điện Foucault chỉ xuất hiện trong chất cách điện.
\item (Đúng) Khi giải cần xác định đúng dữ kiện, phương chiều, điều kiện áp dụng và đơn vị.
\item (Sai) Kết quả không phụ thuộc dữ kiện và có thể bỏ qua điều kiện.
\end{itemchoice}
}
\end{ex}

% ===== Câu 83 =====
% ID: L12C3B18-16-TL
% Mức: VD
\dangbt{Nhận biết cấu tạo, nguyên lí hoạt động của máy biến áp}
\begin{bt}
% ID: L12C3B18-16-TL
% Mức: VD
Đề xuất các bước giải một bài toán thuộc dạng Nhận biết cấu tạo, nguyên lí hoạt động của máy biến áp.
\loigiai{
Cần nêu được: Máy biến áp gồm lõi sắt và hai cuộn dây, hoạt động dựa trên cảm ứng điện từ với dòng điện xoay chiều. Khi vận dụng phải xác định đúng dữ kiện, phương chiều nếu có, điều kiện áp dụng, hệ đơn vị và kiểm tra tính hợp lí. Nhận định “Máy biến áp lí tưởng hoạt động được với nguồn điện một chiều không đổi.” sai.
}
\end{bt}

% ===== Câu 84 =====
% ID: L12C3B18-16-TLN
% Mức: VD
\begin{ex}
% ID: L12C3B18-16-TLN
% Mức: VD
Trong bốn phát biểu về Nhận biết cấu tạo, nguyên lí hoạt động của máy biến áp, có bao nhiêu phát biểu đúng? (Chỉ ghi một số nguyên.) (1) Máy biến áp lí tưởng hoạt động được với nguồn điện một chiều không đổi. (2) Máy biến áp gồm lõi sắt và hai cuộn dây, hoạt động dựa trên cảm ứng điện từ với dòng điện xoay chiều. (3) Kết quả phải phù hợp ý nghĩa vật lí. (4) Mọi đại lượng đều là vô hướng.
\shortans{2}
\loigiai{
Đối chiếu với kiến thức: Máy biến áp gồm lõi sắt và hai cuộn dây, hoạt động dựa trên cảm ứng điện từ với dòng điện xoay chiều. Có 2 phát biểu đúng.
}
\end{ex}

% ===== Câu 85 =====
% ID: L12C3B18-16-TN
% Mức: TH
\dangbt{Giải thích nguyên lí bếp từ, phanh điện từ và gia nhiệt cảm ứng}
\begin{ex}
% ID: L12C3B18-16-TN
% Mức: TH
Một học sinh ôn tập Giải thích nguyên lí bếp từ, phanh điện từ và gia nhiệt cảm ứng?
\choice
{\True Bếp từ làm nóng đáy nồi dẫn điện nhờ dòng điện Foucault và sự tỏa nhiệt Joule.}
{Bếp từ làm nóng trực tiếp mọi vật liệu bằng điện trường tĩnh.}
{Có thể bỏ qua phương, chiều và đơn vị trong mọi trường hợp.}
{Kết quả không phụ thuộc dữ kiện và điều kiện áp dụng.}
\loigiai{
Bếp từ làm nóng đáy nồi dẫn điện nhờ dòng điện Foucault và sự tỏa nhiệt Joule. Vì vậy phương án $A$ đúng; các phương án còn lại sai.
}
\end{ex}

% ===== Câu 86 =====
% ID: L12C3B18-17-DS
% Mức: TH
\dangbt{Nhận biết dòng điện Foucault và các ứng dụng của cảm ứng điện từ}
\begin{ex}
% ID: L12C3B18-17-DS
% Mức: TH
Xét Nhận biết dòng điện Foucault và các ứng dụng của cảm ứng điện từ (Tình huống 2). Hãy xác định tính đúng, sai của bốn phát biểu sau.
\choiceTF
{Dòng điện Foucault chỉ xuất hiện trong chất cách điện.}
{\True Dòng điện Foucault là dòng điện cảm ứng xuất hiện trong khối vật dẫn khi từ thông qua nó biến thiên.}
{Kết quả không phụ thuộc dữ kiện và có thể bỏ qua điều kiện.}
{\True Khi giải cần xác định đúng dữ kiện, phương chiều, điều kiện áp dụng và đơn vị.}
\loigiai{
\begin{itemchoice}
\item (Sai) Dòng điện Foucault chỉ xuất hiện trong chất cách điện. (Vì): Dòng điện Foucault chỉ xuất hiện trong chất cách điện. b) Đúng: Dòng điện Foucault là dòng điện cảm ứng xuất hiện trong khối vật dẫn khi từ thông qua nó biến thiên. c) Sai: Kết quả không phụ thuộc dữ kiện và có thể bỏ qua điều kiện. d) Đúng: Khi giải cần xác định đúng dữ kiện, phương chiều, điều kiện áp dụng và đơn vị. Kiến thức cốt lõi: Dòng điện Foucault là dòng điện cảm ứng xuất hiện trong khối vật dẫn khi từ thông qua nó biến thiên
\item (Đúng) Dòng điện Foucault là dòng điện cảm ứng xuất hiện trong khối vật dẫn khi từ thông qua nó biến thiên.
\item (Sai) Kết quả không phụ thuộc dữ kiện và có thể bỏ qua điều kiện.
\item (Đúng) Khi giải cần xác định đúng dữ kiện, phương chiều, điều kiện áp dụng và đơn vị.
\end{itemchoice}
}
\end{ex}

% ===== Câu 87 =====
% ID: L12C3B18-17-TL
% Mức: VD
\dangbt{Nhận biết cấu tạo, nguyên lí hoạt động của máy biến áp}
\begin{bt}
% ID: L12C3B18-17-TL
% Mức: VD
Nêu một tình huống thực tiễn liên quan đến Nhận biết cấu tạo, nguyên lí hoạt động của máy biến áp và giải thích bằng kiến thức Vật lí.
\loigiai{
Cần nêu được: Máy biến áp gồm lõi sắt và hai cuộn dây, hoạt động dựa trên cảm ứng điện từ với dòng điện xoay chiều. Khi vận dụng phải xác định đúng dữ kiện, phương chiều nếu có, điều kiện áp dụng, hệ đơn vị và kiểm tra tính hợp lí. Nhận định “Máy biến áp lí tưởng hoạt động được với nguồn điện một chiều không đổi.” sai.
}
\end{bt}

% ===== Câu 88 =====
% ID: L12C3B18-17-TLN
% Mức: VD
\begin{ex}
% ID: L12C3B18-17-TLN
% Mức: VD
Trong bốn phát biểu về Nhận biết cấu tạo, nguyên lí hoạt động của máy biến áp, có bao nhiêu phát biểu đúng? (Chỉ ghi một số nguyên.) (1) Máy biến áp gồm lõi sắt và hai cuộn dây, hoạt động dựa trên cảm ứng điện từ với dòng điện xoay chiều. (2) Máy biến áp lí tưởng hoạt động được với nguồn điện một chiều không đổi. (3) Cần kiểm tra điều kiện áp dụng. (4) Có thể bỏ qua đơn vị.
\shortans{2}
\loigiai{
Đối chiếu với kiến thức: Máy biến áp gồm lõi sắt và hai cuộn dây, hoạt động dựa trên cảm ứng điện từ với dòng điện xoay chiều. Có 2 phát biểu đúng.
}
\end{ex}

% ===== Câu 89 =====
% ID: L12C3B18-17-TN
% Mức: TH
\dangbt{Giải thích nguyên lí bếp từ, phanh điện từ và gia nhiệt cảm ứng}
\begin{ex}
% ID: L12C3B18-17-TN
% Mức: TH
Chọn phát biểu không trái với kiến thức về Giải thích nguyên lí bếp từ, phanh điện từ và gia nhiệt cảm ứng?
\choice
{Kết quả không phụ thuộc dữ kiện và điều kiện áp dụng.}
{\True Bếp từ làm nóng đáy nồi dẫn điện nhờ dòng điện Foucault và sự tỏa nhiệt Joule.}
{Bếp từ làm nóng trực tiếp mọi vật liệu bằng điện trường tĩnh.}
{Có thể bỏ qua phương, chiều và đơn vị trong mọi trường hợp.}
\loigiai{
Bếp từ làm nóng đáy nồi dẫn điện nhờ dòng điện Foucault và sự tỏa nhiệt Joule. Vì vậy phương án $B$ đúng; các phương án còn lại sai.
}
\end{ex}

% ===== Câu 90 =====
% ID: L12C3B18-18-DS
% Mức: VD
\dangbt{Nhận biết dòng điện Foucault và các ứng dụng của cảm ứng điện từ}
\begin{ex}
% ID: L12C3B18-18-DS
% Mức: VD
Xét Nhận biết dòng điện Foucault và các ứng dụng của cảm ứng điện từ (Tình huống 3). Hãy xác định tính đúng, sai của bốn phát biểu sau.
\choiceTF
{\True Khi giải cần xác định đúng dữ kiện, phương chiều, điều kiện áp dụng và đơn vị.}
{Kết quả không phụ thuộc dữ kiện và có thể bỏ qua điều kiện.}
{\True Dòng điện Foucault là dòng điện cảm ứng xuất hiện trong khối vật dẫn khi từ thông qua nó biến thiên.}
{Dòng điện Foucault chỉ xuất hiện trong chất cách điện.}
\loigiai{
\begin{itemchoice}
\item (Đúng) Khi giải cần xác định đúng dữ kiện, phương chiều, điều kiện áp dụng và đơn vị. (Vì): Khi giải cần xác định đúng dữ kiện, phương chiều, điều kiện áp dụng và đơn vị. b) Sai: Kết quả không phụ thuộc dữ kiện và có thể bỏ qua điều kiện. c) Đúng: Dòng điện Foucault là dòng điện cảm ứng xuất hiện trong khối vật dẫn khi từ thông qua nó biến thiên. d) Sai: Dòng điện Foucault chỉ xuất hiện trong chất cách điện. Kiến thức cốt lõi: Dòng điện Foucault là dòng điện cảm ứng xuất hiện trong khối vật dẫn khi từ thông qua nó biến thiên
\item (Sai) Kết quả không phụ thuộc dữ kiện và có thể bỏ qua điều kiện.
\item (Đúng) Dòng điện Foucault là dòng điện cảm ứng xuất hiện trong khối vật dẫn khi từ thông qua nó biến thiên.
\item (Sai) Dòng điện Foucault chỉ xuất hiện trong chất cách điện.
\end{itemchoice}
}
\end{ex}

% ===== Câu 91 =====
% ID: L12C3B18-18-TL
% Mức: VDC
\dangbt{Nhận biết cấu tạo, nguyên lí hoạt động của máy biến áp}
\begin{bt}
% ID: L12C3B18-18-TL
% Mức: VDC
So sánh nhận định đúng “Máy biến áp gồm lõi sắt và hai cuộn dây, hoạt động dựa trên cảm ứng điện từ với dòng điện xoay chiều.” với nhận định sai “Máy biến áp lí tưởng hoạt động được với nguồn điện một chiều không đổi.”; chỉ rõ căn cứ Vật lí.
\loigiai{
Cần nêu được: Máy biến áp gồm lõi sắt và hai cuộn dây, hoạt động dựa trên cảm ứng điện từ với dòng điện xoay chiều. Khi vận dụng phải xác định đúng dữ kiện, phương chiều nếu có, điều kiện áp dụng, hệ đơn vị và kiểm tra tính hợp lí. Nhận định “Máy biến áp lí tưởng hoạt động được với nguồn điện một chiều không đổi.” sai.
}
\end{bt}

% ===== Câu 92 =====
% ID: L12C3B18-18-TLN
% Mức: VDC
\begin{ex}
% ID: L12C3B18-18-TLN
% Mức: VDC
Trong bốn phát biểu về Nhận biết cấu tạo, nguyên lí hoạt động của máy biến áp, có bao nhiêu phát biểu đúng? (Chỉ ghi một số nguyên.) (1) Khi giải cần xác định đúng dữ kiện, phương chiều, điều kiện áp dụng và đơn vị. (2) Máy biến áp lí tưởng hoạt động được với nguồn điện một chiều không đổi. (3) Phải dùng đúng định luật cho tình huống. (4) Máy biến áp gồm lõi sắt và hai cuộn dây, hoạt động dựa trên cảm ứng điện từ với dòng điện xoay chiều.
\shortans{3}
\loigiai{
Đối chiếu với kiến thức: Máy biến áp gồm lõi sắt và hai cuộn dây, hoạt động dựa trên cảm ứng điện từ với dòng điện xoay chiều. Có 3 phát biểu đúng.
}
\end{ex}

% ===== Câu 93 =====
% ID: L12C3B18-18-TN
% Mức: VD
\dangbt{Giải thích nguyên lí bếp từ, phanh điện từ và gia nhiệt cảm ứng}
\begin{ex}
% ID: L12C3B18-18-TN
% Mức: VD
Cơ sở Vật lí đúng dùng để giải Giải thích nguyên lí bếp từ, phanh điện từ và gia nhiệt cảm ứng?
\choice
{Có thể bỏ qua phương, chiều và đơn vị trong mọi trường hợp.}
{Kết quả không phụ thuộc dữ kiện và điều kiện áp dụng.}
{\True Bếp từ làm nóng đáy nồi dẫn điện nhờ dòng điện Foucault và sự tỏa nhiệt Joule.}
{Bếp từ làm nóng trực tiếp mọi vật liệu bằng điện trường tĩnh.}
\loigiai{
Bếp từ làm nóng đáy nồi dẫn điện nhờ dòng điện Foucault và sự tỏa nhiệt Joule. Vì vậy phương án $C$ đúng; các phương án còn lại sai.
}
\end{ex}

% ===== Câu 94 =====
% ID: L12C3B18-19-DS
% Mức: VD
\dangbt{Nhận biết dòng điện Foucault và các ứng dụng của cảm ứng điện từ}
\begin{ex}
% ID: L12C3B18-19-DS
% Mức: VD
Xét Nhận biết dòng điện Foucault và các ứng dụng của cảm ứng điện từ (Tình huống 4). Hãy xác định tính đúng, sai của bốn phát biểu sau.
\choiceTF
{Kết quả không phụ thuộc dữ kiện và có thể bỏ qua điều kiện.}
{\True Khi giải cần xác định đúng dữ kiện, phương chiều, điều kiện áp dụng và đơn vị.}
{Dòng điện Foucault chỉ xuất hiện trong chất cách điện.}
{\True Dòng điện Foucault là dòng điện cảm ứng xuất hiện trong khối vật dẫn khi từ thông qua nó biến thiên.}
\loigiai{
\begin{itemchoice}
\item (Sai) Kết quả không phụ thuộc dữ kiện và có thể bỏ qua điều kiện. (Vì): Kết quả không phụ thuộc dữ kiện và có thể bỏ qua điều kiện. b) Đúng: Khi giải cần xác định đúng dữ kiện, phương chiều, điều kiện áp dụng và đơn vị. c) Sai: Dòng điện Foucault chỉ xuất hiện trong chất cách điện. d) Đúng: Dòng điện Foucault là dòng điện cảm ứng xuất hiện trong khối vật dẫn khi từ thông qua nó biến thiên. Kiến thức cốt lõi: Dòng điện Foucault là dòng điện cảm ứng xuất hiện trong khối vật dẫn khi từ thông qua nó biến thiên
\item (Đúng) Khi giải cần xác định đúng dữ kiện, phương chiều, điều kiện áp dụng và đơn vị.
\item (Sai) Dòng điện Foucault chỉ xuất hiện trong chất cách điện.
\item (Đúng) Dòng điện Foucault là dòng điện cảm ứng xuất hiện trong khối vật dẫn khi từ thông qua nó biến thiên.
\end{itemchoice}
}
\end{ex}

% ===== Câu 95 =====
% ID: L12C3B18-19-TL
% Mức: TH
\begin{bt}
% ID: L12C3B18-19-TL
% Mức: TH
Trình bày kiến thức cốt lõi của Nhận biết dòng điện Foucault và các ứng dụng của cảm ứng điện từ và nêu điều kiện áp dụng.
\loigiai{
Cần nêu được: Dòng điện Foucault là dòng điện cảm ứng xuất hiện trong khối vật dẫn khi từ thông qua nó biến thiên. Khi vận dụng phải xác định đúng dữ kiện, phương chiều nếu có, điều kiện áp dụng, hệ đơn vị và kiểm tra tính hợp lí. Nhận định “Dòng điện Foucault chỉ xuất hiện trong chất cách điện.” sai.
}
\end{bt}

% ===== Câu 96 =====
% ID: L12C3B18-19-TLN
% Mức: TH
\begin{ex}
% ID: L12C3B18-19-TLN
% Mức: TH
Trong bốn phát biểu về Nhận biết dòng điện Foucault và các ứng dụng của cảm ứng điện từ, có bao nhiêu phát biểu đúng? (Chỉ ghi một số nguyên.) (1) Dòng điện Foucault là dòng điện cảm ứng xuất hiện trong khối vật dẫn khi từ thông qua nó biến thiên. (2) Dòng điện Foucault chỉ xuất hiện trong chất cách điện. (3) Cần kiểm tra điều kiện áp dụng. (4) Có thể bỏ qua đơn vị.
\shortans{2}
\loigiai{
Đối chiếu với kiến thức: Dòng điện Foucault là dòng điện cảm ứng xuất hiện trong khối vật dẫn khi từ thông qua nó biến thiên. Có 2 phát biểu đúng.
}
\end{ex}

% ===== Câu 97 =====
% ID: L12C3B18-19-TN
% Mức: VD
\dangbt{Giải thích nguyên lí bếp từ, phanh điện từ và gia nhiệt cảm ứng}
\begin{ex}
% ID: L12C3B18-19-TN
% Mức: VD
Trong một bài thực hành về Giải thích nguyên lí bếp từ, phanh điện từ và gia nhiệt cảm ứng?
\choice
{Bếp từ làm nóng trực tiếp mọi vật liệu bằng điện trường tĩnh.}
{Có thể bỏ qua phương, chiều và đơn vị trong mọi trường hợp.}
{Kết quả không phụ thuộc dữ kiện và điều kiện áp dụng.}
{\True Bếp từ làm nóng đáy nồi dẫn điện nhờ dòng điện Foucault và sự tỏa nhiệt Joule.}
\loigiai{
Bếp từ làm nóng đáy nồi dẫn điện nhờ dòng điện Foucault và sự tỏa nhiệt Joule. Vì vậy phương án $D$ đúng; các phương án còn lại sai.
}
\end{ex}

% ===== Câu 98 =====
% ID: L12C3B18-20-DS
% Mức: VD
\dangbt{Nhận biết dòng điện Foucault và các ứng dụng của cảm ứng điện từ}
\begin{ex}
% ID: L12C3B18-20-DS
% Mức: VD
Xét Nhận biết dòng điện Foucault và các ứng dụng của cảm ứng điện từ (Tình huống 5). Hãy xác định tính đúng, sai của bốn phát biểu sau.
\choiceTF
{\True Dòng điện Foucault là dòng điện cảm ứng xuất hiện trong khối vật dẫn khi từ thông qua nó biến thiên.}
{Kết quả không phụ thuộc dữ kiện và có thể bỏ qua điều kiện.}
{Dòng điện Foucault chỉ xuất hiện trong chất cách điện.}
{\True Khi giải cần xác định đúng dữ kiện, phương chiều, điều kiện áp dụng và đơn vị.}
\loigiai{
\begin{itemchoice}
\item (Đúng) Dòng điện Foucault là dòng điện cảm ứng xuất hiện trong khối vật dẫn khi từ thông qua nó biến thiên. (Vì): òng điện Foucault là dòng điện cảm ứng xuất hiện trong khối vật dẫn khi từ thông qua nó biến thiên. b) Sai: Kết quả không phụ thuộc dữ kiện và có thể bỏ qua điều kiện. c) Sai: Dòng điện Foucault chỉ xuất hiện trong chất cách điện. d) Đúng: Khi giải cần xác định đúng dữ kiện, phương chiều, điều kiện áp dụng và đơn vị. Kiến thức cốt lõi: Dòng điện Foucault là dòng điện cảm ứng xuất hiện trong khối vật dẫn khi từ thông qua nó biến thiên
\item (Sai) Kết quả không phụ thuộc dữ kiện và có thể bỏ qua điều kiện.
\item (Sai) Dòng điện Foucault chỉ xuất hiện trong chất cách điện.
\item (Đúng) Khi giải cần xác định đúng dữ kiện, phương chiều, điều kiện áp dụng và đơn vị.
\end{itemchoice}
}
\end{ex}

% ===== Câu 99 =====
% ID: L12C3B18-20-TL
% Mức: TH
\begin{bt}
% ID: L12C3B18-20-TL
% Mức: TH
Giải thích vì sao nhận định sau đúng: “Dòng điện Foucault là dòng điện cảm ứng xuất hiện trong khối vật dẫn khi từ thông qua nó biến thiên.”
\loigiai{
Cần nêu được: Dòng điện Foucault là dòng điện cảm ứng xuất hiện trong khối vật dẫn khi từ thông qua nó biến thiên. Khi vận dụng phải xác định đúng dữ kiện, phương chiều nếu có, điều kiện áp dụng, hệ đơn vị và kiểm tra tính hợp lí. Nhận định “Dòng điện Foucault chỉ xuất hiện trong chất cách điện.” sai.
}
\end{bt}

% ===== Câu 100 =====
% ID: L12C3B18-20-TLN
% Mức: TH
\begin{ex}
% ID: L12C3B18-20-TLN
% Mức: TH
Trong bốn phát biểu về Nhận biết dòng điện Foucault và các ứng dụng của cảm ứng điện từ, có bao nhiêu phát biểu đúng? (Chỉ ghi một số nguyên.) (1) Dòng điện Foucault chỉ xuất hiện trong chất cách điện. (2) Dòng điện Foucault là dòng điện cảm ứng xuất hiện trong khối vật dẫn khi từ thông qua nó biến thiên. (3) Kết quả phải phù hợp ý nghĩa vật lí. (4) Mọi đại lượng đều là vô hướng.
\shortans{1}
\loigiai{
Đối chiếu với kiến thức: Dòng điện Foucault là dòng điện cảm ứng xuất hiện trong khối vật dẫn khi từ thông qua nó biến thiên. Có 1 phát biểu đúng.
}
\end{ex}

% ===== Câu 101 =====
% ID: L12C3B18-20-TN
% Mức: VD
\dangbt{Giải thích nguyên lí bếp từ, phanh điện từ và gia nhiệt cảm ứng}
\begin{ex}
% ID: L12C3B18-20-TN
% Mức: VD
Kết luận đúng khi vận dụng Giải thích nguyên lí bếp từ, phanh điện từ và gia nhiệt cảm ứng?
\choice
{\True Bếp từ làm nóng đáy nồi dẫn điện nhờ dòng điện Foucault và sự tỏa nhiệt Joule.}
{Bếp từ làm nóng trực tiếp mọi vật liệu bằng điện trường tĩnh.}
{Có thể bỏ qua phương, chiều và đơn vị trong mọi trường hợp.}
{Kết quả không phụ thuộc dữ kiện và điều kiện áp dụng.}
\loigiai{
Bếp từ làm nóng đáy nồi dẫn điện nhờ dòng điện Foucault và sự tỏa nhiệt Joule. Vì vậy phương án $A$ đúng; các phương án còn lại sai.
}
\end{ex}

% ===== Câu 102 =====
% ID: L12C3B18-21-DS
% Mức: TH
\dangbt{Truyền tải điện năng đi xa và giảm hao phí trên đường dây}
\begin{ex}
% ID: L12C3B18-21-DS
% Mức: TH
Xét Truyền tải điện năng đi xa và giảm hao phí trên đường dây (Tình huống 1). Hãy xác định tính đúng, sai của bốn phát biểu sau.
\choiceTF
{\True Với công suất truyền không đổi, tăng điện áp truyền tải làm giảm dòng điện và giảm hao phí I^2R.}
{Muốn giảm hao phí phải giảm điện áp truyền tải.}
{\True Khi giải cần xác định đúng dữ kiện, phương chiều, điều kiện áp dụng và đơn vị.}
{Kết quả không phụ thuộc dữ kiện và có thể bỏ qua điều kiện.}
\loigiai{
\begin{itemchoice}
\item (Đúng) Với công suất truyền không đổi, tăng điện áp truyền tải làm giảm dòng điện và giảm hao phí I^2R. (Vì): Với công suất truyền không đổi, tăng điện áp truyền tải làm giảm dòng điện và giảm hao phí I^2R. b) Sai: Muốn giảm hao phí phải giảm điện áp truyền tải. c) Đúng: Khi giải cần xác định đúng dữ kiện, phương chiều, điều kiện áp dụng và đơn vị. d) Sai: Kết quả không phụ thuộc dữ kiện và có thể bỏ qua điều kiện. Kiến thức cốt lõi: Với công suất truyền không đổi, tăng điện áp truyền tải làm giảm dòng điện và giảm hao phí I^2R
\item (Sai) Muốn giảm hao phí phải giảm điện áp truyền tải.
\item (Đúng) Khi giải cần xác định đúng dữ kiện, phương chiều, điều kiện áp dụng và đơn vị.
\item (Sai) Kết quả không phụ thuộc dữ kiện và có thể bỏ qua điều kiện.
\end{itemchoice}
}
\end{ex}

% ===== Câu 103 =====
% ID: L12C3B18-21-TL
% Mức: VD
\dangbt{Nhận biết dòng điện Foucault và các ứng dụng của cảm ứng điện từ}
\begin{bt}
% ID: L12C3B18-21-TL
% Mức: VD
Phân tích sai lầm trong nhận định: “Dòng điện Foucault chỉ xuất hiện trong chất cách điện.”
\loigiai{
Cần nêu được: Dòng điện Foucault là dòng điện cảm ứng xuất hiện trong khối vật dẫn khi từ thông qua nó biến thiên. Khi vận dụng phải xác định đúng dữ kiện, phương chiều nếu có, điều kiện áp dụng, hệ đơn vị và kiểm tra tính hợp lí. Nhận định “Dòng điện Foucault chỉ xuất hiện trong chất cách điện.” sai.
}
\end{bt}

% ===== Câu 104 =====
% ID: L12C3B18-21-TLN
% Mức: TH
\begin{ex}
% ID: L12C3B18-21-TLN
% Mức: TH
Trong bốn phát biểu về Nhận biết dòng điện Foucault và các ứng dụng của cảm ứng điện từ, có bao nhiêu phát biểu đúng? (Chỉ ghi một số nguyên.) (1) Dòng điện Foucault là dòng điện cảm ứng xuất hiện trong khối vật dẫn khi từ thông qua nó biến thiên. (2) Dòng điện Foucault chỉ xuất hiện trong chất cách điện. (3) Cần kiểm tra điều kiện áp dụng. (4) Có thể bỏ qua đơn vị.
\shortans{2}
\loigiai{
Đối chiếu với kiến thức: Dòng điện Foucault là dòng điện cảm ứng xuất hiện trong khối vật dẫn khi từ thông qua nó biến thiên. Có 2 phát biểu đúng.
}
\end{ex}

% ===== Câu 105 =====
% ID: L12C3B18-21-TN
% Mức: NB
\dangbt{Nhận biết cấu tạo, nguyên lí hoạt động của máy biến áp}
\begin{ex}
% ID: L12C3B18-21-TN
% Mức: NB
Phát biểu nào sau đây đúng về Nhận biết cấu tạo, nguyên lí hoạt động của máy biến áp?
\choice
{Máy biến áp lí tưởng hoạt động được với nguồn điện một chiều không đổi.}
{Có thể bỏ qua phương, chiều và đơn vị trong mọi trường hợp.}
{Kết quả không phụ thuộc dữ kiện và điều kiện áp dụng.}
{\True Máy biến áp gồm lõi sắt và hai cuộn dây, hoạt động dựa trên cảm ứng điện từ với dòng điện xoay chiều.}
\loigiai{
Máy biến áp gồm lõi sắt và hai cuộn dây, hoạt động dựa trên cảm ứng điện từ với dòng điện xoay chiều. Vì vậy phương án $D$ đúng; các phương án còn lại sai.
}
\end{ex}

% ===== Câu 106 =====
% ID: L12C3B18-22-DS
% Mức: TH
\dangbt{Truyền tải điện năng đi xa và giảm hao phí trên đường dây}
\begin{ex}
% ID: L12C3B18-22-DS
% Mức: TH
Xét Truyền tải điện năng đi xa và giảm hao phí trên đường dây (Tình huống 2). Hãy xác định tính đúng, sai của bốn phát biểu sau.
\choiceTF
{Muốn giảm hao phí phải giảm điện áp truyền tải.}
{\True Với công suất truyền không đổi, tăng điện áp truyền tải làm giảm dòng điện và giảm hao phí I^2R.}
{Kết quả không phụ thuộc dữ kiện và có thể bỏ qua điều kiện.}
{\True Khi giải cần xác định đúng dữ kiện, phương chiều, điều kiện áp dụng và đơn vị.}
\loigiai{
\begin{itemchoice}
\item (Sai) Muốn giảm hao phí phải giảm điện áp truyền tải. (Vì): Muốn giảm hao phí phải giảm điện áp truyền tải. b) Đúng: Với công suất truyền không đổi, tăng điện áp truyền tải làm giảm dòng điện và giảm hao phí I^2R. c) Sai: Kết quả không phụ thuộc dữ kiện và có thể bỏ qua điều kiện. d) Đúng: Khi giải cần xác định đúng dữ kiện, phương chiều, điều kiện áp dụng và đơn vị. Kiến thức cốt lõi: Với công suất truyền không đổi, tăng điện áp truyền tải làm giảm dòng điện và giảm hao phí I^2R
\item (Đúng) Với công suất truyền không đổi, tăng điện áp truyền tải làm giảm dòng điện và giảm hao phí I^2R.
\item (Sai) Kết quả không phụ thuộc dữ kiện và có thể bỏ qua điều kiện.
\item (Đúng) Khi giải cần xác định đúng dữ kiện, phương chiều, điều kiện áp dụng và đơn vị.
\end{itemchoice}
}
\end{ex}

% ===== Câu 107 =====
% ID: L12C3B18-22-TL
% Mức: VD
\dangbt{Nhận biết dòng điện Foucault và các ứng dụng của cảm ứng điện từ}
\begin{bt}
% ID: L12C3B18-22-TL
% Mức: VD
Đề xuất các bước giải một bài toán thuộc dạng Nhận biết dòng điện Foucault và các ứng dụng của cảm ứng điện từ.
\loigiai{
Cần nêu được: Dòng điện Foucault là dòng điện cảm ứng xuất hiện trong khối vật dẫn khi từ thông qua nó biến thiên. Khi vận dụng phải xác định đúng dữ kiện, phương chiều nếu có, điều kiện áp dụng, hệ đơn vị và kiểm tra tính hợp lí. Nhận định “Dòng điện Foucault chỉ xuất hiện trong chất cách điện.” sai.
}
\end{bt}

% ===== Câu 108 =====
% ID: L12C3B18-22-TLN
% Mức: VD
\begin{ex}
% ID: L12C3B18-22-TLN
% Mức: VD
Trong bốn phát biểu về Nhận biết dòng điện Foucault và các ứng dụng của cảm ứng điện từ, có bao nhiêu phát biểu đúng? (Chỉ ghi một số nguyên.) (1) Dòng điện Foucault chỉ xuất hiện trong chất cách điện. (2) Dòng điện Foucault là dòng điện cảm ứng xuất hiện trong khối vật dẫn khi từ thông qua nó biến thiên. (3) Kết quả phải phù hợp ý nghĩa vật lí. (4) Mọi đại lượng đều là vô hướng.
\shortans{2}
\loigiai{
Đối chiếu với kiến thức: Dòng điện Foucault là dòng điện cảm ứng xuất hiện trong khối vật dẫn khi từ thông qua nó biến thiên. Có 2 phát biểu đúng.
}
\end{ex}

% ===== Câu 109 =====
% ID: L12C3B18-22-TN
% Mức: NB
\dangbt{Nhận biết cấu tạo, nguyên lí hoạt động của máy biến áp}
\begin{ex}
% ID: L12C3B18-22-TN
% Mức: NB
Trong nội dung Nhận biết cấu tạo, nguyên lí hoạt động của máy biến áp?
\choice
{\True Máy biến áp gồm lõi sắt và hai cuộn dây, hoạt động dựa trên cảm ứng điện từ với dòng điện xoay chiều.}
{Máy biến áp lí tưởng hoạt động được với nguồn điện một chiều không đổi.}
{Có thể bỏ qua phương, chiều và đơn vị trong mọi trường hợp.}
{Kết quả không phụ thuộc dữ kiện và điều kiện áp dụng.}
\loigiai{
Máy biến áp gồm lõi sắt và hai cuộn dây, hoạt động dựa trên cảm ứng điện từ với dòng điện xoay chiều. Vì vậy phương án $A$ đúng; các phương án còn lại sai.
}
\end{ex}

% ===== Câu 110 =====
% ID: L12C3B18-23-DS
% Mức: VD
\dangbt{Truyền tải điện năng đi xa và giảm hao phí trên đường dây}
\begin{ex}
% ID: L12C3B18-23-DS
% Mức: VD
Xét Truyền tải điện năng đi xa và giảm hao phí trên đường dây (Tình huống 3). Hãy xác định tính đúng, sai của bốn phát biểu sau.
\choiceTF
{\True Khi giải cần xác định đúng dữ kiện, phương chiều, điều kiện áp dụng và đơn vị.}
{Kết quả không phụ thuộc dữ kiện và có thể bỏ qua điều kiện.}
{\True Với công suất truyền không đổi, tăng điện áp truyền tải làm giảm dòng điện và giảm hao phí I^2R.}
{Muốn giảm hao phí phải giảm điện áp truyền tải.}
\loigiai{
\begin{itemchoice}
\item (Đúng) Khi giải cần xác định đúng dữ kiện, phương chiều, điều kiện áp dụng và đơn vị. (Vì): Khi giải cần xác định đúng dữ kiện, phương chiều, điều kiện áp dụng và đơn vị. b) Sai: Kết quả không phụ thuộc dữ kiện và có thể bỏ qua điều kiện. c) Đúng: Với công suất truyền không đổi, tăng điện áp truyền tải làm giảm dòng điện và giảm hao phí I^2R. d) Sai: Muốn giảm hao phí phải giảm điện áp truyền tải. Kiến thức cốt lõi: Với công suất truyền không đổi, tăng điện áp truyền tải làm giảm dòng điện và giảm hao phí I^2R
\item (Sai) Kết quả không phụ thuộc dữ kiện và có thể bỏ qua điều kiện.
\item (Đúng) Với công suất truyền không đổi, tăng điện áp truyền tải làm giảm dòng điện và giảm hao phí I^2R.
\item (Sai) Muốn giảm hao phí phải giảm điện áp truyền tải.
\end{itemchoice}
}
\end{ex}

% ===== Câu 111 =====
% ID: L12C3B18-23-TL
% Mức: VD
\dangbt{Nhận biết dòng điện Foucault và các ứng dụng của cảm ứng điện từ}
\begin{bt}
% ID: L12C3B18-23-TL
% Mức: VD
Nêu một tình huống thực tiễn liên quan đến Nhận biết dòng điện Foucault và các ứng dụng của cảm ứng điện từ và giải thích bằng kiến thức Vật lí.
\loigiai{
Cần nêu được: Dòng điện Foucault là dòng điện cảm ứng xuất hiện trong khối vật dẫn khi từ thông qua nó biến thiên. Khi vận dụng phải xác định đúng dữ kiện, phương chiều nếu có, điều kiện áp dụng, hệ đơn vị và kiểm tra tính hợp lí. Nhận định “Dòng điện Foucault chỉ xuất hiện trong chất cách điện.” sai.
}
\end{bt}

% ===== Câu 112 =====
% ID: L12C3B18-23-TLN
% Mức: VD
\begin{ex}
% ID: L12C3B18-23-TLN
% Mức: VD
Trong bốn phát biểu về Nhận biết dòng điện Foucault và các ứng dụng của cảm ứng điện từ, có bao nhiêu phát biểu đúng? (Chỉ ghi một số nguyên.) (1) Dòng điện Foucault là dòng điện cảm ứng xuất hiện trong khối vật dẫn khi từ thông qua nó biến thiên. (2) Dòng điện Foucault chỉ xuất hiện trong chất cách điện. (3) Cần kiểm tra điều kiện áp dụng. (4) Có thể bỏ qua đơn vị.
\shortans{2}
\loigiai{
Đối chiếu với kiến thức: Dòng điện Foucault là dòng điện cảm ứng xuất hiện trong khối vật dẫn khi từ thông qua nó biến thiên. Có 2 phát biểu đúng.
}
\end{ex}

% ===== Câu 113 =====
% ID: L12C3B18-23-TN
% Mức: NB
\dangbt{Nhận biết cấu tạo, nguyên lí hoạt động của máy biến áp}
\begin{ex}
% ID: L12C3B18-23-TN
% Mức: NB
Tình huống 3: chọn kết luận chính xác về Nhận biết cấu tạo, nguyên lí hoạt động của máy biến áp?
\choice
{Kết quả không phụ thuộc dữ kiện và điều kiện áp dụng.}
{\True Máy biến áp gồm lõi sắt và hai cuộn dây, hoạt động dựa trên cảm ứng điện từ với dòng điện xoay chiều.}
{Máy biến áp lí tưởng hoạt động được với nguồn điện một chiều không đổi.}
{Có thể bỏ qua phương, chiều và đơn vị trong mọi trường hợp.}
\loigiai{
Máy biến áp gồm lõi sắt và hai cuộn dây, hoạt động dựa trên cảm ứng điện từ với dòng điện xoay chiều. Vì vậy phương án $B$ đúng; các phương án còn lại sai.
}
\end{ex}

% ===== Câu 114 =====
% ID: L12C3B18-24-DS
% Mức: VD
\dangbt{Truyền tải điện năng đi xa và giảm hao phí trên đường dây}
\begin{ex}
% ID: L12C3B18-24-DS
% Mức: VD
Xét Truyền tải điện năng đi xa và giảm hao phí trên đường dây (Tình huống 4). Hãy xác định tính đúng, sai của bốn phát biểu sau.
\choiceTF
{Kết quả không phụ thuộc dữ kiện và có thể bỏ qua điều kiện.}
{\True Khi giải cần xác định đúng dữ kiện, phương chiều, điều kiện áp dụng và đơn vị.}
{Muốn giảm hao phí phải giảm điện áp truyền tải.}
{\True Với công suất truyền không đổi, tăng điện áp truyền tải làm giảm dòng điện và giảm hao phí I^2R.}
\loigiai{
\begin{itemchoice}
\item (Sai) Kết quả không phụ thuộc dữ kiện và có thể bỏ qua điều kiện. (Vì): Kết quả không phụ thuộc dữ kiện và có thể bỏ qua điều kiện. b) Đúng: Khi giải cần xác định đúng dữ kiện, phương chiều, điều kiện áp dụng và đơn vị. c) Sai: Muốn giảm hao phí phải giảm điện áp truyền tải. d) Đúng: Với công suất truyền không đổi, tăng điện áp truyền tải làm giảm dòng điện và giảm hao phí I^2R. Kiến thức cốt lõi: Với công suất truyền không đổi, tăng điện áp truyền tải làm giảm dòng điện và giảm hao phí I^2R
\item (Đúng) Khi giải cần xác định đúng dữ kiện, phương chiều, điều kiện áp dụng và đơn vị.
\item (Sai) Muốn giảm hao phí phải giảm điện áp truyền tải.
\item (Đúng) Với công suất truyền không đổi, tăng điện áp truyền tải làm giảm dòng điện và giảm hao phí I^2R.
\end{itemchoice}
}
\end{ex}

% ===== Câu 115 =====
% ID: L12C3B18-24-TL
% Mức: VDC
\dangbt{Nhận biết dòng điện Foucault và các ứng dụng của cảm ứng điện từ}
\begin{bt}
% ID: L12C3B18-24-TL
% Mức: VDC
So sánh nhận định đúng “Dòng điện Foucault là dòng điện cảm ứng xuất hiện trong khối vật dẫn khi từ thông qua nó biến thiên.” với nhận định sai “Dòng điện Foucault chỉ xuất hiện trong chất cách điện.”; chỉ rõ căn cứ Vật lí.
\loigiai{
Cần nêu được: Dòng điện Foucault là dòng điện cảm ứng xuất hiện trong khối vật dẫn khi từ thông qua nó biến thiên. Khi vận dụng phải xác định đúng dữ kiện, phương chiều nếu có, điều kiện áp dụng, hệ đơn vị và kiểm tra tính hợp lí. Nhận định “Dòng điện Foucault chỉ xuất hiện trong chất cách điện.” sai.
}
\end{bt}

% ===== Câu 116 =====
% ID: L12C3B18-24-TLN
% Mức: VDC
\begin{ex}
% ID: L12C3B18-24-TLN
% Mức: VDC
Trong bốn phát biểu về Nhận biết dòng điện Foucault và các ứng dụng của cảm ứng điện từ, có bao nhiêu phát biểu đúng? (Chỉ ghi một số nguyên.) (1) Khi giải cần xác định đúng dữ kiện, phương chiều, điều kiện áp dụng và đơn vị. (2) Dòng điện Foucault chỉ xuất hiện trong chất cách điện. (3) Phải dùng đúng định luật cho tình huống. (4) Dòng điện Foucault là dòng điện cảm ứng xuất hiện trong khối vật dẫn khi từ thông qua nó biến thiên.
\shortans{3}
\loigiai{
Đối chiếu với kiến thức: Dòng điện Foucault là dòng điện cảm ứng xuất hiện trong khối vật dẫn khi từ thông qua nó biến thiên. Có 3 phát biểu đúng.
}
\end{ex}

% ===== Câu 117 =====
% ID: L12C3B18-24-TN
% Mức: TH
\dangbt{Nhận biết cấu tạo, nguyên lí hoạt động của máy biến áp}
\begin{ex}
% ID: L12C3B18-24-TN
% Mức: TH
Nội dung nào mô tả đúng nhất Nhận biết cấu tạo, nguyên lí hoạt động của máy biến áp?
\choice
{Có thể bỏ qua phương, chiều và đơn vị trong mọi trường hợp.}
{Kết quả không phụ thuộc dữ kiện và điều kiện áp dụng.}
{\True Máy biến áp gồm lõi sắt và hai cuộn dây, hoạt động dựa trên cảm ứng điện từ với dòng điện xoay chiều.}
{Máy biến áp lí tưởng hoạt động được với nguồn điện một chiều không đổi.}
\loigiai{
Máy biến áp gồm lõi sắt và hai cuộn dây, hoạt động dựa trên cảm ứng điện từ với dòng điện xoay chiều. Vì vậy phương án $C$ đúng; các phương án còn lại sai.
}
\end{ex}

% ===== Câu 118 =====
% ID: L12C3B18-25-DS
% Mức: VD
\dangbt{Truyền tải điện năng đi xa và giảm hao phí trên đường dây}
\begin{ex}
% ID: L12C3B18-25-DS
% Mức: VD
Xét Truyền tải điện năng đi xa và giảm hao phí trên đường dây (Tình huống 5). Hãy xác định tính đúng, sai của bốn phát biểu sau.
\choiceTF
{\True Với công suất truyền không đổi, tăng điện áp truyền tải làm giảm dòng điện và giảm hao phí I^2R.}
{Kết quả không phụ thuộc dữ kiện và có thể bỏ qua điều kiện.}
{Muốn giảm hao phí phải giảm điện áp truyền tải.}
{\True Khi giải cần xác định đúng dữ kiện, phương chiều, điều kiện áp dụng và đơn vị.}
\loigiai{
\begin{itemchoice}
\item (Đúng) Với công suất truyền không đổi, tăng điện áp truyền tải làm giảm dòng điện và giảm hao phí I^2R. (Vì): Với công suất truyền không đổi, tăng điện áp truyền tải làm giảm dòng điện và giảm hao phí I^2R. b) Sai: Kết quả không phụ thuộc dữ kiện và có thể bỏ qua điều kiện. c) Sai: Muốn giảm hao phí phải giảm điện áp truyền tải. d) Đúng: Khi giải cần xác định đúng dữ kiện, phương chiều, điều kiện áp dụng và đơn vị. Kiến thức cốt lõi: Với công suất truyền không đổi, tăng điện áp truyền tải làm giảm dòng điện và giảm hao phí I^2R
\item (Sai) Kết quả không phụ thuộc dữ kiện và có thể bỏ qua điều kiện.
\item (Sai) Muốn giảm hao phí phải giảm điện áp truyền tải.
\item (Đúng) Khi giải cần xác định đúng dữ kiện, phương chiều, điều kiện áp dụng và đơn vị.
\end{itemchoice}
}
\end{ex}

% ===== Câu 119 =====
% ID: L12C3B18-25-TL
% Mức: TH
\begin{bt}
% ID: L12C3B18-25-TL
% Mức: TH
Trình bày kiến thức cốt lõi của Truyền tải điện năng đi xa và giảm hao phí trên đường dây và nêu điều kiện áp dụng.
\loigiai{
Cần nêu được: Với công suất truyền không đổi, tăng điện áp truyền tải làm giảm dòng điện và giảm hao phí I^2R. Khi vận dụng phải xác định đúng dữ kiện, phương chiều nếu có, điều kiện áp dụng, hệ đơn vị và kiểm tra tính hợp lí. Nhận định “Muốn giảm hao phí phải giảm điện áp truyền tải.” sai.
}
\end{bt}

% ===== Câu 120 =====
% ID: L12C3B18-25-TLN
% Mức: TH
\begin{ex}
% ID: L12C3B18-25-TLN
% Mức: TH
Trong bốn phát biểu về Truyền tải điện năng đi xa và giảm hao phí trên đường dây, có bao nhiêu phát biểu đúng? (Chỉ ghi một số nguyên.) (1) Với công suất truyền không đổi, tăng điện áp truyền tải làm giảm dòng điện và giảm hao phí I^2R. (2) Muốn giảm hao phí phải giảm điện áp truyền tải. (3) Cần kiểm tra điều kiện áp dụng. (4) Có thể bỏ qua đơn vị.
\shortans{2}
\loigiai{
Đối chiếu với kiến thức: Với công suất truyền không đổi, tăng điện áp truyền tải làm giảm dòng điện và giảm hao phí I^2R. Có 2 phát biểu đúng.
}
\end{ex}

% ===== Câu 121 =====
% ID: L12C3B18-25-TN
% Mức: TH
\dangbt{Nhận biết cấu tạo, nguyên lí hoạt động của máy biến áp}
\begin{ex}
% ID: L12C3B18-25-TN
% Mức: TH
Khi phân tích Nhận biết cấu tạo, nguyên lí hoạt động của máy biến áp?
\choice
{Máy biến áp lí tưởng hoạt động được với nguồn điện một chiều không đổi.}
{Có thể bỏ qua phương, chiều và đơn vị trong mọi trường hợp.}
{Kết quả không phụ thuộc dữ kiện và điều kiện áp dụng.}
{\True Máy biến áp gồm lõi sắt và hai cuộn dây, hoạt động dựa trên cảm ứng điện từ với dòng điện xoay chiều.}
\loigiai{
Máy biến áp gồm lõi sắt và hai cuộn dây, hoạt động dựa trên cảm ứng điện từ với dòng điện xoay chiều. Vì vậy phương án $D$ đúng; các phương án còn lại sai.
}
\end{ex}

% ===== Câu 122 =====
% ID: L12C3B18-26-DS
% Mức: TH
\dangbt{Tính điện áp, số vòng dây và cường độ dòng điện của máy biến áp lí tưởng}
\begin{ex}
% ID: L12C3B18-26-DS
% Mức: TH
Xét Tính điện áp, số vòng dây và cường độ dòng điện của máy biến áp lí tưởng (Tình huống 1). Hãy xác định tính đúng, sai của bốn phát biểu sau.
\choiceTF
{\True Máy biến áp lí tưởng thỏa U2/U1 = N2/N1 và U1I1 = U2I2.}
{Máy biến áp lí tưởng luôn có U2 = U1 bất kể số vòng dây.}
{\True Khi giải cần xác định đúng dữ kiện, phương chiều, điều kiện áp dụng và đơn vị.}
{Kết quả không phụ thuộc dữ kiện và có thể bỏ qua điều kiện.}
\loigiai{
\begin{itemchoice}
\item (Đúng) Máy biến áp lí tưởng thỏa U2/U1 = N2/N1 và U1I1 = U2I2. (Vì): Máy biến áp lí tưởng thỏa U2/U1 = N2/N1 và U1I1 = U2I2. b) Sai: Máy biến áp lí tưởng luôn có U2 = U1 bất kể số vòng dây. c) Đúng: Khi giải cần xác định đúng dữ kiện, phương chiều, điều kiện áp dụng và đơn vị. d) Sai: Kết quả không phụ thuộc dữ kiện và có thể bỏ qua điều kiện. Kiến thức cốt lõi: Máy biến áp lí tưởng thỏa U2/U1 = N2/N1 và U1I1 = U2I2
\item (Sai) Máy biến áp lí tưởng luôn có U2 = U1 bất kể số vòng dây.
\item (Đúng) Khi giải cần xác định đúng dữ kiện, phương chiều, điều kiện áp dụng và đơn vị.
\item (Sai) Kết quả không phụ thuộc dữ kiện và có thể bỏ qua điều kiện.
\end{itemchoice}
}
\end{ex}

% ===== Câu 123 =====
% ID: L12C3B18-26-TL
% Mức: TH
\dangbt{Truyền tải điện năng đi xa và giảm hao phí trên đường dây}
\begin{bt}
% ID: L12C3B18-26-TL
% Mức: TH
Giải thích vì sao nhận định sau đúng: “Với công suất truyền không đổi, tăng điện áp truyền tải làm giảm dòng điện và giảm hao phí I^2R.”
\loigiai{
Cần nêu được: Với công suất truyền không đổi, tăng điện áp truyền tải làm giảm dòng điện và giảm hao phí I^2R. Khi vận dụng phải xác định đúng dữ kiện, phương chiều nếu có, điều kiện áp dụng, hệ đơn vị và kiểm tra tính hợp lí. Nhận định “Muốn giảm hao phí phải giảm điện áp truyền tải.” sai.
}
\end{bt}

% ===== Câu 124 =====
% ID: L12C3B18-26-TLN
% Mức: TH
\begin{ex}
% ID: L12C3B18-26-TLN
% Mức: TH
Trong bốn phát biểu về Truyền tải điện năng đi xa và giảm hao phí trên đường dây, có bao nhiêu phát biểu đúng? (Chỉ ghi một số nguyên.) (1) Muốn giảm hao phí phải giảm điện áp truyền tải. (2) Với công suất truyền không đổi, tăng điện áp truyền tải làm giảm dòng điện và giảm hao phí I^2R. (3) Kết quả phải phù hợp ý nghĩa vật lí. (4) Mọi đại lượng đều là vô hướng.
\shortans{1}
\loigiai{
Đối chiếu với kiến thức: Với công suất truyền không đổi, tăng điện áp truyền tải làm giảm dòng điện và giảm hao phí I^2R. Có 1 phát biểu đúng.
}
\end{ex}

% ===== Câu 125 =====
% ID: L12C3B18-26-TN
% Mức: TH
\dangbt{Nhận biết cấu tạo, nguyên lí hoạt động của máy biến áp}
\begin{ex}
% ID: L12C3B18-26-TN
% Mức: TH
Một học sinh ôn tập Nhận biết cấu tạo, nguyên lí hoạt động của máy biến áp?
\choice
{\True Máy biến áp gồm lõi sắt và hai cuộn dây, hoạt động dựa trên cảm ứng điện từ với dòng điện xoay chiều.}
{Máy biến áp lí tưởng hoạt động được với nguồn điện một chiều không đổi.}
{Có thể bỏ qua phương, chiều và đơn vị trong mọi trường hợp.}
{Kết quả không phụ thuộc dữ kiện và điều kiện áp dụng.}
\loigiai{
Máy biến áp gồm lõi sắt và hai cuộn dây, hoạt động dựa trên cảm ứng điện từ với dòng điện xoay chiều. Vì vậy phương án $A$ đúng; các phương án còn lại sai.
}
\end{ex}

% ===== Câu 126 =====
% ID: L12C3B18-27-DS
% Mức: TH
\dangbt{Tính điện áp, số vòng dây và cường độ dòng điện của máy biến áp lí tưởng}
\begin{ex}
% ID: L12C3B18-27-DS
% Mức: TH
Xét Tính điện áp, số vòng dây và cường độ dòng điện của máy biến áp lí tưởng (Tình huống 2). Hãy xác định tính đúng, sai của bốn phát biểu sau.
\choiceTF
{Máy biến áp lí tưởng luôn có U2 = U1 bất kể số vòng dây.}
{\True Máy biến áp lí tưởng thỏa U2/U1 = N2/N1 và U1I1 = U2I2.}
{Kết quả không phụ thuộc dữ kiện và có thể bỏ qua điều kiện.}
{\True Khi giải cần xác định đúng dữ kiện, phương chiều, điều kiện áp dụng và đơn vị.}
\loigiai{
\begin{itemchoice}
\item (Sai) Máy biến áp lí tưởng luôn có U2 = U1 bất kể số vòng dây. (Vì): Máy biến áp lí tưởng luôn có U2 = U1 bất kể số vòng dây. b) Đúng: Máy biến áp lí tưởng thỏa U2/U1 = N2/N1 và U1I1 = U2I2. c) Sai: Kết quả không phụ thuộc dữ kiện và có thể bỏ qua điều kiện. d) Đúng: Khi giải cần xác định đúng dữ kiện, phương chiều, điều kiện áp dụng và đơn vị. Kiến thức cốt lõi: Máy biến áp lí tưởng thỏa U2/U1 = N2/N1 và U1I1 = U2I2
\item (Đúng) Máy biến áp lí tưởng thỏa U2/U1 = N2/N1 và U1I1 = U2I2.
\item (Sai) Kết quả không phụ thuộc dữ kiện và có thể bỏ qua điều kiện.
\item (Đúng) Khi giải cần xác định đúng dữ kiện, phương chiều, điều kiện áp dụng và đơn vị.
\end{itemchoice}
}
\end{ex}

% ===== Câu 127 =====
% ID: L12C3B18-27-TL
% Mức: VD
\dangbt{Truyền tải điện năng đi xa và giảm hao phí trên đường dây}
\begin{bt}
% ID: L12C3B18-27-TL
% Mức: VD
Phân tích sai lầm trong nhận định: “Muốn giảm hao phí phải giảm điện áp truyền tải.”
\loigiai{
Cần nêu được: Với công suất truyền không đổi, tăng điện áp truyền tải làm giảm dòng điện và giảm hao phí I^2R. Khi vận dụng phải xác định đúng dữ kiện, phương chiều nếu có, điều kiện áp dụng, hệ đơn vị và kiểm tra tính hợp lí. Nhận định “Muốn giảm hao phí phải giảm điện áp truyền tải.” sai.
}
\end{bt}

% ===== Câu 128 =====
% ID: L12C3B18-27-TLN
% Mức: TH
\begin{ex}
% ID: L12C3B18-27-TLN
% Mức: TH
Trong bốn phát biểu về Truyền tải điện năng đi xa và giảm hao phí trên đường dây, có bao nhiêu phát biểu đúng? (Chỉ ghi một số nguyên.) (1) Với công suất truyền không đổi, tăng điện áp truyền tải làm giảm dòng điện và giảm hao phí I^2R. (2) Muốn giảm hao phí phải giảm điện áp truyền tải. (3) Cần kiểm tra điều kiện áp dụng. (4) Có thể bỏ qua đơn vị.
\shortans{2}
\loigiai{
Đối chiếu với kiến thức: Với công suất truyền không đổi, tăng điện áp truyền tải làm giảm dòng điện và giảm hao phí I^2R. Có 2 phát biểu đúng.
}
\end{ex}

% ===== Câu 129 =====
% ID: L12C3B18-27-TN
% Mức: TH
\dangbt{Nhận biết cấu tạo, nguyên lí hoạt động của máy biến áp}
\begin{ex}
% ID: L12C3B18-27-TN
% Mức: TH
Chọn phát biểu không trái với kiến thức về Nhận biết cấu tạo, nguyên lí hoạt động của máy biến áp?
\choice
{Kết quả không phụ thuộc dữ kiện và điều kiện áp dụng.}
{\True Máy biến áp gồm lõi sắt và hai cuộn dây, hoạt động dựa trên cảm ứng điện từ với dòng điện xoay chiều.}
{Máy biến áp lí tưởng hoạt động được với nguồn điện một chiều không đổi.}
{Có thể bỏ qua phương, chiều và đơn vị trong mọi trường hợp.}
\loigiai{
Máy biến áp gồm lõi sắt và hai cuộn dây, hoạt động dựa trên cảm ứng điện từ với dòng điện xoay chiều. Vì vậy phương án $B$ đúng; các phương án còn lại sai.
}
\end{ex}

% ===== Câu 130 =====
% ID: L12C3B18-28-DS
% Mức: VD
\dangbt{Tính điện áp, số vòng dây và cường độ dòng điện của máy biến áp lí tưởng}
\begin{ex}
% ID: L12C3B18-28-DS
% Mức: VD
Xét Tính điện áp, số vòng dây và cường độ dòng điện của máy biến áp lí tưởng (Tình huống 3). Hãy xác định tính đúng, sai của bốn phát biểu sau.
\choiceTF
{\True Khi giải cần xác định đúng dữ kiện, phương chiều, điều kiện áp dụng và đơn vị.}
{Kết quả không phụ thuộc dữ kiện và có thể bỏ qua điều kiện.}
{\True Máy biến áp lí tưởng thỏa U2/U1 = N2/N1 và U1I1 = U2I2.}
{Máy biến áp lí tưởng luôn có U2 = U1 bất kể số vòng dây.}
\loigiai{
\begin{itemchoice}
\item (Đúng) Khi giải cần xác định đúng dữ kiện, phương chiều, điều kiện áp dụng và đơn vị. (Vì): Khi giải cần xác định đúng dữ kiện, phương chiều, điều kiện áp dụng và đơn vị. b) Sai: Kết quả không phụ thuộc dữ kiện và có thể bỏ qua điều kiện. c) Đúng: Máy biến áp lí tưởng thỏa U2/U1 = N2/N1 và U1I1 = U2I2. d) Sai: Máy biến áp lí tưởng luôn có U2 = U1 bất kể số vòng dây. Kiến thức cốt lõi: Máy biến áp lí tưởng thỏa U2/U1 = N2/N1 và U1I1 = U2I2
\item (Sai) Kết quả không phụ thuộc dữ kiện và có thể bỏ qua điều kiện.
\item (Đúng) Máy biến áp lí tưởng thỏa U2/U1 = N2/N1 và U1I1 = U2I2.
\item (Sai) Máy biến áp lí tưởng luôn có U2 = U1 bất kể số vòng dây.
\end{itemchoice}
}
\end{ex}

% ===== Câu 131 =====
% ID: L12C3B18-28-TL
% Mức: VD
\dangbt{Truyền tải điện năng đi xa và giảm hao phí trên đường dây}
\begin{bt}
% ID: L12C3B18-28-TL
% Mức: VD
Đề xuất các bước giải một bài toán thuộc dạng Truyền tải điện năng đi xa và giảm hao phí trên đường dây.
\loigiai{
Cần nêu được: Với công suất truyền không đổi, tăng điện áp truyền tải làm giảm dòng điện và giảm hao phí I^2R. Khi vận dụng phải xác định đúng dữ kiện, phương chiều nếu có, điều kiện áp dụng, hệ đơn vị và kiểm tra tính hợp lí. Nhận định “Muốn giảm hao phí phải giảm điện áp truyền tải.” sai.
}
\end{bt}

% ===== Câu 132 =====
% ID: L12C3B18-28-TLN
% Mức: VD
\begin{ex}
% ID: L12C3B18-28-TLN
% Mức: VD
Trong bốn phát biểu về Truyền tải điện năng đi xa và giảm hao phí trên đường dây, có bao nhiêu phát biểu đúng? (Chỉ ghi một số nguyên.) (1) Muốn giảm hao phí phải giảm điện áp truyền tải. (2) Với công suất truyền không đổi, tăng điện áp truyền tải làm giảm dòng điện và giảm hao phí I^2R. (3) Kết quả phải phù hợp ý nghĩa vật lí. (4) Mọi đại lượng đều là vô hướng.
\shortans{2}
\loigiai{
Đối chiếu với kiến thức: Với công suất truyền không đổi, tăng điện áp truyền tải làm giảm dòng điện và giảm hao phí I^2R. Có 2 phát biểu đúng.
}
\end{ex}

% ===== Câu 133 =====
% ID: L12C3B18-28-TN
% Mức: VD
\dangbt{Nhận biết cấu tạo, nguyên lí hoạt động của máy biến áp}
\begin{ex}
% ID: L12C3B18-28-TN
% Mức: VD
Cơ sở Vật lí đúng dùng để giải Nhận biết cấu tạo, nguyên lí hoạt động của máy biến áp?
\choice
{Có thể bỏ qua phương, chiều và đơn vị trong mọi trường hợp.}
{Kết quả không phụ thuộc dữ kiện và điều kiện áp dụng.}
{\True Máy biến áp gồm lõi sắt và hai cuộn dây, hoạt động dựa trên cảm ứng điện từ với dòng điện xoay chiều.}
{Máy biến áp lí tưởng hoạt động được với nguồn điện một chiều không đổi.}
\loigiai{
Máy biến áp gồm lõi sắt và hai cuộn dây, hoạt động dựa trên cảm ứng điện từ với dòng điện xoay chiều. Vì vậy phương án $C$ đúng; các phương án còn lại sai.
}
\end{ex}

% ===== Câu 134 =====
% ID: L12C3B18-29-DS
% Mức: VD
\dangbt{Tính điện áp, số vòng dây và cường độ dòng điện của máy biến áp lí tưởng}
\begin{ex}
% ID: L12C3B18-29-DS
% Mức: VD
Xét Tính điện áp, số vòng dây và cường độ dòng điện của máy biến áp lí tưởng (Tình huống 4). Hãy xác định tính đúng, sai của bốn phát biểu sau.
\choiceTF
{Kết quả không phụ thuộc dữ kiện và có thể bỏ qua điều kiện.}
{\True Khi giải cần xác định đúng dữ kiện, phương chiều, điều kiện áp dụng và đơn vị.}
{Máy biến áp lí tưởng luôn có U2 = U1 bất kể số vòng dây.}
{\True Máy biến áp lí tưởng thỏa U2/U1 = N2/N1 và U1I1 = U2I2.}
\loigiai{
\begin{itemchoice}
\item (Sai) Kết quả không phụ thuộc dữ kiện và có thể bỏ qua điều kiện. (Vì): Kết quả không phụ thuộc dữ kiện và có thể bỏ qua điều kiện. b) Đúng: Khi giải cần xác định đúng dữ kiện, phương chiều, điều kiện áp dụng và đơn vị. c) Sai: Máy biến áp lí tưởng luôn có U2 = U1 bất kể số vòng dây. d) Đúng: Máy biến áp lí tưởng thỏa U2/U1 = N2/N1 và U1I1 = U2I2. Kiến thức cốt lõi: Máy biến áp lí tưởng thỏa U2/U1 = N2/N1 và U1I1 = U2I2
\item (Đúng) Khi giải cần xác định đúng dữ kiện, phương chiều, điều kiện áp dụng và đơn vị.
\item (Sai) Máy biến áp lí tưởng luôn có U2 = U1 bất kể số vòng dây.
\item (Đúng) Máy biến áp lí tưởng thỏa U2/U1 = N2/N1 và U1I1 = U2I2.
\end{itemchoice}
}
\end{ex}

% ===== Câu 135 =====
% ID: L12C3B18-29-TL
% Mức: VD
\dangbt{Truyền tải điện năng đi xa và giảm hao phí trên đường dây}
\begin{bt}
% ID: L12C3B18-29-TL
% Mức: VD
Nêu một tình huống thực tiễn liên quan đến Truyền tải điện năng đi xa và giảm hao phí trên đường dây và giải thích bằng kiến thức Vật lí.
\loigiai{
Cần nêu được: Với công suất truyền không đổi, tăng điện áp truyền tải làm giảm dòng điện và giảm hao phí I^2R. Khi vận dụng phải xác định đúng dữ kiện, phương chiều nếu có, điều kiện áp dụng, hệ đơn vị và kiểm tra tính hợp lí. Nhận định “Muốn giảm hao phí phải giảm điện áp truyền tải.” sai.
}
\end{bt}

% ===== Câu 136 =====
% ID: L12C3B18-29-TLN
% Mức: VD
\begin{ex}
% ID: L12C3B18-29-TLN
% Mức: VD
Trong bốn phát biểu về Truyền tải điện năng đi xa và giảm hao phí trên đường dây, có bao nhiêu phát biểu đúng? (Chỉ ghi một số nguyên.) (1) Với công suất truyền không đổi, tăng điện áp truyền tải làm giảm dòng điện và giảm hao phí I^2R. (2) Muốn giảm hao phí phải giảm điện áp truyền tải. (3) Cần kiểm tra điều kiện áp dụng. (4) Có thể bỏ qua đơn vị.
\shortans{2}
\loigiai{
Đối chiếu với kiến thức: Với công suất truyền không đổi, tăng điện áp truyền tải làm giảm dòng điện và giảm hao phí I^2R. Có 2 phát biểu đúng.
}
\end{ex}

% ===== Câu 137 =====
% ID: L12C3B18-29-TN
% Mức: VD
\dangbt{Nhận biết cấu tạo, nguyên lí hoạt động của máy biến áp}
\begin{ex}
% ID: L12C3B18-29-TN
% Mức: VD
Trong một bài thực hành về Nhận biết cấu tạo, nguyên lí hoạt động của máy biến áp?
\choice
{Máy biến áp lí tưởng hoạt động được với nguồn điện một chiều không đổi.}
{Có thể bỏ qua phương, chiều và đơn vị trong mọi trường hợp.}
{Kết quả không phụ thuộc dữ kiện và điều kiện áp dụng.}
{\True Máy biến áp gồm lõi sắt và hai cuộn dây, hoạt động dựa trên cảm ứng điện từ với dòng điện xoay chiều.}
\loigiai{
Máy biến áp gồm lõi sắt và hai cuộn dây, hoạt động dựa trên cảm ứng điện từ với dòng điện xoay chiều. Vì vậy phương án $D$ đúng; các phương án còn lại sai.
}
\end{ex}

% ===== Câu 138 =====
% ID: L12C3B18-30-DS
% Mức: VD
\dangbt{Tính điện áp, số vòng dây và cường độ dòng điện của máy biến áp lí tưởng}
\begin{ex}
% ID: L12C3B18-30-DS
% Mức: VD
Xét Tính điện áp, số vòng dây và cường độ dòng điện của máy biến áp lí tưởng (Tình huống 5). Hãy xác định tính đúng, sai của bốn phát biểu sau.
\choiceTF
{\True Máy biến áp lí tưởng thỏa U2/U1 = N2/N1 và U1I1 = U2I2.}
{Kết quả không phụ thuộc dữ kiện và có thể bỏ qua điều kiện.}
{Máy biến áp lí tưởng luôn có U2 = U1 bất kể số vòng dây.}
{\True Khi giải cần xác định đúng dữ kiện, phương chiều, điều kiện áp dụng và đơn vị.}
\loigiai{
\begin{itemchoice}
\item (Đúng) Máy biến áp lí tưởng thỏa U2/U1 = N2/N1 và U1I1 = U2I2. (Vì): Máy biến áp lí tưởng thỏa U2/U1 = N2/N1 và U1I1 = U2I2. b) Sai: Kết quả không phụ thuộc dữ kiện và có thể bỏ qua điều kiện. c) Sai: Máy biến áp lí tưởng luôn có U2 = U1 bất kể số vòng dây. d) Đúng: Khi giải cần xác định đúng dữ kiện, phương chiều, điều kiện áp dụng và đơn vị. Kiến thức cốt lõi: Máy biến áp lí tưởng thỏa U2/U1 = N2/N1 và U1I1 = U2I2
\item (Sai) Kết quả không phụ thuộc dữ kiện và có thể bỏ qua điều kiện.
\item (Sai) Máy biến áp lí tưởng luôn có U2 = U1 bất kể số vòng dây.
\item (Đúng) Khi giải cần xác định đúng dữ kiện, phương chiều, điều kiện áp dụng và đơn vị.
\end{itemchoice}
}
\end{ex}

% ===== Câu 139 =====
% ID: L12C3B18-30-TL
% Mức: VDC
\dangbt{Truyền tải điện năng đi xa và giảm hao phí trên đường dây}
\begin{bt}
% ID: L12C3B18-30-TL
% Mức: VDC
So sánh nhận định đúng “Với công suất truyền không đổi, tăng điện áp truyền tải làm giảm dòng điện và giảm hao phí I^2R.” với nhận định sai “Muốn giảm hao phí phải giảm điện áp truyền tải.”; chỉ rõ căn cứ Vật lí.
\loigiai{
Cần nêu được: Với công suất truyền không đổi, tăng điện áp truyền tải làm giảm dòng điện và giảm hao phí I^2R. Khi vận dụng phải xác định đúng dữ kiện, phương chiều nếu có, điều kiện áp dụng, hệ đơn vị và kiểm tra tính hợp lí. Nhận định “Muốn giảm hao phí phải giảm điện áp truyền tải.” sai.
}
\end{bt}

% ===== Câu 140 =====
% ID: L12C3B18-30-TLN
% Mức: VDC
\begin{ex}
% ID: L12C3B18-30-TLN
% Mức: VDC
Trong bốn phát biểu về Truyền tải điện năng đi xa và giảm hao phí trên đường dây, có bao nhiêu phát biểu đúng? (Chỉ ghi một số nguyên.) (1) Khi giải cần xác định đúng dữ kiện, phương chiều, điều kiện áp dụng và đơn vị. (2) Muốn giảm hao phí phải giảm điện áp truyền tải. (3) Phải dùng đúng định luật cho tình huống. (4) Với công suất truyền không đổi, tăng điện áp truyền tải làm giảm dòng điện và giảm hao phí I^2R.
\shortans{3}
\loigiai{
Đối chiếu với kiến thức: Với công suất truyền không đổi, tăng điện áp truyền tải làm giảm dòng điện và giảm hao phí I^2R. Có 3 phát biểu đúng.
}
\end{ex}

% ===== Câu 141 =====
% ID: L12C3B18-30-TN
% Mức: VD
\dangbt{Nhận biết cấu tạo, nguyên lí hoạt động của máy biến áp}
\begin{ex}
% ID: L12C3B18-30-TN
% Mức: VD
Kết luận đúng khi vận dụng Nhận biết cấu tạo, nguyên lí hoạt động của máy biến áp?
\choice
{\True Máy biến áp gồm lõi sắt và hai cuộn dây, hoạt động dựa trên cảm ứng điện từ với dòng điện xoay chiều.}
{Máy biến áp lí tưởng hoạt động được với nguồn điện một chiều không đổi.}
{Có thể bỏ qua phương, chiều và đơn vị trong mọi trường hợp.}
{Kết quả không phụ thuộc dữ kiện và điều kiện áp dụng.}
\loigiai{
Máy biến áp gồm lõi sắt và hai cuộn dây, hoạt động dựa trên cảm ứng điện từ với dòng điện xoay chiều. Vì vậy phương án $A$ đúng; các phương án còn lại sai.
}
\end{ex}

% ===== Câu 142 =====
% ID: L12C3B18-31-DS
% Mức: TH
\dangbt{Nhận biết dòng điện Foucault và các ứng dụng của cảm ứng điện từ}
\begin{ex}
% ID: L12C3B18-31-DS
% Mức: TH
Xét Nhận biết dòng điện Foucault và các ứng dụng của cảm ứng điện từ (Tình huống 1). Hãy xác định tính đúng, sai của bốn phát biểu sau.
\choiceTF
{\True Dòng điện Foucault là dòng điện cảm ứng xuất hiện trong khối vật dẫn khi từ thông qua nó biến thiên.}
{Dòng điện Foucault chỉ xuất hiện trong chất cách điện.}
{\True Khi giải cần xác định đúng dữ kiện, phương chiều, điều kiện áp dụng và đơn vị.}
{Kết quả không phụ thuộc dữ kiện và có thể bỏ qua điều kiện.}
\loigiai{
\begin{itemchoice}
\item (Đúng) Dòng điện Foucault là dòng điện cảm ứng xuất hiện trong khối vật dẫn khi từ thông qua nó biến thiên. (Vì): òng điện Foucault là dòng điện cảm ứng xuất hiện trong khối vật dẫn khi từ thông qua nó biến thiên. b) Sai: Dòng điện Foucault chỉ xuất hiện trong chất cách điện. c) Đúng: Khi giải cần xác định đúng dữ kiện, phương chiều, điều kiện áp dụng và đơn vị. d) Sai: Kết quả không phụ thuộc dữ kiện và có thể bỏ qua điều kiện. Kiến thức cốt lõi: Dòng điện Foucault là dòng điện cảm ứng xuất hiện trong khối vật dẫn khi từ thông qua nó biến thiên
\item (Sai) Dòng điện Foucault chỉ xuất hiện trong chất cách điện.
\item (Đúng) Khi giải cần xác định đúng dữ kiện, phương chiều, điều kiện áp dụng và đơn vị.
\item (Sai) Kết quả không phụ thuộc dữ kiện và có thể bỏ qua điều kiện.
\end{itemchoice}
}
\end{ex}

% ===== Câu 143 =====
% ID: L12C3B18-31-TL
% Mức: TH
\dangbt{Tính điện áp, số vòng dây và cường độ dòng điện của máy biến áp lí tưởng}
\begin{bt}
% ID: L12C3B18-31-TL
% Mức: TH
Trình bày kiến thức cốt lõi của Tính điện áp, số vòng dây và cường độ dòng điện của máy biến áp lí tưởng và nêu điều kiện áp dụng.
\loigiai{
Cần nêu được: Máy biến áp lí tưởng thỏa U2/U1 = N2/N1 và U1I1 = U2I2. Khi vận dụng phải xác định đúng dữ kiện, phương chiều nếu có, điều kiện áp dụng, hệ đơn vị và kiểm tra tính hợp lí. Nhận định “Máy biến áp lí tưởng luôn có U2 = U1 bất kể số vòng dây.” sai.
}
\end{bt}

% ===== Câu 144 =====
% ID: L12C3B18-31-TLN
% Mức: TH
\begin{ex}
% ID: L12C3B18-31-TLN
% Mức: TH
Trong bốn phát biểu về Tính điện áp, số vòng dây và cường độ dòng điện của máy biến áp lí tưởng, có bao nhiêu phát biểu đúng? (Chỉ ghi một số nguyên.) (1) Máy biến áp lí tưởng thỏa U2/U1 = N2/N1 và U1I1 = U2I2. (2) Máy biến áp lí tưởng luôn có U2 = U1 bất kể số vòng dây. (3) Cần kiểm tra điều kiện áp dụng. (4) Có thể bỏ qua đơn vị.
\shortans{2}
\loigiai{
Đối chiếu với kiến thức: Máy biến áp lí tưởng thỏa U2/U1 = N2/N1 và U1I1 = U2I2. Có 2 phát biểu đúng.
}
\end{ex}

% ===== Câu 145 =====
% ID: L12C3B18-31-TN
% Mức: NB
\dangbt{Nhận biết dòng điện Foucault và các ứng dụng của cảm ứng điện từ}
\begin{ex}
% ID: L12C3B18-31-TN
% Mức: NB
Phát biểu nào sau đây đúng về Nhận biết dòng điện Foucault và các ứng dụng của cảm ứng điện từ?
\choice
{Dòng điện Foucault chỉ xuất hiện trong chất cách điện.}
{Có thể bỏ qua phương, chiều và đơn vị trong mọi trường hợp.}
{Kết quả không phụ thuộc dữ kiện và điều kiện áp dụng.}
{\True Dòng điện Foucault là dòng điện cảm ứng xuất hiện trong khối vật dẫn khi từ thông qua nó biến thiên.}
\loigiai{
Dòng điện Foucault là dòng điện cảm ứng xuất hiện trong khối vật dẫn khi từ thông qua nó biến thiên. Vì vậy phương án $D$ đúng; các phương án còn lại sai.
}
\end{ex}

% ===== Câu 146 =====
% ID: L12C3B18-32-DS
% Mức: TH
\begin{ex}
% ID: L12C3B18-32-DS
% Mức: TH
Xét Nhận biết dòng điện Foucault và các ứng dụng của cảm ứng điện từ (Tình huống 2). Hãy xác định tính đúng, sai của bốn phát biểu sau.
\choiceTF
{Dòng điện Foucault chỉ xuất hiện trong chất cách điện.}
{\True Dòng điện Foucault là dòng điện cảm ứng xuất hiện trong khối vật dẫn khi từ thông qua nó biến thiên.}
{Kết quả không phụ thuộc dữ kiện và có thể bỏ qua điều kiện.}
{\True Khi giải cần xác định đúng dữ kiện, phương chiều, điều kiện áp dụng và đơn vị.}
\loigiai{
\begin{itemchoice}
\item (Sai) Dòng điện Foucault chỉ xuất hiện trong chất cách điện. (Vì): Dòng điện Foucault chỉ xuất hiện trong chất cách điện. b) Đúng: Dòng điện Foucault là dòng điện cảm ứng xuất hiện trong khối vật dẫn khi từ thông qua nó biến thiên. c) Sai: Kết quả không phụ thuộc dữ kiện và có thể bỏ qua điều kiện. d) Đúng: Khi giải cần xác định đúng dữ kiện, phương chiều, điều kiện áp dụng và đơn vị. Kiến thức cốt lõi: Dòng điện Foucault là dòng điện cảm ứng xuất hiện trong khối vật dẫn khi từ thông qua nó biến thiên
\item (Đúng) Dòng điện Foucault là dòng điện cảm ứng xuất hiện trong khối vật dẫn khi từ thông qua nó biến thiên.
\item (Sai) Kết quả không phụ thuộc dữ kiện và có thể bỏ qua điều kiện.
\item (Đúng) Khi giải cần xác định đúng dữ kiện, phương chiều, điều kiện áp dụng và đơn vị.
\end{itemchoice}
}
\end{ex}

% ===== Câu 147 =====
% ID: L12C3B18-32-TL
% Mức: TH
\dangbt{Tính điện áp, số vòng dây và cường độ dòng điện của máy biến áp lí tưởng}
\begin{bt}
% ID: L12C3B18-32-TL
% Mức: TH
Giải thích vì sao nhận định sau đúng: “Máy biến áp lí tưởng thỏa U2/U1 = N2/N1 và U1I1 = U2I2.”
\loigiai{
Cần nêu được: Máy biến áp lí tưởng thỏa U2/U1 = N2/N1 và U1I1 = U2I2. Khi vận dụng phải xác định đúng dữ kiện, phương chiều nếu có, điều kiện áp dụng, hệ đơn vị và kiểm tra tính hợp lí. Nhận định “Máy biến áp lí tưởng luôn có U2 = U1 bất kể số vòng dây.” sai.
}
\end{bt}

% ===== Câu 148 =====
% ID: L12C3B18-32-TLN
% Mức: TH
\begin{ex}
% ID: L12C3B18-32-TLN
% Mức: TH
Trong bốn phát biểu về Tính điện áp, số vòng dây và cường độ dòng điện của máy biến áp lí tưởng, có bao nhiêu phát biểu đúng? (Chỉ ghi một số nguyên.) (1) Máy biến áp lí tưởng luôn có U2 = U1 bất kể số vòng dây. (2) Máy biến áp lí tưởng thỏa U2/U1 = N2/N1 và U1I1 = U2I2. (3) Kết quả phải phù hợp ý nghĩa vật lí. (4) Mọi đại lượng đều là vô hướng.
\shortans{1}
\loigiai{
Đối chiếu với kiến thức: Máy biến áp lí tưởng thỏa U2/U1 = N2/N1 và U1I1 = U2I2. Có 1 phát biểu đúng.
}
\end{ex}

% ===== Câu 149 =====
% ID: L12C3B18-32-TN
% Mức: NB
\dangbt{Nhận biết dòng điện Foucault và các ứng dụng của cảm ứng điện từ}
\begin{ex}
% ID: L12C3B18-32-TN
% Mức: NB
Trong nội dung Nhận biết dòng điện Foucault và các ứng dụng của cảm ứng điện từ?
\choice
{\True Dòng điện Foucault là dòng điện cảm ứng xuất hiện trong khối vật dẫn khi từ thông qua nó biến thiên.}
{Dòng điện Foucault chỉ xuất hiện trong chất cách điện.}
{Có thể bỏ qua phương, chiều và đơn vị trong mọi trường hợp.}
{Kết quả không phụ thuộc dữ kiện và điều kiện áp dụng.}
\loigiai{
Dòng điện Foucault là dòng điện cảm ứng xuất hiện trong khối vật dẫn khi từ thông qua nó biến thiên. Vì vậy phương án $A$ đúng; các phương án còn lại sai.
}
\end{ex}

% ===== Câu 150 =====
% ID: L12C3B18-33-DS
% Mức: VD
\begin{ex}
% ID: L12C3B18-33-DS
% Mức: VD
Xét Nhận biết dòng điện Foucault và các ứng dụng của cảm ứng điện từ (Tình huống 3). Hãy xác định tính đúng, sai của bốn phát biểu sau.
\choiceTF
{\True Khi giải cần xác định đúng dữ kiện, phương chiều, điều kiện áp dụng và đơn vị.}
{Kết quả không phụ thuộc dữ kiện và có thể bỏ qua điều kiện.}
{\True Dòng điện Foucault là dòng điện cảm ứng xuất hiện trong khối vật dẫn khi từ thông qua nó biến thiên.}
{Dòng điện Foucault chỉ xuất hiện trong chất cách điện.}
\loigiai{
\begin{itemchoice}
\item (Đúng) Khi giải cần xác định đúng dữ kiện, phương chiều, điều kiện áp dụng và đơn vị. (Vì): Khi giải cần xác định đúng dữ kiện, phương chiều, điều kiện áp dụng và đơn vị. b) Sai: Kết quả không phụ thuộc dữ kiện và có thể bỏ qua điều kiện. c) Đúng: Dòng điện Foucault là dòng điện cảm ứng xuất hiện trong khối vật dẫn khi từ thông qua nó biến thiên. d) Sai: Dòng điện Foucault chỉ xuất hiện trong chất cách điện. Kiến thức cốt lõi: Dòng điện Foucault là dòng điện cảm ứng xuất hiện trong khối vật dẫn khi từ thông qua nó biến thiên
\item (Sai) Kết quả không phụ thuộc dữ kiện và có thể bỏ qua điều kiện.
\item (Đúng) Dòng điện Foucault là dòng điện cảm ứng xuất hiện trong khối vật dẫn khi từ thông qua nó biến thiên.
\item (Sai) Dòng điện Foucault chỉ xuất hiện trong chất cách điện.
\end{itemchoice}
}
\end{ex}

% ===== Câu 151 =====
% ID: L12C3B18-33-TL
% Mức: VD
\dangbt{Tính điện áp, số vòng dây và cường độ dòng điện của máy biến áp lí tưởng}
\begin{bt}
% ID: L12C3B18-33-TL
% Mức: VD
Phân tích sai lầm trong nhận định: “Máy biến áp lí tưởng luôn có U2 = U1 bất kể số vòng dây.”
\loigiai{
Cần nêu được: Máy biến áp lí tưởng thỏa U2/U1 = N2/N1 và U1I1 = U2I2. Khi vận dụng phải xác định đúng dữ kiện, phương chiều nếu có, điều kiện áp dụng, hệ đơn vị và kiểm tra tính hợp lí. Nhận định “Máy biến áp lí tưởng luôn có U2 = U1 bất kể số vòng dây.” sai.
}
\end{bt}

% ===== Câu 152 =====
% ID: L12C3B18-33-TLN
% Mức: TH
\begin{ex}
% ID: L12C3B18-33-TLN
% Mức: TH
Trong bốn phát biểu về Tính điện áp, số vòng dây và cường độ dòng điện của máy biến áp lí tưởng, có bao nhiêu phát biểu đúng? (Chỉ ghi một số nguyên.) (1) Máy biến áp lí tưởng thỏa U2/U1 = N2/N1 và U1I1 = U2I2. (2) Máy biến áp lí tưởng luôn có U2 = U1 bất kể số vòng dây. (3) Cần kiểm tra điều kiện áp dụng. (4) Có thể bỏ qua đơn vị.
\shortans{2}
\loigiai{
Đối chiếu với kiến thức: Máy biến áp lí tưởng thỏa U2/U1 = N2/N1 và U1I1 = U2I2. Có 2 phát biểu đúng.
}
\end{ex}

% ===== Câu 153 =====
% ID: L12C3B18-33-TN
% Mức: NB
\dangbt{Nhận biết dòng điện Foucault và các ứng dụng của cảm ứng điện từ}
\begin{ex}
% ID: L12C3B18-33-TN
% Mức: NB
Tình huống 3: chọn kết luận chính xác về Nhận biết dòng điện Foucault và các ứng dụng của cảm ứng điện từ?
\choice
{Kết quả không phụ thuộc dữ kiện và điều kiện áp dụng.}
{\True Dòng điện Foucault là dòng điện cảm ứng xuất hiện trong khối vật dẫn khi từ thông qua nó biến thiên.}
{Dòng điện Foucault chỉ xuất hiện trong chất cách điện.}
{Có thể bỏ qua phương, chiều và đơn vị trong mọi trường hợp.}
\loigiai{
Dòng điện Foucault là dòng điện cảm ứng xuất hiện trong khối vật dẫn khi từ thông qua nó biến thiên. Vì vậy phương án $B$ đúng; các phương án còn lại sai.
}
\end{ex}

% ===== Câu 154 =====
% ID: L12C3B18-34-DS
% Mức: VD
\begin{ex}
% ID: L12C3B18-34-DS
% Mức: VD
Xét Nhận biết dòng điện Foucault và các ứng dụng của cảm ứng điện từ (Tình huống 4). Hãy xác định tính đúng, sai của bốn phát biểu sau.
\choiceTF
{Kết quả không phụ thuộc dữ kiện và có thể bỏ qua điều kiện.}
{\True Khi giải cần xác định đúng dữ kiện, phương chiều, điều kiện áp dụng và đơn vị.}
{Dòng điện Foucault chỉ xuất hiện trong chất cách điện.}
{\True Dòng điện Foucault là dòng điện cảm ứng xuất hiện trong khối vật dẫn khi từ thông qua nó biến thiên.}
\loigiai{
\begin{itemchoice}
\item (Sai) Kết quả không phụ thuộc dữ kiện và có thể bỏ qua điều kiện. (Vì): Kết quả không phụ thuộc dữ kiện và có thể bỏ qua điều kiện. b) Đúng: Khi giải cần xác định đúng dữ kiện, phương chiều, điều kiện áp dụng và đơn vị. c) Sai: Dòng điện Foucault chỉ xuất hiện trong chất cách điện. d) Đúng: Dòng điện Foucault là dòng điện cảm ứng xuất hiện trong khối vật dẫn khi từ thông qua nó biến thiên. Kiến thức cốt lõi: Dòng điện Foucault là dòng điện cảm ứng xuất hiện trong khối vật dẫn khi từ thông qua nó biến thiên
\item (Đúng) Khi giải cần xác định đúng dữ kiện, phương chiều, điều kiện áp dụng và đơn vị.
\item (Sai) Dòng điện Foucault chỉ xuất hiện trong chất cách điện.
\item (Đúng) Dòng điện Foucault là dòng điện cảm ứng xuất hiện trong khối vật dẫn khi từ thông qua nó biến thiên.
\end{itemchoice}
}
\end{ex}

% ===== Câu 155 =====
% ID: L12C3B18-34-TL
% Mức: VD
\dangbt{Tính điện áp, số vòng dây và cường độ dòng điện của máy biến áp lí tưởng}
\begin{bt}
% ID: L12C3B18-34-TL
% Mức: VD
Đề xuất các bước giải một bài toán thuộc dạng Tính điện áp, số vòng dây và cường độ dòng điện của máy biến áp lí tưởng.
\loigiai{
Cần nêu được: Máy biến áp lí tưởng thỏa U2/U1 = N2/N1 và U1I1 = U2I2. Khi vận dụng phải xác định đúng dữ kiện, phương chiều nếu có, điều kiện áp dụng, hệ đơn vị và kiểm tra tính hợp lí. Nhận định “Máy biến áp lí tưởng luôn có U2 = U1 bất kể số vòng dây.” sai.
}
\end{bt}

% ===== Câu 156 =====
% ID: L12C3B18-34-TLN
% Mức: VD
\begin{ex}
% ID: L12C3B18-34-TLN
% Mức: VD
Trong bốn phát biểu về Tính điện áp, số vòng dây và cường độ dòng điện của máy biến áp lí tưởng, có bao nhiêu phát biểu đúng? (Chỉ ghi một số nguyên.) (1) Máy biến áp lí tưởng luôn có U2 = U1 bất kể số vòng dây. (2) Máy biến áp lí tưởng thỏa U2/U1 = N2/N1 và U1I1 = U2I2. (3) Kết quả phải phù hợp ý nghĩa vật lí. (4) Mọi đại lượng đều là vô hướng.
\shortans{2}
\loigiai{
Đối chiếu với kiến thức: Máy biến áp lí tưởng thỏa U2/U1 = N2/N1 và U1I1 = U2I2. Có 2 phát biểu đúng.
}
\end{ex}

% ===== Câu 157 =====
% ID: L12C3B18-34-TN
% Mức: TH
\dangbt{Nhận biết dòng điện Foucault và các ứng dụng của cảm ứng điện từ}
\begin{ex}
% ID: L12C3B18-34-TN
% Mức: TH
Nội dung nào mô tả đúng nhất Nhận biết dòng điện Foucault và các ứng dụng của cảm ứng điện từ?
\choice
{Có thể bỏ qua phương, chiều và đơn vị trong mọi trường hợp.}
{Kết quả không phụ thuộc dữ kiện và điều kiện áp dụng.}
{\True Dòng điện Foucault là dòng điện cảm ứng xuất hiện trong khối vật dẫn khi từ thông qua nó biến thiên.}
{Dòng điện Foucault chỉ xuất hiện trong chất cách điện.}
\loigiai{
Dòng điện Foucault là dòng điện cảm ứng xuất hiện trong khối vật dẫn khi từ thông qua nó biến thiên. Vì vậy phương án $C$ đúng; các phương án còn lại sai.
}
\end{ex}

% ===== Câu 158 =====
% ID: L12C3B18-35-DS
% Mức: VD
\begin{ex}
% ID: L12C3B18-35-DS
% Mức: VD
Xét Nhận biết dòng điện Foucault và các ứng dụng của cảm ứng điện từ (Tình huống 5). Hãy xác định tính đúng, sai của bốn phát biểu sau.
\choiceTF
{\True Dòng điện Foucault là dòng điện cảm ứng xuất hiện trong khối vật dẫn khi từ thông qua nó biến thiên.}
{Kết quả không phụ thuộc dữ kiện và có thể bỏ qua điều kiện.}
{Dòng điện Foucault chỉ xuất hiện trong chất cách điện.}
{\True Khi giải cần xác định đúng dữ kiện, phương chiều, điều kiện áp dụng và đơn vị.}
\loigiai{
\begin{itemchoice}
\item (Đúng) Dòng điện Foucault là dòng điện cảm ứng xuất hiện trong khối vật dẫn khi từ thông qua nó biến thiên. (Vì): òng điện Foucault là dòng điện cảm ứng xuất hiện trong khối vật dẫn khi từ thông qua nó biến thiên. b) Sai: Kết quả không phụ thuộc dữ kiện và có thể bỏ qua điều kiện. c) Sai: Dòng điện Foucault chỉ xuất hiện trong chất cách điện. d) Đúng: Khi giải cần xác định đúng dữ kiện, phương chiều, điều kiện áp dụng và đơn vị. Kiến thức cốt lõi: Dòng điện Foucault là dòng điện cảm ứng xuất hiện trong khối vật dẫn khi từ thông qua nó biến thiên
\item (Sai) Kết quả không phụ thuộc dữ kiện và có thể bỏ qua điều kiện.
\item (Sai) Dòng điện Foucault chỉ xuất hiện trong chất cách điện.
\item (Đúng) Khi giải cần xác định đúng dữ kiện, phương chiều, điều kiện áp dụng và đơn vị.
\end{itemchoice}
}
\end{ex}

% ===== Câu 159 =====
% ID: L12C3B18-35-TL
% Mức: VD
\dangbt{Tính điện áp, số vòng dây và cường độ dòng điện của máy biến áp lí tưởng}
\begin{bt}
% ID: L12C3B18-35-TL
% Mức: VD
Nêu một tình huống thực tiễn liên quan đến Tính điện áp, số vòng dây và cường độ dòng điện của máy biến áp lí tưởng và giải thích bằng kiến thức Vật lí.
\loigiai{
Cần nêu được: Máy biến áp lí tưởng thỏa U2/U1 = N2/N1 và U1I1 = U2I2. Khi vận dụng phải xác định đúng dữ kiện, phương chiều nếu có, điều kiện áp dụng, hệ đơn vị và kiểm tra tính hợp lí. Nhận định “Máy biến áp lí tưởng luôn có U2 = U1 bất kể số vòng dây.” sai.
}
\end{bt}

% ===== Câu 160 =====
% ID: L12C3B18-35-TLN
% Mức: VD
\begin{ex}
% ID: L12C3B18-35-TLN
% Mức: VD
Trong bốn phát biểu về Tính điện áp, số vòng dây và cường độ dòng điện của máy biến áp lí tưởng, có bao nhiêu phát biểu đúng? (Chỉ ghi một số nguyên.) (1) Máy biến áp lí tưởng thỏa U2/U1 = N2/N1 và U1I1 = U2I2. (2) Máy biến áp lí tưởng luôn có U2 = U1 bất kể số vòng dây. (3) Cần kiểm tra điều kiện áp dụng. (4) Có thể bỏ qua đơn vị.
\shortans{2}
\loigiai{
Đối chiếu với kiến thức: Máy biến áp lí tưởng thỏa U2/U1 = N2/N1 và U1I1 = U2I2. Có 2 phát biểu đúng.
}
\end{ex}

% ===== Câu 161 =====
% ID: L12C3B18-35-TN
% Mức: TH
\dangbt{Nhận biết dòng điện Foucault và các ứng dụng của cảm ứng điện từ}
\begin{ex}
% ID: L12C3B18-35-TN
% Mức: TH
Khi phân tích Nhận biết dòng điện Foucault và các ứng dụng của cảm ứng điện từ?
\choice
{Dòng điện Foucault chỉ xuất hiện trong chất cách điện.}
{Có thể bỏ qua phương, chiều và đơn vị trong mọi trường hợp.}
{Kết quả không phụ thuộc dữ kiện và điều kiện áp dụng.}
{\True Dòng điện Foucault là dòng điện cảm ứng xuất hiện trong khối vật dẫn khi từ thông qua nó biến thiên.}
\loigiai{
Dòng điện Foucault là dòng điện cảm ứng xuất hiện trong khối vật dẫn khi từ thông qua nó biến thiên. Vì vậy phương án $D$ đúng; các phương án còn lại sai.
}
\end{ex}

% ===== Câu 162 =====
% ID: L12C3B18-36-DS
% Mức: TH
\dangbt{Giải thích nguyên lí bếp từ, phanh điện từ và gia nhiệt cảm ứng}
\begin{ex}
% ID: L12C3B18-36-DS
% Mức: TH
Xét Giải thích nguyên lí bếp từ, phanh điện từ và gia nhiệt cảm ứng (Tình huống 1). Hãy xác định tính đúng, sai của bốn phát biểu sau.
\choiceTF
{\True Bếp từ làm nóng đáy nồi dẫn điện nhờ dòng điện Foucault và sự tỏa nhiệt Joule.}
{Bếp từ làm nóng trực tiếp mọi vật liệu bằng điện trường tĩnh.}
{\True Khi giải cần xác định đúng dữ kiện, phương chiều, điều kiện áp dụng và đơn vị.}
{Kết quả không phụ thuộc dữ kiện và có thể bỏ qua điều kiện.}
\loigiai{
\begin{itemchoice}
\item (Đúng) Bếp từ làm nóng đáy nồi dẫn điện nhờ dòng điện Foucault và sự tỏa nhiệt Joule. (Vì): Bếp từ làm nóng đáy nồi dẫn điện nhờ dòng điện Foucault và sự tỏa nhiệt Joule. b) Sai: Bếp từ làm nóng trực tiếp mọi vật liệu bằng điện trường tĩnh. c) Đúng: Khi giải cần xác định đúng dữ kiện, phương chiều, điều kiện áp dụng và đơn vị. d) Sai: Kết quả không phụ thuộc dữ kiện và có thể bỏ qua điều kiện. Kiến thức cốt lõi: Bếp từ làm nóng đáy nồi dẫn điện nhờ dòng điện Foucault và sự tỏa nhiệt Joule
\item (Sai) Bếp từ làm nóng trực tiếp mọi vật liệu bằng điện trường tĩnh.
\item (Đúng) Khi giải cần xác định đúng dữ kiện, phương chiều, điều kiện áp dụng và đơn vị.
\item (Sai) Kết quả không phụ thuộc dữ kiện và có thể bỏ qua điều kiện.
\end{itemchoice}
}
\end{ex}

% ===== Câu 163 =====
% ID: L12C3B18-36-TL
% Mức: VDC
\dangbt{Tính điện áp, số vòng dây và cường độ dòng điện của máy biến áp lí tưởng}
\begin{bt}
% ID: L12C3B18-36-TL
% Mức: VDC
So sánh nhận định đúng “Máy biến áp lí tưởng thỏa U2/U1 = N2/N1 và U1I1 = U2I2.” với nhận định sai “Máy biến áp lí tưởng luôn có U2 = U1 bất kể số vòng dây.”; chỉ rõ căn cứ Vật lí.
\loigiai{
Cần nêu được: Máy biến áp lí tưởng thỏa U2/U1 = N2/N1 và U1I1 = U2I2. Khi vận dụng phải xác định đúng dữ kiện, phương chiều nếu có, điều kiện áp dụng, hệ đơn vị và kiểm tra tính hợp lí. Nhận định “Máy biến áp lí tưởng luôn có U2 = U1 bất kể số vòng dây.” sai.
}
\end{bt}

% ===== Câu 164 =====
% ID: L12C3B18-36-TLN
% Mức: VDC
\begin{ex}
% ID: L12C3B18-36-TLN
% Mức: VDC
Trong bốn phát biểu về Tính điện áp, số vòng dây và cường độ dòng điện của máy biến áp lí tưởng, có bao nhiêu phát biểu đúng? (Chỉ ghi một số nguyên.) (1) Khi giải cần xác định đúng dữ kiện, phương chiều, điều kiện áp dụng và đơn vị. (2) Máy biến áp lí tưởng luôn có U2 = U1 bất kể số vòng dây. (3) Phải dùng đúng định luật cho tình huống. (4) Máy biến áp lí tưởng thỏa U2/U1 = N2/N1 và U1I1 = U2I2.
\shortans{3}
\loigiai{
Đối chiếu với kiến thức: Máy biến áp lí tưởng thỏa U2/U1 = N2/N1 và U1I1 = U2I2. Có 3 phát biểu đúng.
}
\end{ex}

% ===== Câu 165 =====
% ID: L12C3B18-36-TN
% Mức: TH
\dangbt{Nhận biết dòng điện Foucault và các ứng dụng của cảm ứng điện từ}
\begin{ex}
% ID: L12C3B18-36-TN
% Mức: TH
Một học sinh ôn tập Nhận biết dòng điện Foucault và các ứng dụng của cảm ứng điện từ?
\choice
{\True Dòng điện Foucault là dòng điện cảm ứng xuất hiện trong khối vật dẫn khi từ thông qua nó biến thiên.}
{Dòng điện Foucault chỉ xuất hiện trong chất cách điện.}
{Có thể bỏ qua phương, chiều và đơn vị trong mọi trường hợp.}
{Kết quả không phụ thuộc dữ kiện và điều kiện áp dụng.}
\loigiai{
Dòng điện Foucault là dòng điện cảm ứng xuất hiện trong khối vật dẫn khi từ thông qua nó biến thiên. Vì vậy phương án $A$ đúng; các phương án còn lại sai.
}
\end{ex}

% ===== Câu 166 =====
% ID: L12C3B18-37-DS
% Mức: TH
\dangbt{Giải thích nguyên lí bếp từ, phanh điện từ và gia nhiệt cảm ứng}
\begin{ex}
% ID: L12C3B18-37-DS
% Mức: TH
Xét Giải thích nguyên lí bếp từ, phanh điện từ và gia nhiệt cảm ứng (Tình huống 2). Hãy xác định tính đúng, sai của bốn phát biểu sau.
\choiceTF
{Bếp từ làm nóng trực tiếp mọi vật liệu bằng điện trường tĩnh.}
{\True Bếp từ làm nóng đáy nồi dẫn điện nhờ dòng điện Foucault và sự tỏa nhiệt Joule.}
{Kết quả không phụ thuộc dữ kiện và có thể bỏ qua điều kiện.}
{\True Khi giải cần xác định đúng dữ kiện, phương chiều, điều kiện áp dụng và đơn vị.}
\loigiai{
\begin{itemchoice}
\item (Sai) Bếp từ làm nóng trực tiếp mọi vật liệu bằng điện trường tĩnh. (Vì): Bếp từ làm nóng trực tiếp mọi vật liệu bằng điện trường tĩnh. b) Đúng: Bếp từ làm nóng đáy nồi dẫn điện nhờ dòng điện Foucault và sự tỏa nhiệt Joule. c) Sai: Kết quả không phụ thuộc dữ kiện và có thể bỏ qua điều kiện. d) Đúng: Khi giải cần xác định đúng dữ kiện, phương chiều, điều kiện áp dụng và đơn vị. Kiến thức cốt lõi: Bếp từ làm nóng đáy nồi dẫn điện nhờ dòng điện Foucault và sự tỏa nhiệt Joule
\item (Đúng) Bếp từ làm nóng đáy nồi dẫn điện nhờ dòng điện Foucault và sự tỏa nhiệt Joule.
\item (Sai) Kết quả không phụ thuộc dữ kiện và có thể bỏ qua điều kiện.
\item (Đúng) Khi giải cần xác định đúng dữ kiện, phương chiều, điều kiện áp dụng và đơn vị.
\end{itemchoice}
}
\end{ex}

% ===== Câu 167 =====
% ID: L12C3B18-37-TL
% Mức: TH
\dangbt{Nhận biết dòng điện Foucault và các ứng dụng của cảm ứng điện từ}
\begin{bt}
% ID: L12C3B18-37-TL
% Mức: TH
Trình bày kiến thức cốt lõi của Nhận biết dòng điện Foucault và các ứng dụng của cảm ứng điện từ và nêu điều kiện áp dụng.
\loigiai{
Cần nêu được: Dòng điện Foucault là dòng điện cảm ứng xuất hiện trong khối vật dẫn khi từ thông qua nó biến thiên. Khi vận dụng phải xác định đúng dữ kiện, phương chiều nếu có, điều kiện áp dụng, hệ đơn vị và kiểm tra tính hợp lí. Nhận định “Dòng điện Foucault chỉ xuất hiện trong chất cách điện.” sai.
}
\end{bt}

% ===== Câu 168 =====
% ID: L12C3B18-37-TLN
% Mức: TH
\begin{ex}
% ID: L12C3B18-37-TLN
% Mức: TH
Trong bốn phát biểu về Nhận biết dòng điện Foucault và các ứng dụng của cảm ứng điện từ, có bao nhiêu phát biểu đúng? (Chỉ ghi một số nguyên.) (1) Dòng điện Foucault là dòng điện cảm ứng xuất hiện trong khối vật dẫn khi từ thông qua nó biến thiên. (2) Dòng điện Foucault chỉ xuất hiện trong chất cách điện. (3) Cần kiểm tra điều kiện áp dụng. (4) Có thể bỏ qua đơn vị.
\shortans{2}
\loigiai{
Đối chiếu với kiến thức: Dòng điện Foucault là dòng điện cảm ứng xuất hiện trong khối vật dẫn khi từ thông qua nó biến thiên. Có 2 phát biểu đúng.
}
\end{ex}

% ===== Câu 169 =====
% ID: L12C3B18-37-TN
% Mức: TH
\begin{ex}
% ID: L12C3B18-37-TN
% Mức: TH
Chọn phát biểu không trái với kiến thức về Nhận biết dòng điện Foucault và các ứng dụng của cảm ứng điện từ?
\choice
{Kết quả không phụ thuộc dữ kiện và điều kiện áp dụng.}
{\True Dòng điện Foucault là dòng điện cảm ứng xuất hiện trong khối vật dẫn khi từ thông qua nó biến thiên.}
{Dòng điện Foucault chỉ xuất hiện trong chất cách điện.}
{Có thể bỏ qua phương, chiều và đơn vị trong mọi trường hợp.}
\loigiai{
Dòng điện Foucault là dòng điện cảm ứng xuất hiện trong khối vật dẫn khi từ thông qua nó biến thiên. Vì vậy phương án $B$ đúng; các phương án còn lại sai.
}
\end{ex}

% ===== Câu 170 =====
% ID: L12C3B18-38-DS
% Mức: VD
\dangbt{Giải thích nguyên lí bếp từ, phanh điện từ và gia nhiệt cảm ứng}
\begin{ex}
% ID: L12C3B18-38-DS
% Mức: VD
Xét Giải thích nguyên lí bếp từ, phanh điện từ và gia nhiệt cảm ứng (Tình huống 3). Hãy xác định tính đúng, sai của bốn phát biểu sau.
\choiceTF
{\True Khi giải cần xác định đúng dữ kiện, phương chiều, điều kiện áp dụng và đơn vị.}
{Kết quả không phụ thuộc dữ kiện và có thể bỏ qua điều kiện.}
{\True Bếp từ làm nóng đáy nồi dẫn điện nhờ dòng điện Foucault và sự tỏa nhiệt Joule.}
{Bếp từ làm nóng trực tiếp mọi vật liệu bằng điện trường tĩnh.}
\loigiai{
\begin{itemchoice}
\item (Đúng) Khi giải cần xác định đúng dữ kiện, phương chiều, điều kiện áp dụng và đơn vị. (Vì): Khi giải cần xác định đúng dữ kiện, phương chiều, điều kiện áp dụng và đơn vị. b) Sai: Kết quả không phụ thuộc dữ kiện và có thể bỏ qua điều kiện. c) Đúng: Bếp từ làm nóng đáy nồi dẫn điện nhờ dòng điện Foucault và sự tỏa nhiệt Joule. d) Sai: Bếp từ làm nóng trực tiếp mọi vật liệu bằng điện trường tĩnh. Kiến thức cốt lõi: Bếp từ làm nóng đáy nồi dẫn điện nhờ dòng điện Foucault và sự tỏa nhiệt Joule
\item (Sai) Kết quả không phụ thuộc dữ kiện và có thể bỏ qua điều kiện.
\item (Đúng) Bếp từ làm nóng đáy nồi dẫn điện nhờ dòng điện Foucault và sự tỏa nhiệt Joule.
\item (Sai) Bếp từ làm nóng trực tiếp mọi vật liệu bằng điện trường tĩnh.
\end{itemchoice}
}
\end{ex}

% ===== Câu 171 =====
% ID: L12C3B18-38-TL
% Mức: TH
\dangbt{Nhận biết dòng điện Foucault và các ứng dụng của cảm ứng điện từ}
\begin{bt}
% ID: L12C3B18-38-TL
% Mức: TH
Giải thích vì sao nhận định sau đúng: “Dòng điện Foucault là dòng điện cảm ứng xuất hiện trong khối vật dẫn khi từ thông qua nó biến thiên.”
\loigiai{
Cần nêu được: Dòng điện Foucault là dòng điện cảm ứng xuất hiện trong khối vật dẫn khi từ thông qua nó biến thiên. Khi vận dụng phải xác định đúng dữ kiện, phương chiều nếu có, điều kiện áp dụng, hệ đơn vị và kiểm tra tính hợp lí. Nhận định “Dòng điện Foucault chỉ xuất hiện trong chất cách điện.” sai.
}
\end{bt}

% ===== Câu 172 =====
% ID: L12C3B18-38-TLN
% Mức: TH
\begin{ex}
% ID: L12C3B18-38-TLN
% Mức: TH
Trong bốn phát biểu về Nhận biết dòng điện Foucault và các ứng dụng của cảm ứng điện từ, có bao nhiêu phát biểu đúng? (Chỉ ghi một số nguyên.) (1) Dòng điện Foucault chỉ xuất hiện trong chất cách điện. (2) Dòng điện Foucault là dòng điện cảm ứng xuất hiện trong khối vật dẫn khi từ thông qua nó biến thiên. (3) Kết quả phải phù hợp ý nghĩa vật lí. (4) Mọi đại lượng đều là vô hướng.
\shortans{1}
\loigiai{
Đối chiếu với kiến thức: Dòng điện Foucault là dòng điện cảm ứng xuất hiện trong khối vật dẫn khi từ thông qua nó biến thiên. Có 1 phát biểu đúng.
}
\end{ex}

% ===== Câu 173 =====
% ID: L12C3B18-38-TN
% Mức: VD
\begin{ex}
% ID: L12C3B18-38-TN
% Mức: VD
Cơ sở Vật lí đúng dùng để giải Nhận biết dòng điện Foucault và các ứng dụng của cảm ứng điện từ?
\choice
{Có thể bỏ qua phương, chiều và đơn vị trong mọi trường hợp.}
{Kết quả không phụ thuộc dữ kiện và điều kiện áp dụng.}
{\True Dòng điện Foucault là dòng điện cảm ứng xuất hiện trong khối vật dẫn khi từ thông qua nó biến thiên.}
{Dòng điện Foucault chỉ xuất hiện trong chất cách điện.}
\loigiai{
Dòng điện Foucault là dòng điện cảm ứng xuất hiện trong khối vật dẫn khi từ thông qua nó biến thiên. Vì vậy phương án $C$ đúng; các phương án còn lại sai.
}
\end{ex}

% ===== Câu 174 =====
% ID: L12C3B18-39-DS
% Mức: VD
\dangbt{Giải thích nguyên lí bếp từ, phanh điện từ và gia nhiệt cảm ứng}
\begin{ex}
% ID: L12C3B18-39-DS
% Mức: VD
Xét Giải thích nguyên lí bếp từ, phanh điện từ và gia nhiệt cảm ứng (Tình huống 4). Hãy xác định tính đúng, sai của bốn phát biểu sau.
\choiceTF
{Kết quả không phụ thuộc dữ kiện và có thể bỏ qua điều kiện.}
{\True Khi giải cần xác định đúng dữ kiện, phương chiều, điều kiện áp dụng và đơn vị.}
{Bếp từ làm nóng trực tiếp mọi vật liệu bằng điện trường tĩnh.}
{\True Bếp từ làm nóng đáy nồi dẫn điện nhờ dòng điện Foucault và sự tỏa nhiệt Joule.}
\loigiai{
\begin{itemchoice}
\item (Sai) Kết quả không phụ thuộc dữ kiện và có thể bỏ qua điều kiện. (Vì): Kết quả không phụ thuộc dữ kiện và có thể bỏ qua điều kiện. b) Đúng: Khi giải cần xác định đúng dữ kiện, phương chiều, điều kiện áp dụng và đơn vị. c) Sai: Bếp từ làm nóng trực tiếp mọi vật liệu bằng điện trường tĩnh. d) Đúng: Bếp từ làm nóng đáy nồi dẫn điện nhờ dòng điện Foucault và sự tỏa nhiệt Joule. Kiến thức cốt lõi: Bếp từ làm nóng đáy nồi dẫn điện nhờ dòng điện Foucault và sự tỏa nhiệt Joule
\item (Đúng) Khi giải cần xác định đúng dữ kiện, phương chiều, điều kiện áp dụng và đơn vị.
\item (Sai) Bếp từ làm nóng trực tiếp mọi vật liệu bằng điện trường tĩnh.
\item (Đúng) Bếp từ làm nóng đáy nồi dẫn điện nhờ dòng điện Foucault và sự tỏa nhiệt Joule.
\end{itemchoice}
}
\end{ex}

% ===== Câu 175 =====
% ID: L12C3B18-39-TL
% Mức: VD
\dangbt{Nhận biết dòng điện Foucault và các ứng dụng của cảm ứng điện từ}
\begin{bt}
% ID: L12C3B18-39-TL
% Mức: VD
Phân tích sai lầm trong nhận định: “Dòng điện Foucault chỉ xuất hiện trong chất cách điện.”
\loigiai{
Cần nêu được: Dòng điện Foucault là dòng điện cảm ứng xuất hiện trong khối vật dẫn khi từ thông qua nó biến thiên. Khi vận dụng phải xác định đúng dữ kiện, phương chiều nếu có, điều kiện áp dụng, hệ đơn vị và kiểm tra tính hợp lí. Nhận định “Dòng điện Foucault chỉ xuất hiện trong chất cách điện.” sai.
}
\end{bt}

% ===== Câu 176 =====
% ID: L12C3B18-39-TLN
% Mức: TH
\begin{ex}
% ID: L12C3B18-39-TLN
% Mức: TH
Trong bốn phát biểu về Nhận biết dòng điện Foucault và các ứng dụng của cảm ứng điện từ, có bao nhiêu phát biểu đúng? (Chỉ ghi một số nguyên.) (1) Dòng điện Foucault là dòng điện cảm ứng xuất hiện trong khối vật dẫn khi từ thông qua nó biến thiên. (2) Dòng điện Foucault chỉ xuất hiện trong chất cách điện. (3) Cần kiểm tra điều kiện áp dụng. (4) Có thể bỏ qua đơn vị.
\shortans{2}
\loigiai{
Đối chiếu với kiến thức: Dòng điện Foucault là dòng điện cảm ứng xuất hiện trong khối vật dẫn khi từ thông qua nó biến thiên. Có 2 phát biểu đúng.
}
\end{ex}

% ===== Câu 177 =====
% ID: L12C3B18-39-TN
% Mức: VD
\begin{ex}
% ID: L12C3B18-39-TN
% Mức: VD
Trong một bài thực hành về Nhận biết dòng điện Foucault và các ứng dụng của cảm ứng điện từ?
\choice
{Dòng điện Foucault chỉ xuất hiện trong chất cách điện.}
{Có thể bỏ qua phương, chiều và đơn vị trong mọi trường hợp.}
{Kết quả không phụ thuộc dữ kiện và điều kiện áp dụng.}
{\True Dòng điện Foucault là dòng điện cảm ứng xuất hiện trong khối vật dẫn khi từ thông qua nó biến thiên.}
\loigiai{
Dòng điện Foucault là dòng điện cảm ứng xuất hiện trong khối vật dẫn khi từ thông qua nó biến thiên. Vì vậy phương án $D$ đúng; các phương án còn lại sai.
}
\end{ex}

% ===== Câu 178 =====
% ID: L12C3B18-40-DS
% Mức: VD
\dangbt{Giải thích nguyên lí bếp từ, phanh điện từ và gia nhiệt cảm ứng}
\begin{ex}
% ID: L12C3B18-40-DS
% Mức: VD
Xét Giải thích nguyên lí bếp từ, phanh điện từ và gia nhiệt cảm ứng (Tình huống 5). Hãy xác định tính đúng, sai của bốn phát biểu sau.
\choiceTF
{\True Bếp từ làm nóng đáy nồi dẫn điện nhờ dòng điện Foucault và sự tỏa nhiệt Joule.}
{Kết quả không phụ thuộc dữ kiện và có thể bỏ qua điều kiện.}
{Bếp từ làm nóng trực tiếp mọi vật liệu bằng điện trường tĩnh.}
{\True Khi giải cần xác định đúng dữ kiện, phương chiều, điều kiện áp dụng và đơn vị.}
\loigiai{
\begin{itemchoice}
\item (Đúng) Bếp từ làm nóng đáy nồi dẫn điện nhờ dòng điện Foucault và sự tỏa nhiệt Joule. (Vì): Bếp từ làm nóng đáy nồi dẫn điện nhờ dòng điện Foucault và sự tỏa nhiệt Joule. b) Sai: Kết quả không phụ thuộc dữ kiện và có thể bỏ qua điều kiện. c) Sai: Bếp từ làm nóng trực tiếp mọi vật liệu bằng điện trường tĩnh. d) Đúng: Khi giải cần xác định đúng dữ kiện, phương chiều, điều kiện áp dụng và đơn vị. Kiến thức cốt lõi: Bếp từ làm nóng đáy nồi dẫn điện nhờ dòng điện Foucault và sự tỏa nhiệt Joule
\item (Sai) Kết quả không phụ thuộc dữ kiện và có thể bỏ qua điều kiện.
\item (Sai) Bếp từ làm nóng trực tiếp mọi vật liệu bằng điện trường tĩnh.
\item (Đúng) Khi giải cần xác định đúng dữ kiện, phương chiều, điều kiện áp dụng và đơn vị.
\end{itemchoice}
}
\end{ex}

% ===== Câu 179 =====
% ID: L12C3B18-40-TL
% Mức: VD
\dangbt{Nhận biết dòng điện Foucault và các ứng dụng của cảm ứng điện từ}
\begin{bt}
% ID: L12C3B18-40-TL
% Mức: VD
Đề xuất các bước giải một bài toán thuộc dạng Nhận biết dòng điện Foucault và các ứng dụng của cảm ứng điện từ.
\loigiai{
Cần nêu được: Dòng điện Foucault là dòng điện cảm ứng xuất hiện trong khối vật dẫn khi từ thông qua nó biến thiên. Khi vận dụng phải xác định đúng dữ kiện, phương chiều nếu có, điều kiện áp dụng, hệ đơn vị và kiểm tra tính hợp lí. Nhận định “Dòng điện Foucault chỉ xuất hiện trong chất cách điện.” sai.
}
\end{bt}

% ===== Câu 180 =====
% ID: L12C3B18-40-TLN
% Mức: VD
\begin{ex}
% ID: L12C3B18-40-TLN
% Mức: VD
Trong bốn phát biểu về Nhận biết dòng điện Foucault và các ứng dụng của cảm ứng điện từ, có bao nhiêu phát biểu đúng? (Chỉ ghi một số nguyên.) (1) Dòng điện Foucault chỉ xuất hiện trong chất cách điện. (2) Dòng điện Foucault là dòng điện cảm ứng xuất hiện trong khối vật dẫn khi từ thông qua nó biến thiên. (3) Kết quả phải phù hợp ý nghĩa vật lí. (4) Mọi đại lượng đều là vô hướng.
\shortans{2}
\loigiai{
Đối chiếu với kiến thức: Dòng điện Foucault là dòng điện cảm ứng xuất hiện trong khối vật dẫn khi từ thông qua nó biến thiên. Có 2 phát biểu đúng.
}
\end{ex}

% ===== Câu 181 =====
% ID: L12C3B18-40-TN
% Mức: VD
\begin{ex}
% ID: L12C3B18-40-TN
% Mức: VD
Kết luận đúng khi vận dụng Nhận biết dòng điện Foucault và các ứng dụng của cảm ứng điện từ?
\choice
{\True Dòng điện Foucault là dòng điện cảm ứng xuất hiện trong khối vật dẫn khi từ thông qua nó biến thiên.}
{Dòng điện Foucault chỉ xuất hiện trong chất cách điện.}
{Có thể bỏ qua phương, chiều và đơn vị trong mọi trường hợp.}
{Kết quả không phụ thuộc dữ kiện và điều kiện áp dụng.}
\loigiai{
Dòng điện Foucault là dòng điện cảm ứng xuất hiện trong khối vật dẫn khi từ thông qua nó biến thiên. Vì vậy phương án $A$ đúng; các phương án còn lại sai.
}
\end{ex}

% ===== Câu 182 =====
% ID: L12C3B18-41-DS
% Mức: TH
\dangbt{Nhận biết cấu tạo, nguyên lí hoạt động của máy biến áp}
\begin{ex}
% ID: L12C3B18-41-DS
% Mức: TH
Xét Nhận biết cấu tạo, nguyên lí hoạt động của máy biến áp (Tình huống 1). Hãy xác định tính đúng, sai của bốn phát biểu sau.
\choiceTF
{\True Máy biến áp gồm lõi sắt và hai cuộn dây, hoạt động dựa trên cảm ứng điện từ với dòng điện xoay chiều.}
{Máy biến áp lí tưởng hoạt động được với nguồn điện một chiều không đổi.}
{\True Khi giải cần xác định đúng dữ kiện, phương chiều, điều kiện áp dụng và đơn vị.}
{Kết quả không phụ thuộc dữ kiện và có thể bỏ qua điều kiện.}
\loigiai{
\begin{itemchoice}
\item (Đúng) Máy biến áp gồm lõi sắt và hai cuộn dây, hoạt động dựa trên cảm ứng điện từ với dòng điện xoay chiều. (Vì): Máy biến áp gồm lõi sắt và hai cuộn dây, hoạt động dựa trên cảm ứng điện từ với dòng điện xoay chiều. b) Sai: Máy biến áp lí tưởng hoạt động được với nguồn điện một chiều không đổi. c) Đúng: Khi giải cần xác định đúng dữ kiện, phương chiều, điều kiện áp dụng và đơn vị. d) Sai: Kết quả không phụ thuộc dữ kiện và có thể bỏ qua điều kiện. Kiến thức cốt lõi: Máy biến áp gồm lõi sắt và hai cuộn dây, hoạt động dựa trên cảm ứng điện từ với dòng điện xoay chiều
\item (Sai) Máy biến áp lí tưởng hoạt động được với nguồn điện một chiều không đổi.
\item (Đúng) Khi giải cần xác định đúng dữ kiện, phương chiều, điều kiện áp dụng và đơn vị.
\item (Sai) Kết quả không phụ thuộc dữ kiện và có thể bỏ qua điều kiện.
\end{itemchoice}
}
\end{ex}

% ===== Câu 183 =====
% ID: L12C3B18-41-TL
% Mức: VD
\dangbt{Nhận biết dòng điện Foucault và các ứng dụng của cảm ứng điện từ}
\begin{bt}
% ID: L12C3B18-41-TL
% Mức: VD
Nêu một tình huống thực tiễn liên quan đến Nhận biết dòng điện Foucault và các ứng dụng của cảm ứng điện từ và giải thích bằng kiến thức Vật lí.
\loigiai{
Cần nêu được: Dòng điện Foucault là dòng điện cảm ứng xuất hiện trong khối vật dẫn khi từ thông qua nó biến thiên. Khi vận dụng phải xác định đúng dữ kiện, phương chiều nếu có, điều kiện áp dụng, hệ đơn vị và kiểm tra tính hợp lí. Nhận định “Dòng điện Foucault chỉ xuất hiện trong chất cách điện.” sai.
}
\end{bt}

% ===== Câu 184 =====
% ID: L12C3B18-41-TLN
% Mức: VD
\begin{ex}
% ID: L12C3B18-41-TLN
% Mức: VD
Trong bốn phát biểu về Nhận biết dòng điện Foucault và các ứng dụng của cảm ứng điện từ, có bao nhiêu phát biểu đúng? (Chỉ ghi một số nguyên.) (1) Dòng điện Foucault là dòng điện cảm ứng xuất hiện trong khối vật dẫn khi từ thông qua nó biến thiên. (2) Dòng điện Foucault chỉ xuất hiện trong chất cách điện. (3) Cần kiểm tra điều kiện áp dụng. (4) Có thể bỏ qua đơn vị.
\shortans{2}
\loigiai{
Đối chiếu với kiến thức: Dòng điện Foucault là dòng điện cảm ứng xuất hiện trong khối vật dẫn khi từ thông qua nó biến thiên. Có 2 phát biểu đúng.
}
\end{ex}

% ===== Câu 185 =====
% ID: L12C3B18-41-TN
% Mức: NB
\dangbt{Truyền tải điện năng đi xa và giảm hao phí trên đường dây}
\begin{ex}
% ID: L12C3B18-41-TN
% Mức: NB
Phát biểu nào sau đây đúng về Truyền tải điện năng đi xa và giảm hao phí trên đường dây?
\choice
{Muốn giảm hao phí phải giảm điện áp truyền tải.}
{Có thể bỏ qua phương, chiều và đơn vị trong mọi trường hợp.}
{Kết quả không phụ thuộc dữ kiện và điều kiện áp dụng.}
{\True Với công suất truyền không đổi, tăng điện áp truyền tải làm giảm dòng điện và giảm hao phí I^2R.}
\loigiai{
Với công suất truyền không đổi, tăng điện áp truyền tải làm giảm dòng điện và giảm hao phí I^2R. Vì vậy phương án $D$ đúng; các phương án còn lại sai.
}
\end{ex}

% ===== Câu 186 =====
% ID: L12C3B18-42-DS
% Mức: TH
\dangbt{Nhận biết cấu tạo, nguyên lí hoạt động của máy biến áp}
\begin{ex}
% ID: L12C3B18-42-DS
% Mức: TH
Xét Nhận biết cấu tạo, nguyên lí hoạt động của máy biến áp (Tình huống 2). Hãy xác định tính đúng, sai của bốn phát biểu sau.
\choiceTF
{Máy biến áp lí tưởng hoạt động được với nguồn điện một chiều không đổi.}
{\True Máy biến áp gồm lõi sắt và hai cuộn dây, hoạt động dựa trên cảm ứng điện từ với dòng điện xoay chiều.}
{Kết quả không phụ thuộc dữ kiện và có thể bỏ qua điều kiện.}
{\True Khi giải cần xác định đúng dữ kiện, phương chiều, điều kiện áp dụng và đơn vị.}
\loigiai{
\begin{itemchoice}
\item (Sai) Máy biến áp lí tưởng hoạt động được với nguồn điện một chiều không đổi. (Vì): Máy biến áp lí tưởng hoạt động được với nguồn điện một chiều không đổi. b) Đúng: Máy biến áp gồm lõi sắt và hai cuộn dây, hoạt động dựa trên cảm ứng điện từ với dòng điện xoay chiều. c) Sai: Kết quả không phụ thuộc dữ kiện và có thể bỏ qua điều kiện. d) Đúng: Khi giải cần xác định đúng dữ kiện, phương chiều, điều kiện áp dụng và đơn vị. Kiến thức cốt lõi: Máy biến áp gồm lõi sắt và hai cuộn dây, hoạt động dựa trên cảm ứng điện từ với dòng điện xoay chiều
\item (Đúng) Máy biến áp gồm lõi sắt và hai cuộn dây, hoạt động dựa trên cảm ứng điện từ với dòng điện xoay chiều.
\item (Sai) Kết quả không phụ thuộc dữ kiện và có thể bỏ qua điều kiện.
\item (Đúng) Khi giải cần xác định đúng dữ kiện, phương chiều, điều kiện áp dụng và đơn vị.
\end{itemchoice}
}
\end{ex}

% ===== Câu 187 =====
% ID: L12C3B18-42-TL
% Mức: VDC
\dangbt{Nhận biết dòng điện Foucault và các ứng dụng của cảm ứng điện từ}
\begin{bt}
% ID: L12C3B18-42-TL
% Mức: VDC
So sánh nhận định đúng “Dòng điện Foucault là dòng điện cảm ứng xuất hiện trong khối vật dẫn khi từ thông qua nó biến thiên.” với nhận định sai “Dòng điện Foucault chỉ xuất hiện trong chất cách điện.”; chỉ rõ căn cứ Vật lí.
\loigiai{
Cần nêu được: Dòng điện Foucault là dòng điện cảm ứng xuất hiện trong khối vật dẫn khi từ thông qua nó biến thiên. Khi vận dụng phải xác định đúng dữ kiện, phương chiều nếu có, điều kiện áp dụng, hệ đơn vị và kiểm tra tính hợp lí. Nhận định “Dòng điện Foucault chỉ xuất hiện trong chất cách điện.” sai.
}
\end{bt}

% ===== Câu 188 =====
% ID: L12C3B18-42-TLN
% Mức: VDC
\begin{ex}
% ID: L12C3B18-42-TLN
% Mức: VDC
Trong bốn phát biểu về Nhận biết dòng điện Foucault và các ứng dụng của cảm ứng điện từ, có bao nhiêu phát biểu đúng? (Chỉ ghi một số nguyên.) (1) Khi giải cần xác định đúng dữ kiện, phương chiều, điều kiện áp dụng và đơn vị. (2) Dòng điện Foucault chỉ xuất hiện trong chất cách điện. (3) Phải dùng đúng định luật cho tình huống. (4) Dòng điện Foucault là dòng điện cảm ứng xuất hiện trong khối vật dẫn khi từ thông qua nó biến thiên.
\shortans{3}
\loigiai{
Đối chiếu với kiến thức: Dòng điện Foucault là dòng điện cảm ứng xuất hiện trong khối vật dẫn khi từ thông qua nó biến thiên. Có 3 phát biểu đúng.
}
\end{ex}

% ===== Câu 189 =====
% ID: L12C3B18-42-TN
% Mức: NB
\dangbt{Truyền tải điện năng đi xa và giảm hao phí trên đường dây}
\begin{ex}
% ID: L12C3B18-42-TN
% Mức: NB
Trong nội dung Truyền tải điện năng đi xa và giảm hao phí trên đường dây?
\choice
{\True Với công suất truyền không đổi, tăng điện áp truyền tải làm giảm dòng điện và giảm hao phí I^2R.}
{Muốn giảm hao phí phải giảm điện áp truyền tải.}
{Có thể bỏ qua phương, chiều và đơn vị trong mọi trường hợp.}
{Kết quả không phụ thuộc dữ kiện và điều kiện áp dụng.}
\loigiai{
Với công suất truyền không đổi, tăng điện áp truyền tải làm giảm dòng điện và giảm hao phí I^2R. Vì vậy phương án $A$ đúng; các phương án còn lại sai.
}
\end{ex}

% ===== Câu 190 =====
% ID: L12C3B18-43-DS
% Mức: VD
\dangbt{Nhận biết cấu tạo, nguyên lí hoạt động của máy biến áp}
\begin{ex}
% ID: L12C3B18-43-DS
% Mức: VD
Xét Nhận biết cấu tạo, nguyên lí hoạt động của máy biến áp (Tình huống 3). Hãy xác định tính đúng, sai của bốn phát biểu sau.
\choiceTF
{\True Khi giải cần xác định đúng dữ kiện, phương chiều, điều kiện áp dụng và đơn vị.}
{Kết quả không phụ thuộc dữ kiện và có thể bỏ qua điều kiện.}
{\True Máy biến áp gồm lõi sắt và hai cuộn dây, hoạt động dựa trên cảm ứng điện từ với dòng điện xoay chiều.}
{Máy biến áp lí tưởng hoạt động được với nguồn điện một chiều không đổi.}
\loigiai{
\begin{itemchoice}
\item (Đúng) Khi giải cần xác định đúng dữ kiện, phương chiều, điều kiện áp dụng và đơn vị. (Vì): Khi giải cần xác định đúng dữ kiện, phương chiều, điều kiện áp dụng và đơn vị. b) Sai: Kết quả không phụ thuộc dữ kiện và có thể bỏ qua điều kiện. c) Đúng: Máy biến áp gồm lõi sắt và hai cuộn dây, hoạt động dựa trên cảm ứng điện từ với dòng điện xoay chiều. d) Sai: Máy biến áp lí tưởng hoạt động được với nguồn điện một chiều không đổi. Kiến thức cốt lõi: Máy biến áp gồm lõi sắt và hai cuộn dây, hoạt động dựa trên cảm ứng điện từ với dòng điện xoay chiều
\item (Sai) Kết quả không phụ thuộc dữ kiện và có thể bỏ qua điều kiện.
\item (Đúng) Máy biến áp gồm lõi sắt và hai cuộn dây, hoạt động dựa trên cảm ứng điện từ với dòng điện xoay chiều.
\item (Sai) Máy biến áp lí tưởng hoạt động được với nguồn điện một chiều không đổi.
\end{itemchoice}
}
\end{ex}

% ===== Câu 191 =====
% ID: L12C3B18-43-TL
% Mức: TH
\dangbt{Giải thích nguyên lí bếp từ, phanh điện từ và gia nhiệt cảm ứng}
\begin{bt}
% ID: L12C3B18-43-TL
% Mức: TH
Trình bày kiến thức cốt lõi của Giải thích nguyên lí bếp từ, phanh điện từ và gia nhiệt cảm ứng và nêu điều kiện áp dụng.
\loigiai{
Cần nêu được: Bếp từ làm nóng đáy nồi dẫn điện nhờ dòng điện Foucault và sự tỏa nhiệt Joule. Khi vận dụng phải xác định đúng dữ kiện, phương chiều nếu có, điều kiện áp dụng, hệ đơn vị và kiểm tra tính hợp lí. Nhận định “Bếp từ làm nóng trực tiếp mọi vật liệu bằng điện trường tĩnh.” sai.
}
\end{bt}

% ===== Câu 192 =====
% ID: L12C3B18-43-TLN
% Mức: TH
\begin{ex}
% ID: L12C3B18-43-TLN
% Mức: TH
Trong bốn phát biểu về Giải thích nguyên lí bếp từ, phanh điện từ và gia nhiệt cảm ứng, có bao nhiêu phát biểu đúng? (Chỉ ghi một số nguyên.) (1) Bếp từ làm nóng đáy nồi dẫn điện nhờ dòng điện Foucault và sự tỏa nhiệt Joule. (2) Bếp từ làm nóng trực tiếp mọi vật liệu bằng điện trường tĩnh. (3) Cần kiểm tra điều kiện áp dụng. (4) Có thể bỏ qua đơn vị.
\shortans{2}
\loigiai{
Đối chiếu với kiến thức: Bếp từ làm nóng đáy nồi dẫn điện nhờ dòng điện Foucault và sự tỏa nhiệt Joule. Có 2 phát biểu đúng.
}
\end{ex}

% ===== Câu 193 =====
% ID: L12C3B18-43-TN
% Mức: NB
\dangbt{Truyền tải điện năng đi xa và giảm hao phí trên đường dây}
\begin{ex}
% ID: L12C3B18-43-TN
% Mức: NB
Tình huống 3: chọn kết luận chính xác về Truyền tải điện năng đi xa và giảm hao phí trên đường dây?
\choice
{Kết quả không phụ thuộc dữ kiện và điều kiện áp dụng.}
{\True Với công suất truyền không đổi, tăng điện áp truyền tải làm giảm dòng điện và giảm hao phí I^2R.}
{Muốn giảm hao phí phải giảm điện áp truyền tải.}
{Có thể bỏ qua phương, chiều và đơn vị trong mọi trường hợp.}
\loigiai{
Với công suất truyền không đổi, tăng điện áp truyền tải làm giảm dòng điện và giảm hao phí I^2R. Vì vậy phương án $B$ đúng; các phương án còn lại sai.
}
\end{ex}

% ===== Câu 194 =====
% ID: L12C3B18-44-DS
% Mức: VD
\dangbt{Nhận biết cấu tạo, nguyên lí hoạt động của máy biến áp}
\begin{ex}
% ID: L12C3B18-44-DS
% Mức: VD
Xét Nhận biết cấu tạo, nguyên lí hoạt động của máy biến áp (Tình huống 4). Hãy xác định tính đúng, sai của bốn phát biểu sau.
\choiceTF
{Kết quả không phụ thuộc dữ kiện và có thể bỏ qua điều kiện.}
{\True Khi giải cần xác định đúng dữ kiện, phương chiều, điều kiện áp dụng và đơn vị.}
{Máy biến áp lí tưởng hoạt động được với nguồn điện một chiều không đổi.}
{\True Máy biến áp gồm lõi sắt và hai cuộn dây, hoạt động dựa trên cảm ứng điện từ với dòng điện xoay chiều.}
\loigiai{
\begin{itemchoice}
\item (Sai) Kết quả không phụ thuộc dữ kiện và có thể bỏ qua điều kiện. (Vì): Kết quả không phụ thuộc dữ kiện và có thể bỏ qua điều kiện. b) Đúng: Khi giải cần xác định đúng dữ kiện, phương chiều, điều kiện áp dụng và đơn vị. c) Sai: Máy biến áp lí tưởng hoạt động được với nguồn điện một chiều không đổi. d) Đúng: Máy biến áp gồm lõi sắt và hai cuộn dây, hoạt động dựa trên cảm ứng điện từ với dòng điện xoay chiều. Kiến thức cốt lõi: Máy biến áp gồm lõi sắt và hai cuộn dây, hoạt động dựa trên cảm ứng điện từ với dòng điện xoay chiều
\item (Đúng) Khi giải cần xác định đúng dữ kiện, phương chiều, điều kiện áp dụng và đơn vị.
\item (Sai) Máy biến áp lí tưởng hoạt động được với nguồn điện một chiều không đổi.
\item (Đúng) Máy biến áp gồm lõi sắt và hai cuộn dây, hoạt động dựa trên cảm ứng điện từ với dòng điện xoay chiều.
\end{itemchoice}
}
\end{ex}

% ===== Câu 195 =====
% ID: L12C3B18-44-TL
% Mức: TH
\dangbt{Giải thích nguyên lí bếp từ, phanh điện từ và gia nhiệt cảm ứng}
\begin{bt}
% ID: L12C3B18-44-TL
% Mức: TH
Giải thích vì sao nhận định sau đúng: “Bếp từ làm nóng đáy nồi dẫn điện nhờ dòng điện Foucault và sự tỏa nhiệt Joule.”
\loigiai{
Cần nêu được: Bếp từ làm nóng đáy nồi dẫn điện nhờ dòng điện Foucault và sự tỏa nhiệt Joule. Khi vận dụng phải xác định đúng dữ kiện, phương chiều nếu có, điều kiện áp dụng, hệ đơn vị và kiểm tra tính hợp lí. Nhận định “Bếp từ làm nóng trực tiếp mọi vật liệu bằng điện trường tĩnh.” sai.
}
\end{bt}

% ===== Câu 196 =====
% ID: L12C3B18-44-TLN
% Mức: TH
\begin{ex}
% ID: L12C3B18-44-TLN
% Mức: TH
Trong bốn phát biểu về Giải thích nguyên lí bếp từ, phanh điện từ và gia nhiệt cảm ứng, có bao nhiêu phát biểu đúng? (Chỉ ghi một số nguyên.) (1) Bếp từ làm nóng trực tiếp mọi vật liệu bằng điện trường tĩnh. (2) Bếp từ làm nóng đáy nồi dẫn điện nhờ dòng điện Foucault và sự tỏa nhiệt Joule. (3) Kết quả phải phù hợp ý nghĩa vật lí. (4) Mọi đại lượng đều là vô hướng.
\shortans{1}
\loigiai{
Đối chiếu với kiến thức: Bếp từ làm nóng đáy nồi dẫn điện nhờ dòng điện Foucault và sự tỏa nhiệt Joule. Có 1 phát biểu đúng.
}
\end{ex}

% ===== Câu 197 =====
% ID: L12C3B18-44-TN
% Mức: TH
\dangbt{Truyền tải điện năng đi xa và giảm hao phí trên đường dây}
\begin{ex}
% ID: L12C3B18-44-TN
% Mức: TH
Nội dung nào mô tả đúng nhất Truyền tải điện năng đi xa và giảm hao phí trên đường dây?
\choice
{Có thể bỏ qua phương, chiều và đơn vị trong mọi trường hợp.}
{Kết quả không phụ thuộc dữ kiện và điều kiện áp dụng.}
{\True Với công suất truyền không đổi, tăng điện áp truyền tải làm giảm dòng điện và giảm hao phí I^2R.}
{Muốn giảm hao phí phải giảm điện áp truyền tải.}
\loigiai{
Với công suất truyền không đổi, tăng điện áp truyền tải làm giảm dòng điện và giảm hao phí I^2R. Vì vậy phương án $C$ đúng; các phương án còn lại sai.
}
\end{ex}

% ===== Câu 198 =====
% ID: L12C3B18-45-DS
% Mức: VD
\dangbt{Nhận biết cấu tạo, nguyên lí hoạt động của máy biến áp}
\begin{ex}
% ID: L12C3B18-45-DS
% Mức: VD
Xét Nhận biết cấu tạo, nguyên lí hoạt động của máy biến áp (Tình huống 5). Hãy xác định tính đúng, sai của bốn phát biểu sau.
\choiceTF
{\True Máy biến áp gồm lõi sắt và hai cuộn dây, hoạt động dựa trên cảm ứng điện từ với dòng điện xoay chiều.}
{Kết quả không phụ thuộc dữ kiện và có thể bỏ qua điều kiện.}
{Máy biến áp lí tưởng hoạt động được với nguồn điện một chiều không đổi.}
{\True Khi giải cần xác định đúng dữ kiện, phương chiều, điều kiện áp dụng và đơn vị.}
\loigiai{
\begin{itemchoice}
\item (Đúng) Máy biến áp gồm lõi sắt và hai cuộn dây, hoạt động dựa trên cảm ứng điện từ với dòng điện xoay chiều. (Vì): Máy biến áp gồm lõi sắt và hai cuộn dây, hoạt động dựa trên cảm ứng điện từ với dòng điện xoay chiều. b) Sai: Kết quả không phụ thuộc dữ kiện và có thể bỏ qua điều kiện. c) Sai: Máy biến áp lí tưởng hoạt động được với nguồn điện một chiều không đổi. d) Đúng: Khi giải cần xác định đúng dữ kiện, phương chiều, điều kiện áp dụng và đơn vị. Kiến thức cốt lõi: Máy biến áp gồm lõi sắt và hai cuộn dây, hoạt động dựa trên cảm ứng điện từ với dòng điện xoay chiều
\item (Sai) Kết quả không phụ thuộc dữ kiện và có thể bỏ qua điều kiện.
\item (Sai) Máy biến áp lí tưởng hoạt động được với nguồn điện một chiều không đổi.
\item (Đúng) Khi giải cần xác định đúng dữ kiện, phương chiều, điều kiện áp dụng và đơn vị.
\end{itemchoice}
}
\end{ex}

% ===== Câu 199 =====
% ID: L12C3B18-45-TL
% Mức: VD
\dangbt{Giải thích nguyên lí bếp từ, phanh điện từ và gia nhiệt cảm ứng}
\begin{bt}
% ID: L12C3B18-45-TL
% Mức: VD
Phân tích sai lầm trong nhận định: “Bếp từ làm nóng trực tiếp mọi vật liệu bằng điện trường tĩnh.”
\loigiai{
Cần nêu được: Bếp từ làm nóng đáy nồi dẫn điện nhờ dòng điện Foucault và sự tỏa nhiệt Joule. Khi vận dụng phải xác định đúng dữ kiện, phương chiều nếu có, điều kiện áp dụng, hệ đơn vị và kiểm tra tính hợp lí. Nhận định “Bếp từ làm nóng trực tiếp mọi vật liệu bằng điện trường tĩnh.” sai.
}
\end{bt}

% ===== Câu 200 =====
% ID: L12C3B18-45-TLN
% Mức: TH
\begin{ex}
% ID: L12C3B18-45-TLN
% Mức: TH
Trong bốn phát biểu về Giải thích nguyên lí bếp từ, phanh điện từ và gia nhiệt cảm ứng, có bao nhiêu phát biểu đúng? (Chỉ ghi một số nguyên.) (1) Bếp từ làm nóng đáy nồi dẫn điện nhờ dòng điện Foucault và sự tỏa nhiệt Joule. (2) Bếp từ làm nóng trực tiếp mọi vật liệu bằng điện trường tĩnh. (3) Cần kiểm tra điều kiện áp dụng. (4) Có thể bỏ qua đơn vị.
\shortans{2}
\loigiai{
Đối chiếu với kiến thức: Bếp từ làm nóng đáy nồi dẫn điện nhờ dòng điện Foucault và sự tỏa nhiệt Joule. Có 2 phát biểu đúng.
}
\end{ex}

% ===== Câu 201 =====
% ID: L12C3B18-45-TN
% Mức: TH
\dangbt{Truyền tải điện năng đi xa và giảm hao phí trên đường dây}
\begin{ex}
% ID: L12C3B18-45-TN
% Mức: TH
Khi phân tích Truyền tải điện năng đi xa và giảm hao phí trên đường dây?
\choice
{Muốn giảm hao phí phải giảm điện áp truyền tải.}
{Có thể bỏ qua phương, chiều và đơn vị trong mọi trường hợp.}
{Kết quả không phụ thuộc dữ kiện và điều kiện áp dụng.}
{\True Với công suất truyền không đổi, tăng điện áp truyền tải làm giảm dòng điện và giảm hao phí I^2R.}
\loigiai{
Với công suất truyền không đổi, tăng điện áp truyền tải làm giảm dòng điện và giảm hao phí I^2R. Vì vậy phương án $D$ đúng; các phương án còn lại sai.
}
\end{ex}

% ===== Câu 202 =====
% ID: L12C3B18-46-DS
% Mức: TH
\dangbt{Tính điện áp, số vòng dây và cường độ dòng điện của máy biến áp lí tưởng}
\begin{ex}
% ID: L12C3B18-46-DS
% Mức: TH
Xét Tính điện áp, số vòng dây và cường độ dòng điện của máy biến áp lí tưởng (Tình huống 1). Hãy xác định tính đúng, sai của bốn phát biểu sau.
\choiceTF
{\True Máy biến áp lí tưởng thỏa U2/U1 = N2/N1 và U1I1 = U2I2.}
{Máy biến áp lí tưởng luôn có U2 = U1 bất kể số vòng dây.}
{\True Khi giải cần xác định đúng dữ kiện, phương chiều, điều kiện áp dụng và đơn vị.}
{Kết quả không phụ thuộc dữ kiện và có thể bỏ qua điều kiện.}
\loigiai{
\begin{itemchoice}
\item (Đúng) Máy biến áp lí tưởng thỏa U2/U1 = N2/N1 và U1I1 = U2I2. (Vì): Máy biến áp lí tưởng thỏa U2/U1 = N2/N1 và U1I1 = U2I2. b) Sai: Máy biến áp lí tưởng luôn có U2 = U1 bất kể số vòng dây. c) Đúng: Khi giải cần xác định đúng dữ kiện, phương chiều, điều kiện áp dụng và đơn vị. d) Sai: Kết quả không phụ thuộc dữ kiện và có thể bỏ qua điều kiện. Kiến thức cốt lõi: Máy biến áp lí tưởng thỏa U2/U1 = N2/N1 và U1I1 = U2I2
\item (Sai) Máy biến áp lí tưởng luôn có U2 = U1 bất kể số vòng dây.
\item (Đúng) Khi giải cần xác định đúng dữ kiện, phương chiều, điều kiện áp dụng và đơn vị.
\item (Sai) Kết quả không phụ thuộc dữ kiện và có thể bỏ qua điều kiện.
\end{itemchoice}
}
\end{ex}

% ===== Câu 203 =====
% ID: L12C3B18-46-TL
% Mức: VD
\dangbt{Giải thích nguyên lí bếp từ, phanh điện từ và gia nhiệt cảm ứng}
\begin{bt}
% ID: L12C3B18-46-TL
% Mức: VD
Đề xuất các bước giải một bài toán thuộc dạng Giải thích nguyên lí bếp từ, phanh điện từ và gia nhiệt cảm ứng.
\loigiai{
Cần nêu được: Bếp từ làm nóng đáy nồi dẫn điện nhờ dòng điện Foucault và sự tỏa nhiệt Joule. Khi vận dụng phải xác định đúng dữ kiện, phương chiều nếu có, điều kiện áp dụng, hệ đơn vị và kiểm tra tính hợp lí. Nhận định “Bếp từ làm nóng trực tiếp mọi vật liệu bằng điện trường tĩnh.” sai.
}
\end{bt}

% ===== Câu 204 =====
% ID: L12C3B18-46-TLN
% Mức: VD
\begin{ex}
% ID: L12C3B18-46-TLN
% Mức: VD
Trong bốn phát biểu về Giải thích nguyên lí bếp từ, phanh điện từ và gia nhiệt cảm ứng, có bao nhiêu phát biểu đúng? (Chỉ ghi một số nguyên.) (1) Bếp từ làm nóng trực tiếp mọi vật liệu bằng điện trường tĩnh. (2) Bếp từ làm nóng đáy nồi dẫn điện nhờ dòng điện Foucault và sự tỏa nhiệt Joule. (3) Kết quả phải phù hợp ý nghĩa vật lí. (4) Mọi đại lượng đều là vô hướng.
\shortans{2}
\loigiai{
Đối chiếu với kiến thức: Bếp từ làm nóng đáy nồi dẫn điện nhờ dòng điện Foucault và sự tỏa nhiệt Joule. Có 2 phát biểu đúng.
}
\end{ex}

% ===== Câu 205 =====
% ID: L12C3B18-46-TN
% Mức: TH
\dangbt{Truyền tải điện năng đi xa và giảm hao phí trên đường dây}
\begin{ex}
% ID: L12C3B18-46-TN
% Mức: TH
Một học sinh ôn tập Truyền tải điện năng đi xa và giảm hao phí trên đường dây?
\choice
{\True Với công suất truyền không đổi, tăng điện áp truyền tải làm giảm dòng điện và giảm hao phí I^2R.}
{Muốn giảm hao phí phải giảm điện áp truyền tải.}
{Có thể bỏ qua phương, chiều và đơn vị trong mọi trường hợp.}
{Kết quả không phụ thuộc dữ kiện và điều kiện áp dụng.}
\loigiai{
Với công suất truyền không đổi, tăng điện áp truyền tải làm giảm dòng điện và giảm hao phí I^2R. Vì vậy phương án $A$ đúng; các phương án còn lại sai.
}
\end{ex}

% ===== Câu 206 =====
% ID: L12C3B18-47-DS
% Mức: TH
\dangbt{Tính điện áp, số vòng dây và cường độ dòng điện của máy biến áp lí tưởng}
\begin{ex}
% ID: L12C3B18-47-DS
% Mức: TH
Xét Tính điện áp, số vòng dây và cường độ dòng điện của máy biến áp lí tưởng (Tình huống 2). Hãy xác định tính đúng, sai của bốn phát biểu sau.
\choiceTF
{Máy biến áp lí tưởng luôn có U2 = U1 bất kể số vòng dây.}
{\True Máy biến áp lí tưởng thỏa U2/U1 = N2/N1 và U1I1 = U2I2.}
{Kết quả không phụ thuộc dữ kiện và có thể bỏ qua điều kiện.}
{\True Khi giải cần xác định đúng dữ kiện, phương chiều, điều kiện áp dụng và đơn vị.}
\loigiai{
\begin{itemchoice}
\item (Sai) Máy biến áp lí tưởng luôn có U2 = U1 bất kể số vòng dây. (Vì): Máy biến áp lí tưởng luôn có U2 = U1 bất kể số vòng dây. b) Đúng: Máy biến áp lí tưởng thỏa U2/U1 = N2/N1 và U1I1 = U2I2. c) Sai: Kết quả không phụ thuộc dữ kiện và có thể bỏ qua điều kiện. d) Đúng: Khi giải cần xác định đúng dữ kiện, phương chiều, điều kiện áp dụng và đơn vị. Kiến thức cốt lõi: Máy biến áp lí tưởng thỏa U2/U1 = N2/N1 và U1I1 = U2I2
\item (Đúng) Máy biến áp lí tưởng thỏa U2/U1 = N2/N1 và U1I1 = U2I2.
\item (Sai) Kết quả không phụ thuộc dữ kiện và có thể bỏ qua điều kiện.
\item (Đúng) Khi giải cần xác định đúng dữ kiện, phương chiều, điều kiện áp dụng và đơn vị.
\end{itemchoice}
}
\end{ex}

% ===== Câu 207 =====
% ID: L12C3B18-47-TL
% Mức: VD
\dangbt{Giải thích nguyên lí bếp từ, phanh điện từ và gia nhiệt cảm ứng}
\begin{bt}
% ID: L12C3B18-47-TL
% Mức: VD
Nêu một tình huống thực tiễn liên quan đến Giải thích nguyên lí bếp từ, phanh điện từ và gia nhiệt cảm ứng và giải thích bằng kiến thức Vật lí.
\loigiai{
Cần nêu được: Bếp từ làm nóng đáy nồi dẫn điện nhờ dòng điện Foucault và sự tỏa nhiệt Joule. Khi vận dụng phải xác định đúng dữ kiện, phương chiều nếu có, điều kiện áp dụng, hệ đơn vị và kiểm tra tính hợp lí. Nhận định “Bếp từ làm nóng trực tiếp mọi vật liệu bằng điện trường tĩnh.” sai.
}
\end{bt}

% ===== Câu 208 =====
% ID: L12C3B18-47-TLN
% Mức: VD
\begin{ex}
% ID: L12C3B18-47-TLN
% Mức: VD
Trong bốn phát biểu về Giải thích nguyên lí bếp từ, phanh điện từ và gia nhiệt cảm ứng, có bao nhiêu phát biểu đúng? (Chỉ ghi một số nguyên.) (1) Bếp từ làm nóng đáy nồi dẫn điện nhờ dòng điện Foucault và sự tỏa nhiệt Joule. (2) Bếp từ làm nóng trực tiếp mọi vật liệu bằng điện trường tĩnh. (3) Cần kiểm tra điều kiện áp dụng. (4) Có thể bỏ qua đơn vị.
\shortans{2}
\loigiai{
Đối chiếu với kiến thức: Bếp từ làm nóng đáy nồi dẫn điện nhờ dòng điện Foucault và sự tỏa nhiệt Joule. Có 2 phát biểu đúng.
}
\end{ex}

% ===== Câu 209 =====
% ID: L12C3B18-47-TN
% Mức: TH
\dangbt{Truyền tải điện năng đi xa và giảm hao phí trên đường dây}
\begin{ex}
% ID: L12C3B18-47-TN
% Mức: TH
Chọn phát biểu không trái với kiến thức về Truyền tải điện năng đi xa và giảm hao phí trên đường dây?
\choice
{Kết quả không phụ thuộc dữ kiện và điều kiện áp dụng.}
{\True Với công suất truyền không đổi, tăng điện áp truyền tải làm giảm dòng điện và giảm hao phí I^2R.}
{Muốn giảm hao phí phải giảm điện áp truyền tải.}
{Có thể bỏ qua phương, chiều và đơn vị trong mọi trường hợp.}
\loigiai{
Với công suất truyền không đổi, tăng điện áp truyền tải làm giảm dòng điện và giảm hao phí I^2R. Vì vậy phương án $B$ đúng; các phương án còn lại sai.
}
\end{ex}

% ===== Câu 210 =====
% ID: L12C3B18-48-DS
% Mức: VD
\dangbt{Tính điện áp, số vòng dây và cường độ dòng điện của máy biến áp lí tưởng}
\begin{ex}
% ID: L12C3B18-48-DS
% Mức: VD
Xét Tính điện áp, số vòng dây và cường độ dòng điện của máy biến áp lí tưởng (Tình huống 3). Hãy xác định tính đúng, sai của bốn phát biểu sau.
\choiceTF
{\True Khi giải cần xác định đúng dữ kiện, phương chiều, điều kiện áp dụng và đơn vị.}
{Kết quả không phụ thuộc dữ kiện và có thể bỏ qua điều kiện.}
{\True Máy biến áp lí tưởng thỏa U2/U1 = N2/N1 và U1I1 = U2I2.}
{Máy biến áp lí tưởng luôn có U2 = U1 bất kể số vòng dây.}
\loigiai{
\begin{itemchoice}
\item (Đúng) Khi giải cần xác định đúng dữ kiện, phương chiều, điều kiện áp dụng và đơn vị. (Vì): Khi giải cần xác định đúng dữ kiện, phương chiều, điều kiện áp dụng và đơn vị. b) Sai: Kết quả không phụ thuộc dữ kiện và có thể bỏ qua điều kiện. c) Đúng: Máy biến áp lí tưởng thỏa U2/U1 = N2/N1 và U1I1 = U2I2. d) Sai: Máy biến áp lí tưởng luôn có U2 = U1 bất kể số vòng dây. Kiến thức cốt lõi: Máy biến áp lí tưởng thỏa U2/U1 = N2/N1 và U1I1 = U2I2
\item (Sai) Kết quả không phụ thuộc dữ kiện và có thể bỏ qua điều kiện.
\item (Đúng) Máy biến áp lí tưởng thỏa U2/U1 = N2/N1 và U1I1 = U2I2.
\item (Sai) Máy biến áp lí tưởng luôn có U2 = U1 bất kể số vòng dây.
\end{itemchoice}
}
\end{ex}

% ===== Câu 211 =====
% ID: L12C3B18-48-TL
% Mức: VDC
\dangbt{Giải thích nguyên lí bếp từ, phanh điện từ và gia nhiệt cảm ứng}
\begin{bt}
% ID: L12C3B18-48-TL
% Mức: VDC
So sánh nhận định đúng “Bếp từ làm nóng đáy nồi dẫn điện nhờ dòng điện Foucault và sự tỏa nhiệt Joule.” với nhận định sai “Bếp từ làm nóng trực tiếp mọi vật liệu bằng điện trường tĩnh.”; chỉ rõ căn cứ Vật lí.
\loigiai{
Cần nêu được: Bếp từ làm nóng đáy nồi dẫn điện nhờ dòng điện Foucault và sự tỏa nhiệt Joule. Khi vận dụng phải xác định đúng dữ kiện, phương chiều nếu có, điều kiện áp dụng, hệ đơn vị và kiểm tra tính hợp lí. Nhận định “Bếp từ làm nóng trực tiếp mọi vật liệu bằng điện trường tĩnh.” sai.
}
\end{bt}

% ===== Câu 212 =====
% ID: L12C3B18-48-TLN
% Mức: VDC
\begin{ex}
% ID: L12C3B18-48-TLN
% Mức: VDC
Trong bốn phát biểu về Giải thích nguyên lí bếp từ, phanh điện từ và gia nhiệt cảm ứng, có bao nhiêu phát biểu đúng? (Chỉ ghi một số nguyên.) (1) Khi giải cần xác định đúng dữ kiện, phương chiều, điều kiện áp dụng và đơn vị. (2) Bếp từ làm nóng trực tiếp mọi vật liệu bằng điện trường tĩnh. (3) Phải dùng đúng định luật cho tình huống. (4) Bếp từ làm nóng đáy nồi dẫn điện nhờ dòng điện Foucault và sự tỏa nhiệt Joule.
\shortans{3}
\loigiai{
Đối chiếu với kiến thức: Bếp từ làm nóng đáy nồi dẫn điện nhờ dòng điện Foucault và sự tỏa nhiệt Joule. Có 3 phát biểu đúng.
}
\end{ex}

% ===== Câu 213 =====
% ID: L12C3B18-48-TN
% Mức: VD
\dangbt{Truyền tải điện năng đi xa và giảm hao phí trên đường dây}
\begin{ex}
% ID: L12C3B18-48-TN
% Mức: VD
Cơ sở Vật lí đúng dùng để giải Truyền tải điện năng đi xa và giảm hao phí trên đường dây?
\choice
{Có thể bỏ qua phương, chiều và đơn vị trong mọi trường hợp.}
{Kết quả không phụ thuộc dữ kiện và điều kiện áp dụng.}
{\True Với công suất truyền không đổi, tăng điện áp truyền tải làm giảm dòng điện và giảm hao phí I^2R.}
{Muốn giảm hao phí phải giảm điện áp truyền tải.}
\loigiai{
Với công suất truyền không đổi, tăng điện áp truyền tải làm giảm dòng điện và giảm hao phí I^2R. Vì vậy phương án $C$ đúng; các phương án còn lại sai.
}
\end{ex}

% ===== Câu 214 =====
% ID: L12C3B18-49-DS
% Mức: VD
\dangbt{Tính điện áp, số vòng dây và cường độ dòng điện của máy biến áp lí tưởng}
\begin{ex}
% ID: L12C3B18-49-DS
% Mức: VD
Xét Tính điện áp, số vòng dây và cường độ dòng điện của máy biến áp lí tưởng (Tình huống 4). Hãy xác định tính đúng, sai của bốn phát biểu sau.
\choiceTF
{Kết quả không phụ thuộc dữ kiện và có thể bỏ qua điều kiện.}
{\True Khi giải cần xác định đúng dữ kiện, phương chiều, điều kiện áp dụng và đơn vị.}
{Máy biến áp lí tưởng luôn có U2 = U1 bất kể số vòng dây.}
{\True Máy biến áp lí tưởng thỏa U2/U1 = N2/N1 và U1I1 = U2I2.}
\loigiai{
\begin{itemchoice}
\item (Sai) Kết quả không phụ thuộc dữ kiện và có thể bỏ qua điều kiện. (Vì): Kết quả không phụ thuộc dữ kiện và có thể bỏ qua điều kiện. b) Đúng: Khi giải cần xác định đúng dữ kiện, phương chiều, điều kiện áp dụng và đơn vị. c) Sai: Máy biến áp lí tưởng luôn có U2 = U1 bất kể số vòng dây. d) Đúng: Máy biến áp lí tưởng thỏa U2/U1 = N2/N1 và U1I1 = U2I2. Kiến thức cốt lõi: Máy biến áp lí tưởng thỏa U2/U1 = N2/N1 và U1I1 = U2I2
\item (Đúng) Khi giải cần xác định đúng dữ kiện, phương chiều, điều kiện áp dụng và đơn vị.
\item (Sai) Máy biến áp lí tưởng luôn có U2 = U1 bất kể số vòng dây.
\item (Đúng) Máy biến áp lí tưởng thỏa U2/U1 = N2/N1 và U1I1 = U2I2.
\end{itemchoice}
}
\end{ex}

% ===== Câu 215 =====
% ID: L12C3B18-49-TL
% Mức: TH
\dangbt{Nhận biết cấu tạo, nguyên lí hoạt động của máy biến áp}
\begin{bt}
% ID: L12C3B18-49-TL
% Mức: TH
Trình bày kiến thức cốt lõi của Nhận biết cấu tạo, nguyên lí hoạt động của máy biến áp và nêu điều kiện áp dụng.
\loigiai{
Cần nêu được: Máy biến áp gồm lõi sắt và hai cuộn dây, hoạt động dựa trên cảm ứng điện từ với dòng điện xoay chiều. Khi vận dụng phải xác định đúng dữ kiện, phương chiều nếu có, điều kiện áp dụng, hệ đơn vị và kiểm tra tính hợp lí. Nhận định “Máy biến áp lí tưởng hoạt động được với nguồn điện một chiều không đổi.” sai.
}
\end{bt}

% ===== Câu 216 =====
% ID: L12C3B18-49-TLN
% Mức: TH
\begin{ex}
% ID: L12C3B18-49-TLN
% Mức: TH
Trong bốn phát biểu về Nhận biết cấu tạo, nguyên lí hoạt động của máy biến áp, có bao nhiêu phát biểu đúng? (Chỉ ghi một số nguyên.) (1) Máy biến áp gồm lõi sắt và hai cuộn dây, hoạt động dựa trên cảm ứng điện từ với dòng điện xoay chiều. (2) Máy biến áp lí tưởng hoạt động được với nguồn điện một chiều không đổi. (3) Cần kiểm tra điều kiện áp dụng. (4) Có thể bỏ qua đơn vị.
\shortans{2}
\loigiai{
Đối chiếu với kiến thức: Máy biến áp gồm lõi sắt và hai cuộn dây, hoạt động dựa trên cảm ứng điện từ với dòng điện xoay chiều. Có 2 phát biểu đúng.
}
\end{ex}

% ===== Câu 217 =====
% ID: L12C3B18-49-TN
% Mức: VD
\dangbt{Truyền tải điện năng đi xa và giảm hao phí trên đường dây}
\begin{ex}
% ID: L12C3B18-49-TN
% Mức: VD
Trong một bài thực hành về Truyền tải điện năng đi xa và giảm hao phí trên đường dây?
\choice
{Muốn giảm hao phí phải giảm điện áp truyền tải.}
{Có thể bỏ qua phương, chiều và đơn vị trong mọi trường hợp.}
{Kết quả không phụ thuộc dữ kiện và điều kiện áp dụng.}
{\True Với công suất truyền không đổi, tăng điện áp truyền tải làm giảm dòng điện và giảm hao phí I^2R.}
\loigiai{
Với công suất truyền không đổi, tăng điện áp truyền tải làm giảm dòng điện và giảm hao phí I^2R. Vì vậy phương án $D$ đúng; các phương án còn lại sai.
}
\end{ex}

% ===== Câu 218 =====
% ID: L12C3B18-50-DS
% Mức: VD
\dangbt{Tính điện áp, số vòng dây và cường độ dòng điện của máy biến áp lí tưởng}
\begin{ex}
% ID: L12C3B18-50-DS
% Mức: VD
Xét Tính điện áp, số vòng dây và cường độ dòng điện của máy biến áp lí tưởng (Tình huống 5). Hãy xác định tính đúng, sai của bốn phát biểu sau.
\choiceTF
{\True Máy biến áp lí tưởng thỏa U2/U1 = N2/N1 và U1I1 = U2I2.}
{Kết quả không phụ thuộc dữ kiện và có thể bỏ qua điều kiện.}
{Máy biến áp lí tưởng luôn có U2 = U1 bất kể số vòng dây.}
{\True Khi giải cần xác định đúng dữ kiện, phương chiều, điều kiện áp dụng và đơn vị.}
\loigiai{
\begin{itemchoice}
\item (Đúng) Máy biến áp lí tưởng thỏa U2/U1 = N2/N1 và U1I1 = U2I2. (Vì): Máy biến áp lí tưởng thỏa U2/U1 = N2/N1 và U1I1 = U2I2. b) Sai: Kết quả không phụ thuộc dữ kiện và có thể bỏ qua điều kiện. c) Sai: Máy biến áp lí tưởng luôn có U2 = U1 bất kể số vòng dây. d) Đúng: Khi giải cần xác định đúng dữ kiện, phương chiều, điều kiện áp dụng và đơn vị. Kiến thức cốt lõi: Máy biến áp lí tưởng thỏa U2/U1 = N2/N1 và U1I1 = U2I2
\item (Sai) Kết quả không phụ thuộc dữ kiện và có thể bỏ qua điều kiện.
\item (Sai) Máy biến áp lí tưởng luôn có U2 = U1 bất kể số vòng dây.
\item (Đúng) Khi giải cần xác định đúng dữ kiện, phương chiều, điều kiện áp dụng và đơn vị.
\end{itemchoice}
}
\end{ex}

% ===== Câu 219 =====
% ID: L12C3B18-50-TL
% Mức: TH
\dangbt{Nhận biết cấu tạo, nguyên lí hoạt động của máy biến áp}
\begin{bt}
% ID: L12C3B18-50-TL
% Mức: TH
Giải thích vì sao nhận định sau đúng: “Máy biến áp gồm lõi sắt và hai cuộn dây, hoạt động dựa trên cảm ứng điện từ với dòng điện xoay chiều.”
\loigiai{
Cần nêu được: Máy biến áp gồm lõi sắt và hai cuộn dây, hoạt động dựa trên cảm ứng điện từ với dòng điện xoay chiều. Khi vận dụng phải xác định đúng dữ kiện, phương chiều nếu có, điều kiện áp dụng, hệ đơn vị và kiểm tra tính hợp lí. Nhận định “Máy biến áp lí tưởng hoạt động được với nguồn điện một chiều không đổi.” sai.
}
\end{bt}

% ===== Câu 220 =====
% ID: L12C3B18-50-TLN
% Mức: TH
\begin{ex}
% ID: L12C3B18-50-TLN
% Mức: TH
Trong bốn phát biểu về Nhận biết cấu tạo, nguyên lí hoạt động của máy biến áp, có bao nhiêu phát biểu đúng? (Chỉ ghi một số nguyên.) (1) Máy biến áp lí tưởng hoạt động được với nguồn điện một chiều không đổi. (2) Máy biến áp gồm lõi sắt và hai cuộn dây, hoạt động dựa trên cảm ứng điện từ với dòng điện xoay chiều. (3) Kết quả phải phù hợp ý nghĩa vật lí. (4) Mọi đại lượng đều là vô hướng.
\shortans{1}
\loigiai{
Đối chiếu với kiến thức: Máy biến áp gồm lõi sắt và hai cuộn dây, hoạt động dựa trên cảm ứng điện từ với dòng điện xoay chiều. Có 1 phát biểu đúng.
}
\end{ex}

% ===== Câu 221 =====
% ID: L12C3B18-50-TN
% Mức: VD
\dangbt{Truyền tải điện năng đi xa và giảm hao phí trên đường dây}
\begin{ex}
% ID: L12C3B18-50-TN
% Mức: VD
Kết luận đúng khi vận dụng Truyền tải điện năng đi xa và giảm hao phí trên đường dây?
\choice
{\True Với công suất truyền không đổi, tăng điện áp truyền tải làm giảm dòng điện và giảm hao phí I^2R.}
{Muốn giảm hao phí phải giảm điện áp truyền tải.}
{Có thể bỏ qua phương, chiều và đơn vị trong mọi trường hợp.}
{Kết quả không phụ thuộc dữ kiện và điều kiện áp dụng.}
\loigiai{
Với công suất truyền không đổi, tăng điện áp truyền tải làm giảm dòng điện và giảm hao phí I^2R. Vì vậy phương án $A$ đúng; các phương án còn lại sai.
}
\end{ex}

% ===== Câu 222 =====
% ID: L12C3B18-51-DS
% Mức: TH
\begin{ex}
% ID: L12C3B18-51-DS
% Mức: TH
Xét Truyền tải điện năng đi xa và giảm hao phí trên đường dây (Tình huống 1). Hãy xác định tính đúng, sai của bốn phát biểu sau.
\choiceTF
{\True Với công suất truyền không đổi, tăng điện áp truyền tải làm giảm dòng điện và giảm hao phí I^2R.}
{Muốn giảm hao phí phải giảm điện áp truyền tải.}
{\True Khi giải cần xác định đúng dữ kiện, phương chiều, điều kiện áp dụng và đơn vị.}
{Kết quả không phụ thuộc dữ kiện và có thể bỏ qua điều kiện.}
\loigiai{
\begin{itemchoice}
\item (Đúng) Với công suất truyền không đổi, tăng điện áp truyền tải làm giảm dòng điện và giảm hao phí I^2R. (Vì): Với công suất truyền không đổi, tăng điện áp truyền tải làm giảm dòng điện và giảm hao phí I^2R. b) Sai: Muốn giảm hao phí phải giảm điện áp truyền tải. c) Đúng: Khi giải cần xác định đúng dữ kiện, phương chiều, điều kiện áp dụng và đơn vị. d) Sai: Kết quả không phụ thuộc dữ kiện và có thể bỏ qua điều kiện. Kiến thức cốt lõi: Với công suất truyền không đổi, tăng điện áp truyền tải làm giảm dòng điện và giảm hao phí I^2R
\item (Sai) Muốn giảm hao phí phải giảm điện áp truyền tải.
\item (Đúng) Khi giải cần xác định đúng dữ kiện, phương chiều, điều kiện áp dụng và đơn vị.
\item (Sai) Kết quả không phụ thuộc dữ kiện và có thể bỏ qua điều kiện.
\end{itemchoice}
}
\end{ex}

% ===== Câu 223 =====
% ID: L12C3B18-51-TL
% Mức: VD
\dangbt{Nhận biết cấu tạo, nguyên lí hoạt động của máy biến áp}
\begin{bt}
% ID: L12C3B18-51-TL
% Mức: VD
Phân tích sai lầm trong nhận định: “Máy biến áp lí tưởng hoạt động được với nguồn điện một chiều không đổi.”
\loigiai{
Cần nêu được: Máy biến áp gồm lõi sắt và hai cuộn dây, hoạt động dựa trên cảm ứng điện từ với dòng điện xoay chiều. Khi vận dụng phải xác định đúng dữ kiện, phương chiều nếu có, điều kiện áp dụng, hệ đơn vị và kiểm tra tính hợp lí. Nhận định “Máy biến áp lí tưởng hoạt động được với nguồn điện một chiều không đổi.” sai.
}
\end{bt}

% ===== Câu 224 =====
% ID: L12C3B18-51-TLN
% Mức: TH
\begin{ex}
% ID: L12C3B18-51-TLN
% Mức: TH
Trong bốn phát biểu về Nhận biết cấu tạo, nguyên lí hoạt động của máy biến áp, có bao nhiêu phát biểu đúng? (Chỉ ghi một số nguyên.) (1) Máy biến áp gồm lõi sắt và hai cuộn dây, hoạt động dựa trên cảm ứng điện từ với dòng điện xoay chiều. (2) Máy biến áp lí tưởng hoạt động được với nguồn điện một chiều không đổi. (3) Cần kiểm tra điều kiện áp dụng. (4) Có thể bỏ qua đơn vị.
\shortans{2}
\loigiai{
Đối chiếu với kiến thức: Máy biến áp gồm lõi sắt và hai cuộn dây, hoạt động dựa trên cảm ứng điện từ với dòng điện xoay chiều. Có 2 phát biểu đúng.
}
\end{ex}

% ===== Câu 225 =====
% ID: L12C3B18-51-TN
% Mức: NB
\dangbt{Tính điện áp, số vòng dây và cường độ dòng điện của máy biến áp lí tưởng}
\begin{ex}
% ID: L12C3B18-51-TN
% Mức: NB
Phát biểu nào sau đây đúng về Tính điện áp, số vòng dây và cường độ dòng điện của máy biến áp lí tưởng?
\choice
{Máy biến áp lí tưởng luôn có U2 = U1 bất kể số vòng dây.}
{Có thể bỏ qua phương, chiều và đơn vị trong mọi trường hợp.}
{Kết quả không phụ thuộc dữ kiện và điều kiện áp dụng.}
{\True Máy biến áp lí tưởng thỏa U2/U1 = N2/N1 và U1I1 = U2I2.}
\loigiai{
Máy biến áp lí tưởng thỏa U2/U1 = N2/N1 và U1I1 = U2I2. Vì vậy phương án $D$ đúng; các phương án còn lại sai.
}
\end{ex}

% ===== Câu 226 =====
% ID: L12C3B18-52-DS
% Mức: TH
\dangbt{Truyền tải điện năng đi xa và giảm hao phí trên đường dây}
\begin{ex}
% ID: L12C3B18-52-DS
% Mức: TH
Xét Truyền tải điện năng đi xa và giảm hao phí trên đường dây (Tình huống 2). Hãy xác định tính đúng, sai của bốn phát biểu sau.
\choiceTF
{Muốn giảm hao phí phải giảm điện áp truyền tải.}
{\True Với công suất truyền không đổi, tăng điện áp truyền tải làm giảm dòng điện và giảm hao phí I^2R.}
{Kết quả không phụ thuộc dữ kiện và có thể bỏ qua điều kiện.}
{\True Khi giải cần xác định đúng dữ kiện, phương chiều, điều kiện áp dụng và đơn vị.}
\loigiai{
\begin{itemchoice}
\item (Sai) Muốn giảm hao phí phải giảm điện áp truyền tải. (Vì): Muốn giảm hao phí phải giảm điện áp truyền tải. b) Đúng: Với công suất truyền không đổi, tăng điện áp truyền tải làm giảm dòng điện và giảm hao phí I^2R. c) Sai: Kết quả không phụ thuộc dữ kiện và có thể bỏ qua điều kiện. d) Đúng: Khi giải cần xác định đúng dữ kiện, phương chiều, điều kiện áp dụng và đơn vị. Kiến thức cốt lõi: Với công suất truyền không đổi, tăng điện áp truyền tải làm giảm dòng điện và giảm hao phí I^2R
\item (Đúng) Với công suất truyền không đổi, tăng điện áp truyền tải làm giảm dòng điện và giảm hao phí I^2R.
\item (Sai) Kết quả không phụ thuộc dữ kiện và có thể bỏ qua điều kiện.
\item (Đúng) Khi giải cần xác định đúng dữ kiện, phương chiều, điều kiện áp dụng và đơn vị.
\end{itemchoice}
}
\end{ex}

% ===== Câu 227 =====
% ID: L12C3B18-52-TL
% Mức: VD
\dangbt{Nhận biết cấu tạo, nguyên lí hoạt động của máy biến áp}
\begin{bt}
% ID: L12C3B18-52-TL
% Mức: VD
Đề xuất các bước giải một bài toán thuộc dạng Nhận biết cấu tạo, nguyên lí hoạt động của máy biến áp.
\loigiai{
Cần nêu được: Máy biến áp gồm lõi sắt và hai cuộn dây, hoạt động dựa trên cảm ứng điện từ với dòng điện xoay chiều. Khi vận dụng phải xác định đúng dữ kiện, phương chiều nếu có, điều kiện áp dụng, hệ đơn vị và kiểm tra tính hợp lí. Nhận định “Máy biến áp lí tưởng hoạt động được với nguồn điện một chiều không đổi.” sai.
}
\end{bt}

% ===== Câu 228 =====
% ID: L12C3B18-52-TLN
% Mức: VD
\begin{ex}
% ID: L12C3B18-52-TLN
% Mức: VD
Trong bốn phát biểu về Nhận biết cấu tạo, nguyên lí hoạt động của máy biến áp, có bao nhiêu phát biểu đúng? (Chỉ ghi một số nguyên.) (1) Máy biến áp lí tưởng hoạt động được với nguồn điện một chiều không đổi. (2) Máy biến áp gồm lõi sắt và hai cuộn dây, hoạt động dựa trên cảm ứng điện từ với dòng điện xoay chiều. (3) Kết quả phải phù hợp ý nghĩa vật lí. (4) Mọi đại lượng đều là vô hướng.
\shortans{2}
\loigiai{
Đối chiếu với kiến thức: Máy biến áp gồm lõi sắt và hai cuộn dây, hoạt động dựa trên cảm ứng điện từ với dòng điện xoay chiều. Có 2 phát biểu đúng.
}
\end{ex}

% ===== Câu 229 =====
% ID: L12C3B18-52-TN
% Mức: NB
\dangbt{Tính điện áp, số vòng dây và cường độ dòng điện của máy biến áp lí tưởng}
\begin{ex}
% ID: L12C3B18-52-TN
% Mức: NB
Trong nội dung Tính điện áp, số vòng dây và cường độ dòng điện của máy biến áp lí tưởng?
\choice
{\True Máy biến áp lí tưởng thỏa U2/U1 = N2/N1 và U1I1 = U2I2.}
{Máy biến áp lí tưởng luôn có U2 = U1 bất kể số vòng dây.}
{Có thể bỏ qua phương, chiều và đơn vị trong mọi trường hợp.}
{Kết quả không phụ thuộc dữ kiện và điều kiện áp dụng.}
\loigiai{
Máy biến áp lí tưởng thỏa U2/U1 = N2/N1 và U1I1 = U2I2. Vì vậy phương án $A$ đúng; các phương án còn lại sai.
}
\end{ex}

% ===== Câu 230 =====
% ID: L12C3B18-53-DS
% Mức: VD
\dangbt{Truyền tải điện năng đi xa và giảm hao phí trên đường dây}
\begin{ex}
% ID: L12C3B18-53-DS
% Mức: VD
Xét Truyền tải điện năng đi xa và giảm hao phí trên đường dây (Tình huống 3). Hãy xác định tính đúng, sai của bốn phát biểu sau.
\choiceTF
{\True Khi giải cần xác định đúng dữ kiện, phương chiều, điều kiện áp dụng và đơn vị.}
{Kết quả không phụ thuộc dữ kiện và có thể bỏ qua điều kiện.}
{\True Với công suất truyền không đổi, tăng điện áp truyền tải làm giảm dòng điện và giảm hao phí I^2R.}
{Muốn giảm hao phí phải giảm điện áp truyền tải.}
\loigiai{
\begin{itemchoice}
\item (Đúng) Khi giải cần xác định đúng dữ kiện, phương chiều, điều kiện áp dụng và đơn vị. (Vì): Khi giải cần xác định đúng dữ kiện, phương chiều, điều kiện áp dụng và đơn vị. b) Sai: Kết quả không phụ thuộc dữ kiện và có thể bỏ qua điều kiện. c) Đúng: Với công suất truyền không đổi, tăng điện áp truyền tải làm giảm dòng điện và giảm hao phí I^2R. d) Sai: Muốn giảm hao phí phải giảm điện áp truyền tải. Kiến thức cốt lõi: Với công suất truyền không đổi, tăng điện áp truyền tải làm giảm dòng điện và giảm hao phí I^2R
\item (Sai) Kết quả không phụ thuộc dữ kiện và có thể bỏ qua điều kiện.
\item (Đúng) Với công suất truyền không đổi, tăng điện áp truyền tải làm giảm dòng điện và giảm hao phí I^2R.
\item (Sai) Muốn giảm hao phí phải giảm điện áp truyền tải.
\end{itemchoice}
}
\end{ex}

% ===== Câu 231 =====
% ID: L12C3B18-53-TL
% Mức: VD
\dangbt{Nhận biết cấu tạo, nguyên lí hoạt động của máy biến áp}
\begin{bt}
% ID: L12C3B18-53-TL
% Mức: VD
Nêu một tình huống thực tiễn liên quan đến Nhận biết cấu tạo, nguyên lí hoạt động của máy biến áp và giải thích bằng kiến thức Vật lí.
\loigiai{
Cần nêu được: Máy biến áp gồm lõi sắt và hai cuộn dây, hoạt động dựa trên cảm ứng điện từ với dòng điện xoay chiều. Khi vận dụng phải xác định đúng dữ kiện, phương chiều nếu có, điều kiện áp dụng, hệ đơn vị và kiểm tra tính hợp lí. Nhận định “Máy biến áp lí tưởng hoạt động được với nguồn điện một chiều không đổi.” sai.
}
\end{bt}

% ===== Câu 232 =====
% ID: L12C3B18-53-TLN
% Mức: VD
\begin{ex}
% ID: L12C3B18-53-TLN
% Mức: VD
Trong bốn phát biểu về Nhận biết cấu tạo, nguyên lí hoạt động của máy biến áp, có bao nhiêu phát biểu đúng? (Chỉ ghi một số nguyên.) (1) Máy biến áp gồm lõi sắt và hai cuộn dây, hoạt động dựa trên cảm ứng điện từ với dòng điện xoay chiều. (2) Máy biến áp lí tưởng hoạt động được với nguồn điện một chiều không đổi. (3) Cần kiểm tra điều kiện áp dụng. (4) Có thể bỏ qua đơn vị.
\shortans{2}
\loigiai{
Đối chiếu với kiến thức: Máy biến áp gồm lõi sắt và hai cuộn dây, hoạt động dựa trên cảm ứng điện từ với dòng điện xoay chiều. Có 2 phát biểu đúng.
}
\end{ex}

% ===== Câu 233 =====
% ID: L12C3B18-53-TN
% Mức: NB
\dangbt{Tính điện áp, số vòng dây và cường độ dòng điện của máy biến áp lí tưởng}
\begin{ex}
% ID: L12C3B18-53-TN
% Mức: NB
Tình huống 3: chọn kết luận chính xác về Tính điện áp, số vòng dây và cường độ dòng điện của máy biến áp lí tưởng?
\choice
{Kết quả không phụ thuộc dữ kiện và điều kiện áp dụng.}
{\True Máy biến áp lí tưởng thỏa U2/U1 = N2/N1 và U1I1 = U2I2.}
{Máy biến áp lí tưởng luôn có U2 = U1 bất kể số vòng dây.}
{Có thể bỏ qua phương, chiều và đơn vị trong mọi trường hợp.}
\loigiai{
Máy biến áp lí tưởng thỏa U2/U1 = N2/N1 và U1I1 = U2I2. Vì vậy phương án $B$ đúng; các phương án còn lại sai.
}
\end{ex}

% ===== Câu 234 =====
% ID: L12C3B18-54-DS
% Mức: VD
\dangbt{Truyền tải điện năng đi xa và giảm hao phí trên đường dây}
\begin{ex}
% ID: L12C3B18-54-DS
% Mức: VD
Xét Truyền tải điện năng đi xa và giảm hao phí trên đường dây (Tình huống 4). Hãy xác định tính đúng, sai của bốn phát biểu sau.
\choiceTF
{Kết quả không phụ thuộc dữ kiện và có thể bỏ qua điều kiện.}
{\True Khi giải cần xác định đúng dữ kiện, phương chiều, điều kiện áp dụng và đơn vị.}
{Muốn giảm hao phí phải giảm điện áp truyền tải.}
{\True Với công suất truyền không đổi, tăng điện áp truyền tải làm giảm dòng điện và giảm hao phí I^2R.}
\loigiai{
\begin{itemchoice}
\item (Sai) Kết quả không phụ thuộc dữ kiện và có thể bỏ qua điều kiện. (Vì): Kết quả không phụ thuộc dữ kiện và có thể bỏ qua điều kiện. b) Đúng: Khi giải cần xác định đúng dữ kiện, phương chiều, điều kiện áp dụng và đơn vị. c) Sai: Muốn giảm hao phí phải giảm điện áp truyền tải. d) Đúng: Với công suất truyền không đổi, tăng điện áp truyền tải làm giảm dòng điện và giảm hao phí I^2R. Kiến thức cốt lõi: Với công suất truyền không đổi, tăng điện áp truyền tải làm giảm dòng điện và giảm hao phí I^2R
\item (Đúng) Khi giải cần xác định đúng dữ kiện, phương chiều, điều kiện áp dụng và đơn vị.
\item (Sai) Muốn giảm hao phí phải giảm điện áp truyền tải.
\item (Đúng) Với công suất truyền không đổi, tăng điện áp truyền tải làm giảm dòng điện và giảm hao phí I^2R.
\end{itemchoice}
}
\end{ex}

% ===== Câu 235 =====
% ID: L12C3B18-54-TL
% Mức: VDC
\dangbt{Nhận biết cấu tạo, nguyên lí hoạt động của máy biến áp}
\begin{bt}
% ID: L12C3B18-54-TL
% Mức: VDC
So sánh nhận định đúng “Máy biến áp gồm lõi sắt và hai cuộn dây, hoạt động dựa trên cảm ứng điện từ với dòng điện xoay chiều.” với nhận định sai “Máy biến áp lí tưởng hoạt động được với nguồn điện một chiều không đổi.”; chỉ rõ căn cứ Vật lí.
\loigiai{
Cần nêu được: Máy biến áp gồm lõi sắt và hai cuộn dây, hoạt động dựa trên cảm ứng điện từ với dòng điện xoay chiều. Khi vận dụng phải xác định đúng dữ kiện, phương chiều nếu có, điều kiện áp dụng, hệ đơn vị và kiểm tra tính hợp lí. Nhận định “Máy biến áp lí tưởng hoạt động được với nguồn điện một chiều không đổi.” sai.
}
\end{bt}

% ===== Câu 236 =====
% ID: L12C3B18-54-TLN
% Mức: VDC
\begin{ex}
% ID: L12C3B18-54-TLN
% Mức: VDC
Trong bốn phát biểu về Nhận biết cấu tạo, nguyên lí hoạt động của máy biến áp, có bao nhiêu phát biểu đúng? (Chỉ ghi một số nguyên.) (1) Khi giải cần xác định đúng dữ kiện, phương chiều, điều kiện áp dụng và đơn vị. (2) Máy biến áp lí tưởng hoạt động được với nguồn điện một chiều không đổi. (3) Phải dùng đúng định luật cho tình huống. (4) Máy biến áp gồm lõi sắt và hai cuộn dây, hoạt động dựa trên cảm ứng điện từ với dòng điện xoay chiều.
\shortans{3}
\loigiai{
Đối chiếu với kiến thức: Máy biến áp gồm lõi sắt và hai cuộn dây, hoạt động dựa trên cảm ứng điện từ với dòng điện xoay chiều. Có 3 phát biểu đúng.
}
\end{ex}

% ===== Câu 237 =====
% ID: L12C3B18-54-TN
% Mức: TH
\dangbt{Tính điện áp, số vòng dây và cường độ dòng điện của máy biến áp lí tưởng}
\begin{ex}
% ID: L12C3B18-54-TN
% Mức: TH
Nội dung nào mô tả đúng nhất Tính điện áp, số vòng dây và cường độ dòng điện của máy biến áp lí tưởng?
\choice
{Có thể bỏ qua phương, chiều và đơn vị trong mọi trường hợp.}
{Kết quả không phụ thuộc dữ kiện và điều kiện áp dụng.}
{\True Máy biến áp lí tưởng thỏa U2/U1 = N2/N1 và U1I1 = U2I2.}
{Máy biến áp lí tưởng luôn có U2 = U1 bất kể số vòng dây.}
\loigiai{
Máy biến áp lí tưởng thỏa U2/U1 = N2/N1 và U1I1 = U2I2. Vì vậy phương án $C$ đúng; các phương án còn lại sai.
}
\end{ex}

% ===== Câu 238 =====
% ID: L12C3B18-55-DS
% Mức: VD
\dangbt{Truyền tải điện năng đi xa và giảm hao phí trên đường dây}
\begin{ex}
% ID: L12C3B18-55-DS
% Mức: VD
Xét Truyền tải điện năng đi xa và giảm hao phí trên đường dây (Tình huống 5). Hãy xác định tính đúng, sai của bốn phát biểu sau.
\choiceTF
{\True Với công suất truyền không đổi, tăng điện áp truyền tải làm giảm dòng điện và giảm hao phí I^2R.}
{Kết quả không phụ thuộc dữ kiện và có thể bỏ qua điều kiện.}
{Muốn giảm hao phí phải giảm điện áp truyền tải.}
{\True Khi giải cần xác định đúng dữ kiện, phương chiều, điều kiện áp dụng và đơn vị.}
\loigiai{
\begin{itemchoice}
\item (Đúng) Với công suất truyền không đổi, tăng điện áp truyền tải làm giảm dòng điện và giảm hao phí I^2R. (Vì): Với công suất truyền không đổi, tăng điện áp truyền tải làm giảm dòng điện và giảm hao phí I^2R. b) Sai: Kết quả không phụ thuộc dữ kiện và có thể bỏ qua điều kiện. c) Sai: Muốn giảm hao phí phải giảm điện áp truyền tải. d) Đúng: Khi giải cần xác định đúng dữ kiện, phương chiều, điều kiện áp dụng và đơn vị. Kiến thức cốt lõi: Với công suất truyền không đổi, tăng điện áp truyền tải làm giảm dòng điện và giảm hao phí I^2R
\item (Sai) Kết quả không phụ thuộc dữ kiện và có thể bỏ qua điều kiện.
\item (Sai) Muốn giảm hao phí phải giảm điện áp truyền tải.
\item (Đúng) Khi giải cần xác định đúng dữ kiện, phương chiều, điều kiện áp dụng và đơn vị.
\end{itemchoice}
}
\end{ex}

% ===== Câu 239 =====
% ID: L12C3B18-55-TL
% Mức: TH
\dangbt{Tính điện áp, số vòng dây và cường độ dòng điện của máy biến áp lí tưởng}
\begin{bt}
% ID: L12C3B18-55-TL
% Mức: TH
Trình bày kiến thức cốt lõi của Tính điện áp, số vòng dây và cường độ dòng điện của máy biến áp lí tưởng và nêu điều kiện áp dụng.
\loigiai{
Cần nêu được: Máy biến áp lí tưởng thỏa U2/U1 = N2/N1 và U1I1 = U2I2. Khi vận dụng phải xác định đúng dữ kiện, phương chiều nếu có, điều kiện áp dụng, hệ đơn vị và kiểm tra tính hợp lí. Nhận định “Máy biến áp lí tưởng luôn có U2 = U1 bất kể số vòng dây.” sai.
}
\end{bt}

% ===== Câu 240 =====
% ID: L12C3B18-55-TLN
% Mức: TH
\begin{ex}
% ID: L12C3B18-55-TLN
% Mức: TH
Trong bốn phát biểu về Tính điện áp, số vòng dây và cường độ dòng điện của máy biến áp lí tưởng, có bao nhiêu phát biểu đúng? (Chỉ ghi một số nguyên.) (1) Máy biến áp lí tưởng thỏa U2/U1 = N2/N1 và U1I1 = U2I2. (2) Máy biến áp lí tưởng luôn có U2 = U1 bất kể số vòng dây. (3) Cần kiểm tra điều kiện áp dụng. (4) Có thể bỏ qua đơn vị.
\shortans{2}
\loigiai{
Đối chiếu với kiến thức: Máy biến áp lí tưởng thỏa U2/U1 = N2/N1 và U1I1 = U2I2. Có 2 phát biểu đúng.
}
\end{ex}

% ===== Câu 241 =====
% ID: L12C3B18-55-TN
% Mức: TH
\begin{ex}
% ID: L12C3B18-55-TN
% Mức: TH
Khi phân tích Tính điện áp, số vòng dây và cường độ dòng điện của máy biến áp lí tưởng?
\choice
{Máy biến áp lí tưởng luôn có U2 = U1 bất kể số vòng dây.}
{Có thể bỏ qua phương, chiều và đơn vị trong mọi trường hợp.}
{Kết quả không phụ thuộc dữ kiện và điều kiện áp dụng.}
{\True Máy biến áp lí tưởng thỏa U2/U1 = N2/N1 và U1I1 = U2I2.}
\loigiai{
Máy biến áp lí tưởng thỏa U2/U1 = N2/N1 và U1I1 = U2I2. Vì vậy phương án $D$ đúng; các phương án còn lại sai.
}
\end{ex}

% ===== Câu 242 =====
% ID: L12C3B18-56-DS
% Mức: TH
\dangbt{Bài toán tổng hợp về máy biến áp, truyền tải điện và an toàn điện}
\begin{ex}
% ID: L12C3B18-56-DS
% Mức: TH
Xét Bài toán tổng hợp về máy biến áp, truyền tải điện và an toàn điện (Tình huống 1). Hãy xác định tính đúng, sai của bốn phát biểu sau.
\choiceTF
{\True An toàn điện yêu cầu cách điện, nối đất phù hợp và không chạm phần dẫn điện đang có điện áp nguy hiểm.}
{Có thể bỏ qua nối đất và thiết bị bảo vệ khi dùng điện áp cao.}
{\True Khi giải cần xác định đúng dữ kiện, phương chiều, điều kiện áp dụng và đơn vị.}
{Kết quả không phụ thuộc dữ kiện và có thể bỏ qua điều kiện.}
\loigiai{
\begin{itemchoice}
\item (Đúng) An toàn điện yêu cầu cách điện, nối đất phù hợp và không chạm phần dẫn điện đang có điện áp nguy hiểm. (Vì): An toàn điện yêu cầu cách điện, nối đất phù hợp và không chạm phần dẫn điện đang có điện áp nguy hiểm. b) Sai: Có thể bỏ qua nối đất và thiết bị bảo vệ khi dùng điện áp cao. c) Đúng: Khi giải cần xác định đúng dữ kiện, phương chiều, điều kiện áp dụng và đơn vị. d) Sai: Kết quả không phụ thuộc dữ kiện và có thể bỏ qua điều kiện. Kiến thức cốt lõi: An toàn điện yêu cầu cách điện, nối đất phù hợp và không chạm phần dẫn điện đang có điện áp nguy hiểm
\item (Sai) Có thể bỏ qua nối đất và thiết bị bảo vệ khi dùng điện áp cao.
\item (Đúng) Khi giải cần xác định đúng dữ kiện, phương chiều, điều kiện áp dụng và đơn vị.
\item (Sai) Kết quả không phụ thuộc dữ kiện và có thể bỏ qua điều kiện.
\end{itemchoice}
}
\end{ex}

% ===== Câu 243 =====
% ID: L12C3B18-56-TL
% Mức: TH
\dangbt{Tính điện áp, số vòng dây và cường độ dòng điện của máy biến áp lí tưởng}
\begin{bt}
% ID: L12C3B18-56-TL
% Mức: TH
Giải thích vì sao nhận định sau đúng: “Máy biến áp lí tưởng thỏa U2/U1 = N2/N1 và U1I1 = U2I2.”
\loigiai{
Cần nêu được: Máy biến áp lí tưởng thỏa U2/U1 = N2/N1 và U1I1 = U2I2. Khi vận dụng phải xác định đúng dữ kiện, phương chiều nếu có, điều kiện áp dụng, hệ đơn vị và kiểm tra tính hợp lí. Nhận định “Máy biến áp lí tưởng luôn có U2 = U1 bất kể số vòng dây.” sai.
}
\end{bt}

% ===== Câu 244 =====
% ID: L12C3B18-56-TLN
% Mức: TH
\begin{ex}
% ID: L12C3B18-56-TLN
% Mức: TH
Trong bốn phát biểu về Tính điện áp, số vòng dây và cường độ dòng điện của máy biến áp lí tưởng, có bao nhiêu phát biểu đúng? (Chỉ ghi một số nguyên.) (1) Máy biến áp lí tưởng luôn có U2 = U1 bất kể số vòng dây. (2) Máy biến áp lí tưởng thỏa U2/U1 = N2/N1 và U1I1 = U2I2. (3) Kết quả phải phù hợp ý nghĩa vật lí. (4) Mọi đại lượng đều là vô hướng.
\shortans{1}
\loigiai{
Đối chiếu với kiến thức: Máy biến áp lí tưởng thỏa U2/U1 = N2/N1 và U1I1 = U2I2. Có 1 phát biểu đúng.
}
\end{ex}

% ===== Câu 245 =====
% ID: L12C3B18-56-TN
% Mức: TH
\begin{ex}
% ID: L12C3B18-56-TN
% Mức: TH
Một học sinh ôn tập Tính điện áp, số vòng dây và cường độ dòng điện của máy biến áp lí tưởng?
\choice
{\True Máy biến áp lí tưởng thỏa U2/U1 = N2/N1 và U1I1 = U2I2.}
{Máy biến áp lí tưởng luôn có U2 = U1 bất kể số vòng dây.}
{Có thể bỏ qua phương, chiều và đơn vị trong mọi trường hợp.}
{Kết quả không phụ thuộc dữ kiện và điều kiện áp dụng.}
\loigiai{
Máy biến áp lí tưởng thỏa U2/U1 = N2/N1 và U1I1 = U2I2. Vì vậy phương án $A$ đúng; các phương án còn lại sai.
}
\end{ex}

% ===== Câu 246 =====
% ID: L12C3B18-57-DS
% Mức: TH
\dangbt{Bài toán tổng hợp về máy biến áp, truyền tải điện và an toàn điện}
\begin{ex}
% ID: L12C3B18-57-DS
% Mức: TH
Xét Bài toán tổng hợp về máy biến áp, truyền tải điện và an toàn điện (Tình huống 2). Hãy xác định tính đúng, sai của bốn phát biểu sau.
\choiceTF
{Có thể bỏ qua nối đất và thiết bị bảo vệ khi dùng điện áp cao.}
{\True An toàn điện yêu cầu cách điện, nối đất phù hợp và không chạm phần dẫn điện đang có điện áp nguy hiểm.}
{Kết quả không phụ thuộc dữ kiện và có thể bỏ qua điều kiện.}
{\True Khi giải cần xác định đúng dữ kiện, phương chiều, điều kiện áp dụng và đơn vị.}
\loigiai{
\begin{itemchoice}
\item (Sai) Có thể bỏ qua nối đất và thiết bị bảo vệ khi dùng điện áp cao. (Vì): Có thể bỏ qua nối đất và thiết bị bảo vệ khi dùng điện áp cao. b) Đúng: An toàn điện yêu cầu cách điện, nối đất phù hợp và không chạm phần dẫn điện đang có điện áp nguy hiểm. c) Sai: Kết quả không phụ thuộc dữ kiện và có thể bỏ qua điều kiện. d) Đúng: Khi giải cần xác định đúng dữ kiện, phương chiều, điều kiện áp dụng và đơn vị. Kiến thức cốt lõi: An toàn điện yêu cầu cách điện, nối đất phù hợp và không chạm phần dẫn điện đang có điện áp nguy hiểm
\item (Đúng) An toàn điện yêu cầu cách điện, nối đất phù hợp và không chạm phần dẫn điện đang có điện áp nguy hiểm.
\item (Sai) Kết quả không phụ thuộc dữ kiện và có thể bỏ qua điều kiện.
\item (Đúng) Khi giải cần xác định đúng dữ kiện, phương chiều, điều kiện áp dụng và đơn vị.
\end{itemchoice}
}
\end{ex}

% ===== Câu 247 =====
% ID: L12C3B18-57-TL
% Mức: VD
\dangbt{Tính điện áp, số vòng dây và cường độ dòng điện của máy biến áp lí tưởng}
\begin{bt}
% ID: L12C3B18-57-TL
% Mức: VD
Phân tích sai lầm trong nhận định: “Máy biến áp lí tưởng luôn có U2 = U1 bất kể số vòng dây.”
\loigiai{
Cần nêu được: Máy biến áp lí tưởng thỏa U2/U1 = N2/N1 và U1I1 = U2I2. Khi vận dụng phải xác định đúng dữ kiện, phương chiều nếu có, điều kiện áp dụng, hệ đơn vị và kiểm tra tính hợp lí. Nhận định “Máy biến áp lí tưởng luôn có U2 = U1 bất kể số vòng dây.” sai.
}
\end{bt}

% ===== Câu 248 =====
% ID: L12C3B18-57-TLN
% Mức: TH
\begin{ex}
% ID: L12C3B18-57-TLN
% Mức: TH
Trong bốn phát biểu về Tính điện áp, số vòng dây và cường độ dòng điện của máy biến áp lí tưởng, có bao nhiêu phát biểu đúng? (Chỉ ghi một số nguyên.) (1) Máy biến áp lí tưởng thỏa U2/U1 = N2/N1 và U1I1 = U2I2. (2) Máy biến áp lí tưởng luôn có U2 = U1 bất kể số vòng dây. (3) Cần kiểm tra điều kiện áp dụng. (4) Có thể bỏ qua đơn vị.
\shortans{2}
\loigiai{
Đối chiếu với kiến thức: Máy biến áp lí tưởng thỏa U2/U1 = N2/N1 và U1I1 = U2I2. Có 2 phát biểu đúng.
}
\end{ex}

% ===== Câu 249 =====
% ID: L12C3B18-57-TN
% Mức: TH
\begin{ex}
% ID: L12C3B18-57-TN
% Mức: TH
Chọn phát biểu không trái với kiến thức về Tính điện áp, số vòng dây và cường độ dòng điện của máy biến áp lí tưởng?
\choice
{Kết quả không phụ thuộc dữ kiện và điều kiện áp dụng.}
{\True Máy biến áp lí tưởng thỏa U2/U1 = N2/N1 và U1I1 = U2I2.}
{Máy biến áp lí tưởng luôn có U2 = U1 bất kể số vòng dây.}
{Có thể bỏ qua phương, chiều và đơn vị trong mọi trường hợp.}
\loigiai{
Máy biến áp lí tưởng thỏa U2/U1 = N2/N1 và U1I1 = U2I2. Vì vậy phương án $B$ đúng; các phương án còn lại sai.
}
\end{ex}

% ===== Câu 250 =====
% ID: L12C3B18-58-DS
% Mức: VD
\dangbt{Bài toán tổng hợp về máy biến áp, truyền tải điện và an toàn điện}
\begin{ex}
% ID: L12C3B18-58-DS
% Mức: VD
Xét Bài toán tổng hợp về máy biến áp, truyền tải điện và an toàn điện (Tình huống 3). Hãy xác định tính đúng, sai của bốn phát biểu sau.
\choiceTF
{\True Khi giải cần xác định đúng dữ kiện, phương chiều, điều kiện áp dụng và đơn vị.}
{Kết quả không phụ thuộc dữ kiện và có thể bỏ qua điều kiện.}
{\True An toàn điện yêu cầu cách điện, nối đất phù hợp và không chạm phần dẫn điện đang có điện áp nguy hiểm.}
{Có thể bỏ qua nối đất và thiết bị bảo vệ khi dùng điện áp cao.}
\loigiai{
\begin{itemchoice}
\item (Đúng) Khi giải cần xác định đúng dữ kiện, phương chiều, điều kiện áp dụng và đơn vị. (Vì): Khi giải cần xác định đúng dữ kiện, phương chiều, điều kiện áp dụng và đơn vị. b) Sai: Kết quả không phụ thuộc dữ kiện và có thể bỏ qua điều kiện. c) Đúng: An toàn điện yêu cầu cách điện, nối đất phù hợp và không chạm phần dẫn điện đang có điện áp nguy hiểm. d) Sai: Có thể bỏ qua nối đất và thiết bị bảo vệ khi dùng điện áp cao. Kiến thức cốt lõi: An toàn điện yêu cầu cách điện, nối đất phù hợp và không chạm phần dẫn điện đang có điện áp nguy hiểm
\item (Sai) Kết quả không phụ thuộc dữ kiện và có thể bỏ qua điều kiện.
\item (Đúng) An toàn điện yêu cầu cách điện, nối đất phù hợp và không chạm phần dẫn điện đang có điện áp nguy hiểm.
\item (Sai) Có thể bỏ qua nối đất và thiết bị bảo vệ khi dùng điện áp cao.
\end{itemchoice}
}
\end{ex}

% ===== Câu 251 =====
% ID: L12C3B18-58-TL
% Mức: VD
\dangbt{Tính điện áp, số vòng dây và cường độ dòng điện của máy biến áp lí tưởng}
\begin{bt}
% ID: L12C3B18-58-TL
% Mức: VD
Đề xuất các bước giải một bài toán thuộc dạng Tính điện áp, số vòng dây và cường độ dòng điện của máy biến áp lí tưởng.
\loigiai{
Cần nêu được: Máy biến áp lí tưởng thỏa U2/U1 = N2/N1 và U1I1 = U2I2. Khi vận dụng phải xác định đúng dữ kiện, phương chiều nếu có, điều kiện áp dụng, hệ đơn vị và kiểm tra tính hợp lí. Nhận định “Máy biến áp lí tưởng luôn có U2 = U1 bất kể số vòng dây.” sai.
}
\end{bt}

% ===== Câu 252 =====
% ID: L12C3B18-58-TLN
% Mức: VD
\begin{ex}
% ID: L12C3B18-58-TLN
% Mức: VD
Trong bốn phát biểu về Tính điện áp, số vòng dây và cường độ dòng điện của máy biến áp lí tưởng, có bao nhiêu phát biểu đúng? (Chỉ ghi một số nguyên.) (1) Máy biến áp lí tưởng luôn có U2 = U1 bất kể số vòng dây. (2) Máy biến áp lí tưởng thỏa U2/U1 = N2/N1 và U1I1 = U2I2. (3) Kết quả phải phù hợp ý nghĩa vật lí. (4) Mọi đại lượng đều là vô hướng.
\shortans{2}
\loigiai{
Đối chiếu với kiến thức: Máy biến áp lí tưởng thỏa U2/U1 = N2/N1 và U1I1 = U2I2. Có 2 phát biểu đúng.
}
\end{ex}

% ===== Câu 253 =====
% ID: L12C3B18-58-TN
% Mức: VD
\begin{ex}
% ID: L12C3B18-58-TN
% Mức: VD
Cơ sở Vật lí đúng dùng để giải Tính điện áp, số vòng dây và cường độ dòng điện của máy biến áp lí tưởng?
\choice
{Có thể bỏ qua phương, chiều và đơn vị trong mọi trường hợp.}
{Kết quả không phụ thuộc dữ kiện và điều kiện áp dụng.}
{\True Máy biến áp lí tưởng thỏa U2/U1 = N2/N1 và U1I1 = U2I2.}
{Máy biến áp lí tưởng luôn có U2 = U1 bất kể số vòng dây.}
\loigiai{
Máy biến áp lí tưởng thỏa U2/U1 = N2/N1 và U1I1 = U2I2. Vì vậy phương án $C$ đúng; các phương án còn lại sai.
}
\end{ex}

% ===== Câu 254 =====
% ID: L12C3B18-59-DS
% Mức: VD
\dangbt{Bài toán tổng hợp về máy biến áp, truyền tải điện và an toàn điện}
\begin{ex}
% ID: L12C3B18-59-DS
% Mức: VD
Xét Bài toán tổng hợp về máy biến áp, truyền tải điện và an toàn điện (Tình huống 4). Hãy xác định tính đúng, sai của bốn phát biểu sau.
\choiceTF
{Kết quả không phụ thuộc dữ kiện và có thể bỏ qua điều kiện.}
{\True Khi giải cần xác định đúng dữ kiện, phương chiều, điều kiện áp dụng và đơn vị.}
{Có thể bỏ qua nối đất và thiết bị bảo vệ khi dùng điện áp cao.}
{\True An toàn điện yêu cầu cách điện, nối đất phù hợp và không chạm phần dẫn điện đang có điện áp nguy hiểm.}
\loigiai{
\begin{itemchoice}
\item (Sai) Kết quả không phụ thuộc dữ kiện và có thể bỏ qua điều kiện. (Vì): Kết quả không phụ thuộc dữ kiện và có thể bỏ qua điều kiện. b) Đúng: Khi giải cần xác định đúng dữ kiện, phương chiều, điều kiện áp dụng và đơn vị. c) Sai: Có thể bỏ qua nối đất và thiết bị bảo vệ khi dùng điện áp cao. d) Đúng: An toàn điện yêu cầu cách điện, nối đất phù hợp và không chạm phần dẫn điện đang có điện áp nguy hiểm. Kiến thức cốt lõi: An toàn điện yêu cầu cách điện, nối đất phù hợp và không chạm phần dẫn điện đang có điện áp nguy hiểm
\item (Đúng) Khi giải cần xác định đúng dữ kiện, phương chiều, điều kiện áp dụng và đơn vị.
\item (Sai) Có thể bỏ qua nối đất và thiết bị bảo vệ khi dùng điện áp cao.
\item (Đúng) An toàn điện yêu cầu cách điện, nối đất phù hợp và không chạm phần dẫn điện đang có điện áp nguy hiểm.
\end{itemchoice}
}
\end{ex}

% ===== Câu 255 =====
% ID: L12C3B18-59-TL
% Mức: VD
\dangbt{Tính điện áp, số vòng dây và cường độ dòng điện của máy biến áp lí tưởng}
\begin{bt}
% ID: L12C3B18-59-TL
% Mức: VD
Nêu một tình huống thực tiễn liên quan đến Tính điện áp, số vòng dây và cường độ dòng điện của máy biến áp lí tưởng và giải thích bằng kiến thức Vật lí.
\loigiai{
Cần nêu được: Máy biến áp lí tưởng thỏa U2/U1 = N2/N1 và U1I1 = U2I2. Khi vận dụng phải xác định đúng dữ kiện, phương chiều nếu có, điều kiện áp dụng, hệ đơn vị và kiểm tra tính hợp lí. Nhận định “Máy biến áp lí tưởng luôn có U2 = U1 bất kể số vòng dây.” sai.
}
\end{bt}

% ===== Câu 256 =====
% ID: L12C3B18-59-TLN
% Mức: VD
\begin{ex}
% ID: L12C3B18-59-TLN
% Mức: VD
Trong bốn phát biểu về Tính điện áp, số vòng dây và cường độ dòng điện của máy biến áp lí tưởng, có bao nhiêu phát biểu đúng? (Chỉ ghi một số nguyên.) (1) Máy biến áp lí tưởng thỏa U2/U1 = N2/N1 và U1I1 = U2I2. (2) Máy biến áp lí tưởng luôn có U2 = U1 bất kể số vòng dây. (3) Cần kiểm tra điều kiện áp dụng. (4) Có thể bỏ qua đơn vị.
\shortans{2}
\loigiai{
Đối chiếu với kiến thức: Máy biến áp lí tưởng thỏa U2/U1 = N2/N1 và U1I1 = U2I2. Có 2 phát biểu đúng.
}
\end{ex}

% ===== Câu 257 =====
% ID: L12C3B18-59-TN
% Mức: VD
\begin{ex}
% ID: L12C3B18-59-TN
% Mức: VD
Trong một bài thực hành về Tính điện áp, số vòng dây và cường độ dòng điện của máy biến áp lí tưởng?
\choice
{Máy biến áp lí tưởng luôn có U2 = U1 bất kể số vòng dây.}
{Có thể bỏ qua phương, chiều và đơn vị trong mọi trường hợp.}
{Kết quả không phụ thuộc dữ kiện và điều kiện áp dụng.}
{\True Máy biến áp lí tưởng thỏa U2/U1 = N2/N1 và U1I1 = U2I2.}
\loigiai{
Máy biến áp lí tưởng thỏa U2/U1 = N2/N1 và U1I1 = U2I2. Vì vậy phương án $D$ đúng; các phương án còn lại sai.
}
\end{ex}

% ===== Câu 258 =====
% ID: L12C3B18-60-DS
% Mức: VD
\dangbt{Bài toán tổng hợp về máy biến áp, truyền tải điện và an toàn điện}
\begin{ex}
% ID: L12C3B18-60-DS
% Mức: VD
Xét Bài toán tổng hợp về máy biến áp, truyền tải điện và an toàn điện (Tình huống 5). Hãy xác định tính đúng, sai của bốn phát biểu sau.
\choiceTF
{\True An toàn điện yêu cầu cách điện, nối đất phù hợp và không chạm phần dẫn điện đang có điện áp nguy hiểm.}
{Kết quả không phụ thuộc dữ kiện và có thể bỏ qua điều kiện.}
{Có thể bỏ qua nối đất và thiết bị bảo vệ khi dùng điện áp cao.}
{\True Khi giải cần xác định đúng dữ kiện, phương chiều, điều kiện áp dụng và đơn vị.}
\loigiai{
\begin{itemchoice}
\item (Đúng) An toàn điện yêu cầu cách điện, nối đất phù hợp và không chạm phần dẫn điện đang có điện áp nguy hiểm. (Vì): An toàn điện yêu cầu cách điện, nối đất phù hợp và không chạm phần dẫn điện đang có điện áp nguy hiểm. b) Sai: Kết quả không phụ thuộc dữ kiện và có thể bỏ qua điều kiện. c) Sai: Có thể bỏ qua nối đất và thiết bị bảo vệ khi dùng điện áp cao. d) Đúng: Khi giải cần xác định đúng dữ kiện, phương chiều, điều kiện áp dụng và đơn vị. Kiến thức cốt lõi: An toàn điện yêu cầu cách điện, nối đất phù hợp và không chạm phần dẫn điện đang có điện áp nguy hiểm
\item (Sai) Kết quả không phụ thuộc dữ kiện và có thể bỏ qua điều kiện.
\item (Sai) Có thể bỏ qua nối đất và thiết bị bảo vệ khi dùng điện áp cao.
\item (Đúng) Khi giải cần xác định đúng dữ kiện, phương chiều, điều kiện áp dụng và đơn vị.
\end{itemchoice}
}
\end{ex}

% ===== Câu 259 =====
% ID: L12C3B18-60-TL
% Mức: VDC
\dangbt{Tính điện áp, số vòng dây và cường độ dòng điện của máy biến áp lí tưởng}
\begin{bt}
% ID: L12C3B18-60-TL
% Mức: VDC
So sánh nhận định đúng “Máy biến áp lí tưởng thỏa U2/U1 = N2/N1 và U1I1 = U2I2.” với nhận định sai “Máy biến áp lí tưởng luôn có U2 = U1 bất kể số vòng dây.”; chỉ rõ căn cứ Vật lí.
\loigiai{
Cần nêu được: Máy biến áp lí tưởng thỏa U2/U1 = N2/N1 và U1I1 = U2I2. Khi vận dụng phải xác định đúng dữ kiện, phương chiều nếu có, điều kiện áp dụng, hệ đơn vị và kiểm tra tính hợp lí. Nhận định “Máy biến áp lí tưởng luôn có U2 = U1 bất kể số vòng dây.” sai.
}
\end{bt}

% ===== Câu 260 =====
% ID: L12C3B18-60-TLN
% Mức: VDC
\begin{ex}
% ID: L12C3B18-60-TLN
% Mức: VDC
Trong bốn phát biểu về Tính điện áp, số vòng dây và cường độ dòng điện của máy biến áp lí tưởng, có bao nhiêu phát biểu đúng? (Chỉ ghi một số nguyên.) (1) Khi giải cần xác định đúng dữ kiện, phương chiều, điều kiện áp dụng và đơn vị. (2) Máy biến áp lí tưởng luôn có U2 = U1 bất kể số vòng dây. (3) Phải dùng đúng định luật cho tình huống. (4) Máy biến áp lí tưởng thỏa U2/U1 = N2/N1 và U1I1 = U2I2.
\shortans{3}
\loigiai{
Đối chiếu với kiến thức: Máy biến áp lí tưởng thỏa U2/U1 = N2/N1 và U1I1 = U2I2. Có 3 phát biểu đúng.
}
\end{ex}

% ===== Câu 261 =====
% ID: L12C3B18-60-TN
% Mức: VD
\begin{ex}
% ID: L12C3B18-60-TN
% Mức: VD
Kết luận đúng khi vận dụng Tính điện áp, số vòng dây và cường độ dòng điện của máy biến áp lí tưởng?
\choice
{\True Máy biến áp lí tưởng thỏa U2/U1 = N2/N1 và U1I1 = U2I2.}
{Máy biến áp lí tưởng luôn có U2 = U1 bất kể số vòng dây.}
{Có thể bỏ qua phương, chiều và đơn vị trong mọi trường hợp.}
{Kết quả không phụ thuộc dữ kiện và điều kiện áp dụng.}
\loigiai{
Máy biến áp lí tưởng thỏa U2/U1 = N2/N1 và U1I1 = U2I2. Vì vậy phương án $A$ đúng; các phương án còn lại sai.
}
\end{ex}

% ===== Câu 262 =====
% ID: L12C3B18-61-TL
% Mức: TH
\dangbt{Truyền tải điện năng đi xa và giảm hao phí trên đường dây}
\begin{bt}
% ID: L12C3B18-61-TL
% Mức: TH
Trình bày kiến thức cốt lõi của Truyền tải điện năng đi xa và giảm hao phí trên đường dây và nêu điều kiện áp dụng.
\loigiai{
Cần nêu được: Với công suất truyền không đổi, tăng điện áp truyền tải làm giảm dòng điện và giảm hao phí I^2R. Khi vận dụng phải xác định đúng dữ kiện, phương chiều nếu có, điều kiện áp dụng, hệ đơn vị và kiểm tra tính hợp lí. Nhận định “Muốn giảm hao phí phải giảm điện áp truyền tải.” sai.
}
\end{bt}

% ===== Câu 263 =====
% ID: L12C3B18-61-TLN
% Mức: TH
\begin{ex}
% ID: L12C3B18-61-TLN
% Mức: TH
Trong bốn phát biểu về Truyền tải điện năng đi xa và giảm hao phí trên đường dây, có bao nhiêu phát biểu đúng? (Chỉ ghi một số nguyên.) (1) Với công suất truyền không đổi, tăng điện áp truyền tải làm giảm dòng điện và giảm hao phí I^2R. (2) Muốn giảm hao phí phải giảm điện áp truyền tải. (3) Cần kiểm tra điều kiện áp dụng. (4) Có thể bỏ qua đơn vị.
\shortans{2}
\loigiai{
Đối chiếu với kiến thức: Với công suất truyền không đổi, tăng điện áp truyền tải làm giảm dòng điện và giảm hao phí I^2R. Có 2 phát biểu đúng.
}
\end{ex}

% ===== Câu 264 =====
% ID: L12C3B18-61-TN
% Mức: NB
\dangbt{Nhận biết dòng điện Foucault và các ứng dụng của cảm ứng điện từ}
\begin{ex}
% ID: L12C3B18-61-TN
% Mức: NB
Phát biểu nào sau đây đúng về Nhận biết dòng điện Foucault và các ứng dụng của cảm ứng điện từ?
\choice
{Dòng điện Foucault chỉ xuất hiện trong chất cách điện.}
{Có thể bỏ qua phương, chiều và đơn vị trong mọi trường hợp.}
{Kết quả không phụ thuộc dữ kiện và điều kiện áp dụng.}
{\True Dòng điện Foucault là dòng điện cảm ứng xuất hiện trong khối vật dẫn khi từ thông qua nó biến thiên.}
\loigiai{
Dòng điện Foucault là dòng điện cảm ứng xuất hiện trong khối vật dẫn khi từ thông qua nó biến thiên. Vì vậy phương án $D$ đúng; các phương án còn lại sai.
}
\end{ex}

% ===== Câu 265 =====
% ID: L12C3B18-62-TL
% Mức: TH
\dangbt{Truyền tải điện năng đi xa và giảm hao phí trên đường dây}
\begin{bt}
% ID: L12C3B18-62-TL
% Mức: TH
Giải thích vì sao nhận định sau đúng: “Với công suất truyền không đổi, tăng điện áp truyền tải làm giảm dòng điện và giảm hao phí I^2R.”
\loigiai{
Cần nêu được: Với công suất truyền không đổi, tăng điện áp truyền tải làm giảm dòng điện và giảm hao phí I^2R. Khi vận dụng phải xác định đúng dữ kiện, phương chiều nếu có, điều kiện áp dụng, hệ đơn vị và kiểm tra tính hợp lí. Nhận định “Muốn giảm hao phí phải giảm điện áp truyền tải.” sai.
}
\end{bt}

% ===== Câu 266 =====
% ID: L12C3B18-62-TLN
% Mức: TH
\begin{ex}
% ID: L12C3B18-62-TLN
% Mức: TH
Trong bốn phát biểu về Truyền tải điện năng đi xa và giảm hao phí trên đường dây, có bao nhiêu phát biểu đúng? (Chỉ ghi một số nguyên.) (1) Muốn giảm hao phí phải giảm điện áp truyền tải. (2) Với công suất truyền không đổi, tăng điện áp truyền tải làm giảm dòng điện và giảm hao phí I^2R. (3) Kết quả phải phù hợp ý nghĩa vật lí. (4) Mọi đại lượng đều là vô hướng.
\shortans{1}
\loigiai{
Đối chiếu với kiến thức: Với công suất truyền không đổi, tăng điện áp truyền tải làm giảm dòng điện và giảm hao phí I^2R. Có 1 phát biểu đúng.
}
\end{ex}

% ===== Câu 267 =====
% ID: L12C3B18-62-TN
% Mức: NB
\dangbt{Nhận biết dòng điện Foucault và các ứng dụng của cảm ứng điện từ}
\begin{ex}
% ID: L12C3B18-62-TN
% Mức: NB
Trong nội dung Nhận biết dòng điện Foucault và các ứng dụng của cảm ứng điện từ?
\choice
{\True Dòng điện Foucault là dòng điện cảm ứng xuất hiện trong khối vật dẫn khi từ thông qua nó biến thiên.}
{Dòng điện Foucault chỉ xuất hiện trong chất cách điện.}
{Có thể bỏ qua phương, chiều và đơn vị trong mọi trường hợp.}
{Kết quả không phụ thuộc dữ kiện và điều kiện áp dụng.}
\loigiai{
Dòng điện Foucault là dòng điện cảm ứng xuất hiện trong khối vật dẫn khi từ thông qua nó biến thiên. Vì vậy phương án $A$ đúng; các phương án còn lại sai.
}
\end{ex}

% ===== Câu 268 =====
% ID: L12C3B18-63-TL
% Mức: VD
\dangbt{Truyền tải điện năng đi xa và giảm hao phí trên đường dây}
\begin{bt}
% ID: L12C3B18-63-TL
% Mức: VD
Phân tích sai lầm trong nhận định: “Muốn giảm hao phí phải giảm điện áp truyền tải.”
\loigiai{
Cần nêu được: Với công suất truyền không đổi, tăng điện áp truyền tải làm giảm dòng điện và giảm hao phí I^2R. Khi vận dụng phải xác định đúng dữ kiện, phương chiều nếu có, điều kiện áp dụng, hệ đơn vị và kiểm tra tính hợp lí. Nhận định “Muốn giảm hao phí phải giảm điện áp truyền tải.” sai.
}
\end{bt}

% ===== Câu 269 =====
% ID: L12C3B18-63-TLN
% Mức: TH
\begin{ex}
% ID: L12C3B18-63-TLN
% Mức: TH
Trong bốn phát biểu về Truyền tải điện năng đi xa và giảm hao phí trên đường dây, có bao nhiêu phát biểu đúng? (Chỉ ghi một số nguyên.) (1) Với công suất truyền không đổi, tăng điện áp truyền tải làm giảm dòng điện và giảm hao phí I^2R. (2) Muốn giảm hao phí phải giảm điện áp truyền tải. (3) Cần kiểm tra điều kiện áp dụng. (4) Có thể bỏ qua đơn vị.
\shortans{2}
\loigiai{
Đối chiếu với kiến thức: Với công suất truyền không đổi, tăng điện áp truyền tải làm giảm dòng điện và giảm hao phí I^2R. Có 2 phát biểu đúng.
}
\end{ex}

% ===== Câu 270 =====
% ID: L12C3B18-63-TN
% Mức: NB
\dangbt{Nhận biết dòng điện Foucault và các ứng dụng của cảm ứng điện từ}
\begin{ex}
% ID: L12C3B18-63-TN
% Mức: NB
Tình huống 3: chọn kết luận chính xác về Nhận biết dòng điện Foucault và các ứng dụng của cảm ứng điện từ?
\choice
{Kết quả không phụ thuộc dữ kiện và điều kiện áp dụng.}
{\True Dòng điện Foucault là dòng điện cảm ứng xuất hiện trong khối vật dẫn khi từ thông qua nó biến thiên.}
{Dòng điện Foucault chỉ xuất hiện trong chất cách điện.}
{Có thể bỏ qua phương, chiều và đơn vị trong mọi trường hợp.}
\loigiai{
Dòng điện Foucault là dòng điện cảm ứng xuất hiện trong khối vật dẫn khi từ thông qua nó biến thiên. Vì vậy phương án $B$ đúng; các phương án còn lại sai.
}
\end{ex}

% ===== Câu 271 =====
% ID: L12C3B18-64-TL
% Mức: VD
\dangbt{Truyền tải điện năng đi xa và giảm hao phí trên đường dây}
\begin{bt}
% ID: L12C3B18-64-TL
% Mức: VD
Đề xuất các bước giải một bài toán thuộc dạng Truyền tải điện năng đi xa và giảm hao phí trên đường dây.
\loigiai{
Cần nêu được: Với công suất truyền không đổi, tăng điện áp truyền tải làm giảm dòng điện và giảm hao phí I^2R. Khi vận dụng phải xác định đúng dữ kiện, phương chiều nếu có, điều kiện áp dụng, hệ đơn vị và kiểm tra tính hợp lí. Nhận định “Muốn giảm hao phí phải giảm điện áp truyền tải.” sai.
}
\end{bt}

% ===== Câu 272 =====
% ID: L12C3B18-64-TLN
% Mức: VD
\begin{ex}
% ID: L12C3B18-64-TLN
% Mức: VD
Trong bốn phát biểu về Truyền tải điện năng đi xa và giảm hao phí trên đường dây, có bao nhiêu phát biểu đúng? (Chỉ ghi một số nguyên.) (1) Muốn giảm hao phí phải giảm điện áp truyền tải. (2) Với công suất truyền không đổi, tăng điện áp truyền tải làm giảm dòng điện và giảm hao phí I^2R. (3) Kết quả phải phù hợp ý nghĩa vật lí. (4) Mọi đại lượng đều là vô hướng.
\shortans{2}
\loigiai{
Đối chiếu với kiến thức: Với công suất truyền không đổi, tăng điện áp truyền tải làm giảm dòng điện và giảm hao phí I^2R. Có 2 phát biểu đúng.
}
\end{ex}

% ===== Câu 273 =====
% ID: L12C3B18-64-TN
% Mức: TH
\dangbt{Nhận biết dòng điện Foucault và các ứng dụng của cảm ứng điện từ}
\begin{ex}
% ID: L12C3B18-64-TN
% Mức: TH
Nội dung nào mô tả đúng nhất Nhận biết dòng điện Foucault và các ứng dụng của cảm ứng điện từ?
\choice
{Có thể bỏ qua phương, chiều và đơn vị trong mọi trường hợp.}
{Kết quả không phụ thuộc dữ kiện và điều kiện áp dụng.}
{\True Dòng điện Foucault là dòng điện cảm ứng xuất hiện trong khối vật dẫn khi từ thông qua nó biến thiên.}
{Dòng điện Foucault chỉ xuất hiện trong chất cách điện.}
\loigiai{
Dòng điện Foucault là dòng điện cảm ứng xuất hiện trong khối vật dẫn khi từ thông qua nó biến thiên. Vì vậy phương án $C$ đúng; các phương án còn lại sai.
}
\end{ex}

% ===== Câu 274 =====
% ID: L12C3B18-65-TL
% Mức: VD
\dangbt{Truyền tải điện năng đi xa và giảm hao phí trên đường dây}
\begin{bt}
% ID: L12C3B18-65-TL
% Mức: VD
Nêu một tình huống thực tiễn liên quan đến Truyền tải điện năng đi xa và giảm hao phí trên đường dây và giải thích bằng kiến thức Vật lí.
\loigiai{
Cần nêu được: Với công suất truyền không đổi, tăng điện áp truyền tải làm giảm dòng điện và giảm hao phí I^2R. Khi vận dụng phải xác định đúng dữ kiện, phương chiều nếu có, điều kiện áp dụng, hệ đơn vị và kiểm tra tính hợp lí. Nhận định “Muốn giảm hao phí phải giảm điện áp truyền tải.” sai.
}
\end{bt}

% ===== Câu 275 =====
% ID: L12C3B18-65-TLN
% Mức: VD
\begin{ex}
% ID: L12C3B18-65-TLN
% Mức: VD
Trong bốn phát biểu về Truyền tải điện năng đi xa và giảm hao phí trên đường dây, có bao nhiêu phát biểu đúng? (Chỉ ghi một số nguyên.) (1) Với công suất truyền không đổi, tăng điện áp truyền tải làm giảm dòng điện và giảm hao phí I^2R. (2) Muốn giảm hao phí phải giảm điện áp truyền tải. (3) Cần kiểm tra điều kiện áp dụng. (4) Có thể bỏ qua đơn vị.
\shortans{2}
\loigiai{
Đối chiếu với kiến thức: Với công suất truyền không đổi, tăng điện áp truyền tải làm giảm dòng điện và giảm hao phí I^2R. Có 2 phát biểu đúng.
}
\end{ex}

% ===== Câu 276 =====
% ID: L12C3B18-65-TN
% Mức: TH
\dangbt{Nhận biết dòng điện Foucault và các ứng dụng của cảm ứng điện từ}
\begin{ex}
% ID: L12C3B18-65-TN
% Mức: TH
Khi phân tích Nhận biết dòng điện Foucault và các ứng dụng của cảm ứng điện từ?
\choice
{Dòng điện Foucault chỉ xuất hiện trong chất cách điện.}
{Có thể bỏ qua phương, chiều và đơn vị trong mọi trường hợp.}
{Kết quả không phụ thuộc dữ kiện và điều kiện áp dụng.}
{\True Dòng điện Foucault là dòng điện cảm ứng xuất hiện trong khối vật dẫn khi từ thông qua nó biến thiên.}
\loigiai{
Dòng điện Foucault là dòng điện cảm ứng xuất hiện trong khối vật dẫn khi từ thông qua nó biến thiên. Vì vậy phương án $D$ đúng; các phương án còn lại sai.
}
\end{ex}

% ===== Câu 277 =====
% ID: L12C3B18-66-TL
% Mức: VDC
\dangbt{Truyền tải điện năng đi xa và giảm hao phí trên đường dây}
\begin{bt}
% ID: L12C3B18-66-TL
% Mức: VDC
So sánh nhận định đúng “Với công suất truyền không đổi, tăng điện áp truyền tải làm giảm dòng điện và giảm hao phí I^2R.” với nhận định sai “Muốn giảm hao phí phải giảm điện áp truyền tải.”; chỉ rõ căn cứ Vật lí.
\loigiai{
Cần nêu được: Với công suất truyền không đổi, tăng điện áp truyền tải làm giảm dòng điện và giảm hao phí I^2R. Khi vận dụng phải xác định đúng dữ kiện, phương chiều nếu có, điều kiện áp dụng, hệ đơn vị và kiểm tra tính hợp lí. Nhận định “Muốn giảm hao phí phải giảm điện áp truyền tải.” sai.
}
\end{bt}

% ===== Câu 278 =====
% ID: L12C3B18-66-TLN
% Mức: VDC
\begin{ex}
% ID: L12C3B18-66-TLN
% Mức: VDC
Trong bốn phát biểu về Truyền tải điện năng đi xa và giảm hao phí trên đường dây, có bao nhiêu phát biểu đúng? (Chỉ ghi một số nguyên.) (1) Khi giải cần xác định đúng dữ kiện, phương chiều, điều kiện áp dụng và đơn vị. (2) Muốn giảm hao phí phải giảm điện áp truyền tải. (3) Phải dùng đúng định luật cho tình huống. (4) Với công suất truyền không đổi, tăng điện áp truyền tải làm giảm dòng điện và giảm hao phí I^2R.
\shortans{3}
\loigiai{
Đối chiếu với kiến thức: Với công suất truyền không đổi, tăng điện áp truyền tải làm giảm dòng điện và giảm hao phí I^2R. Có 3 phát biểu đúng.
}
\end{ex}

% ===== Câu 279 =====
% ID: L12C3B18-66-TN
% Mức: TH
\dangbt{Nhận biết dòng điện Foucault và các ứng dụng của cảm ứng điện từ}
\begin{ex}
% ID: L12C3B18-66-TN
% Mức: TH
Một học sinh ôn tập Nhận biết dòng điện Foucault và các ứng dụng của cảm ứng điện từ?
\choice
{\True Dòng điện Foucault là dòng điện cảm ứng xuất hiện trong khối vật dẫn khi từ thông qua nó biến thiên.}
{Dòng điện Foucault chỉ xuất hiện trong chất cách điện.}
{Có thể bỏ qua phương, chiều và đơn vị trong mọi trường hợp.}
{Kết quả không phụ thuộc dữ kiện và điều kiện áp dụng.}
\loigiai{
Dòng điện Foucault là dòng điện cảm ứng xuất hiện trong khối vật dẫn khi từ thông qua nó biến thiên. Vì vậy phương án $A$ đúng; các phương án còn lại sai.
}
\end{ex}

% ===== Câu 280 =====
% ID: L12C3B18-67-TL
% Mức: TH
\dangbt{Bài toán tổng hợp về máy biến áp, truyền tải điện và an toàn điện}
\begin{bt}
% ID: L12C3B18-67-TL
% Mức: TH
Trình bày kiến thức cốt lõi của Bài toán tổng hợp về máy biến áp, truyền tải điện và an toàn điện và nêu điều kiện áp dụng.
\loigiai{
Cần nêu được: An toàn điện yêu cầu cách điện, nối đất phù hợp và không chạm phần dẫn điện đang có điện áp nguy hiểm. Khi vận dụng phải xác định đúng dữ kiện, phương chiều nếu có, điều kiện áp dụng, hệ đơn vị và kiểm tra tính hợp lí. Nhận định “Có thể bỏ qua nối đất và thiết bị bảo vệ khi dùng điện áp cao.” sai.
}
\end{bt}

% ===== Câu 281 =====
% ID: L12C3B18-67-TLN
% Mức: TH
\begin{ex}
% ID: L12C3B18-67-TLN
% Mức: TH
Trong bốn phát biểu về Bài toán tổng hợp về máy biến áp, truyền tải điện và an toàn điện, có bao nhiêu phát biểu đúng? (Chỉ ghi một số nguyên.) (1) An toàn điện yêu cầu cách điện, nối đất phù hợp và không chạm phần dẫn điện đang có điện áp nguy hiểm. (2) Có thể bỏ qua nối đất và thiết bị bảo vệ khi dùng điện áp cao. (3) Cần kiểm tra điều kiện áp dụng. (4) Có thể bỏ qua đơn vị.
\shortans{2}
\loigiai{
Đối chiếu với kiến thức: An toàn điện yêu cầu cách điện, nối đất phù hợp và không chạm phần dẫn điện đang có điện áp nguy hiểm. Có 2 phát biểu đúng.
}
\end{ex}

% ===== Câu 282 =====
% ID: L12C3B18-67-TN
% Mức: TH
\dangbt{Nhận biết dòng điện Foucault và các ứng dụng của cảm ứng điện từ}
\begin{ex}
% ID: L12C3B18-67-TN
% Mức: TH
Chọn phát biểu không trái với kiến thức về Nhận biết dòng điện Foucault và các ứng dụng của cảm ứng điện từ?
\choice
{Kết quả không phụ thuộc dữ kiện và điều kiện áp dụng.}
{\True Dòng điện Foucault là dòng điện cảm ứng xuất hiện trong khối vật dẫn khi từ thông qua nó biến thiên.}
{Dòng điện Foucault chỉ xuất hiện trong chất cách điện.}
{Có thể bỏ qua phương, chiều và đơn vị trong mọi trường hợp.}
\loigiai{
Dòng điện Foucault là dòng điện cảm ứng xuất hiện trong khối vật dẫn khi từ thông qua nó biến thiên. Vì vậy phương án $B$ đúng; các phương án còn lại sai.
}
\end{ex}

% ===== Câu 283 =====
% ID: L12C3B18-68-TL
% Mức: TH
\dangbt{Bài toán tổng hợp về máy biến áp, truyền tải điện và an toàn điện}
\begin{bt}
% ID: L12C3B18-68-TL
% Mức: TH
Giải thích vì sao nhận định sau đúng: “An toàn điện yêu cầu cách điện, nối đất phù hợp và không chạm phần dẫn điện đang có điện áp nguy hiểm.”
\loigiai{
Cần nêu được: An toàn điện yêu cầu cách điện, nối đất phù hợp và không chạm phần dẫn điện đang có điện áp nguy hiểm. Khi vận dụng phải xác định đúng dữ kiện, phương chiều nếu có, điều kiện áp dụng, hệ đơn vị và kiểm tra tính hợp lí. Nhận định “Có thể bỏ qua nối đất và thiết bị bảo vệ khi dùng điện áp cao.” sai.
}
\end{bt}

% ===== Câu 284 =====
% ID: L12C3B18-68-TLN
% Mức: TH
\begin{ex}
% ID: L12C3B18-68-TLN
% Mức: TH
Trong bốn phát biểu về Bài toán tổng hợp về máy biến áp, truyền tải điện và an toàn điện, có bao nhiêu phát biểu đúng? (Chỉ ghi một số nguyên.) (1) Có thể bỏ qua nối đất và thiết bị bảo vệ khi dùng điện áp cao. (2) An toàn điện yêu cầu cách điện, nối đất phù hợp và không chạm phần dẫn điện đang có điện áp nguy hiểm. (3) Kết quả phải phù hợp ý nghĩa vật lí. (4) Mọi đại lượng đều là vô hướng.
\shortans{1}
\loigiai{
Đối chiếu với kiến thức: An toàn điện yêu cầu cách điện, nối đất phù hợp và không chạm phần dẫn điện đang có điện áp nguy hiểm. Có 1 phát biểu đúng.
}
\end{ex}

% ===== Câu 285 =====
% ID: L12C3B18-68-TN
% Mức: VD
\dangbt{Nhận biết dòng điện Foucault và các ứng dụng của cảm ứng điện từ}
\begin{ex}
% ID: L12C3B18-68-TN
% Mức: VD
Cơ sở Vật lí đúng dùng để giải Nhận biết dòng điện Foucault và các ứng dụng của cảm ứng điện từ?
\choice
{Có thể bỏ qua phương, chiều và đơn vị trong mọi trường hợp.}
{Kết quả không phụ thuộc dữ kiện và điều kiện áp dụng.}
{\True Dòng điện Foucault là dòng điện cảm ứng xuất hiện trong khối vật dẫn khi từ thông qua nó biến thiên.}
{Dòng điện Foucault chỉ xuất hiện trong chất cách điện.}
\loigiai{
Dòng điện Foucault là dòng điện cảm ứng xuất hiện trong khối vật dẫn khi từ thông qua nó biến thiên. Vì vậy phương án $C$ đúng; các phương án còn lại sai.
}
\end{ex}

% ===== Câu 286 =====
% ID: L12C3B18-69-TL
% Mức: VD
\dangbt{Bài toán tổng hợp về máy biến áp, truyền tải điện và an toàn điện}
\begin{bt}
% ID: L12C3B18-69-TL
% Mức: VD
Phân tích sai lầm trong nhận định: “Có thể bỏ qua nối đất và thiết bị bảo vệ khi dùng điện áp cao.”
\loigiai{
Cần nêu được: An toàn điện yêu cầu cách điện, nối đất phù hợp và không chạm phần dẫn điện đang có điện áp nguy hiểm. Khi vận dụng phải xác định đúng dữ kiện, phương chiều nếu có, điều kiện áp dụng, hệ đơn vị và kiểm tra tính hợp lí. Nhận định “Có thể bỏ qua nối đất và thiết bị bảo vệ khi dùng điện áp cao.” sai.
}
\end{bt}

% ===== Câu 287 =====
% ID: L12C3B18-69-TLN
% Mức: TH
\begin{ex}
% ID: L12C3B18-69-TLN
% Mức: TH
Trong bốn phát biểu về Bài toán tổng hợp về máy biến áp, truyền tải điện và an toàn điện, có bao nhiêu phát biểu đúng? (Chỉ ghi một số nguyên.) (1) An toàn điện yêu cầu cách điện, nối đất phù hợp và không chạm phần dẫn điện đang có điện áp nguy hiểm. (2) Có thể bỏ qua nối đất và thiết bị bảo vệ khi dùng điện áp cao. (3) Cần kiểm tra điều kiện áp dụng. (4) Có thể bỏ qua đơn vị.
\shortans{2}
\loigiai{
Đối chiếu với kiến thức: An toàn điện yêu cầu cách điện, nối đất phù hợp và không chạm phần dẫn điện đang có điện áp nguy hiểm. Có 2 phát biểu đúng.
}
\end{ex}

% ===== Câu 288 =====
% ID: L12C3B18-69-TN
% Mức: VD
\dangbt{Nhận biết dòng điện Foucault và các ứng dụng của cảm ứng điện từ}
\begin{ex}
% ID: L12C3B18-69-TN
% Mức: VD
Trong một bài thực hành về Nhận biết dòng điện Foucault và các ứng dụng của cảm ứng điện từ?
\choice
{Dòng điện Foucault chỉ xuất hiện trong chất cách điện.}
{Có thể bỏ qua phương, chiều và đơn vị trong mọi trường hợp.}
{Kết quả không phụ thuộc dữ kiện và điều kiện áp dụng.}
{\True Dòng điện Foucault là dòng điện cảm ứng xuất hiện trong khối vật dẫn khi từ thông qua nó biến thiên.}
\loigiai{
Dòng điện Foucault là dòng điện cảm ứng xuất hiện trong khối vật dẫn khi từ thông qua nó biến thiên. Vì vậy phương án $D$ đúng; các phương án còn lại sai.
}
\end{ex}

% ===== Câu 289 =====
% ID: L12C3B18-70-TL
% Mức: VD
\dangbt{Bài toán tổng hợp về máy biến áp, truyền tải điện và an toàn điện}
\begin{bt}
% ID: L12C3B18-70-TL
% Mức: VD
Đề xuất các bước giải một bài toán thuộc dạng Bài toán tổng hợp về máy biến áp, truyền tải điện và an toàn điện.
\loigiai{
Cần nêu được: An toàn điện yêu cầu cách điện, nối đất phù hợp và không chạm phần dẫn điện đang có điện áp nguy hiểm. Khi vận dụng phải xác định đúng dữ kiện, phương chiều nếu có, điều kiện áp dụng, hệ đơn vị và kiểm tra tính hợp lí. Nhận định “Có thể bỏ qua nối đất và thiết bị bảo vệ khi dùng điện áp cao.” sai.
}
\end{bt}

% ===== Câu 290 =====
% ID: L12C3B18-70-TLN
% Mức: VD
\begin{ex}
% ID: L12C3B18-70-TLN
% Mức: VD
Trong bốn phát biểu về Bài toán tổng hợp về máy biến áp, truyền tải điện và an toàn điện, có bao nhiêu phát biểu đúng? (Chỉ ghi một số nguyên.) (1) Có thể bỏ qua nối đất và thiết bị bảo vệ khi dùng điện áp cao. (2) An toàn điện yêu cầu cách điện, nối đất phù hợp và không chạm phần dẫn điện đang có điện áp nguy hiểm. (3) Kết quả phải phù hợp ý nghĩa vật lí. (4) Mọi đại lượng đều là vô hướng.
\shortans{2}
\loigiai{
Đối chiếu với kiến thức: An toàn điện yêu cầu cách điện, nối đất phù hợp và không chạm phần dẫn điện đang có điện áp nguy hiểm. Có 2 phát biểu đúng.
}
\end{ex}

% ===== Câu 291 =====
% ID: L12C3B18-70-TN
% Mức: VD
\dangbt{Nhận biết dòng điện Foucault và các ứng dụng của cảm ứng điện từ}
\begin{ex}
% ID: L12C3B18-70-TN
% Mức: VD
Kết luận đúng khi vận dụng Nhận biết dòng điện Foucault và các ứng dụng của cảm ứng điện từ?
\choice
{\True Dòng điện Foucault là dòng điện cảm ứng xuất hiện trong khối vật dẫn khi từ thông qua nó biến thiên.}
{Dòng điện Foucault chỉ xuất hiện trong chất cách điện.}
{Có thể bỏ qua phương, chiều và đơn vị trong mọi trường hợp.}
{Kết quả không phụ thuộc dữ kiện và điều kiện áp dụng.}
\loigiai{
Dòng điện Foucault là dòng điện cảm ứng xuất hiện trong khối vật dẫn khi từ thông qua nó biến thiên. Vì vậy phương án $A$ đúng; các phương án còn lại sai.
}
\end{ex}

% ===== Câu 292 =====
% ID: L12C3B18-71-TL
% Mức: VD
\dangbt{Bài toán tổng hợp về máy biến áp, truyền tải điện và an toàn điện}
\begin{bt}
% ID: L12C3B18-71-TL
% Mức: VD
Nêu một tình huống thực tiễn liên quan đến Bài toán tổng hợp về máy biến áp, truyền tải điện và an toàn điện và giải thích bằng kiến thức Vật lí.
\loigiai{
Cần nêu được: An toàn điện yêu cầu cách điện, nối đất phù hợp và không chạm phần dẫn điện đang có điện áp nguy hiểm. Khi vận dụng phải xác định đúng dữ kiện, phương chiều nếu có, điều kiện áp dụng, hệ đơn vị và kiểm tra tính hợp lí. Nhận định “Có thể bỏ qua nối đất và thiết bị bảo vệ khi dùng điện áp cao.” sai.
}
\end{bt}

% ===== Câu 293 =====
% ID: L12C3B18-71-TLN
% Mức: VD
\begin{ex}
% ID: L12C3B18-71-TLN
% Mức: VD
Trong bốn phát biểu về Bài toán tổng hợp về máy biến áp, truyền tải điện và an toàn điện, có bao nhiêu phát biểu đúng? (Chỉ ghi một số nguyên.) (1) An toàn điện yêu cầu cách điện, nối đất phù hợp và không chạm phần dẫn điện đang có điện áp nguy hiểm. (2) Có thể bỏ qua nối đất và thiết bị bảo vệ khi dùng điện áp cao. (3) Cần kiểm tra điều kiện áp dụng. (4) Có thể bỏ qua đơn vị.
\shortans{2}
\loigiai{
Đối chiếu với kiến thức: An toàn điện yêu cầu cách điện, nối đất phù hợp và không chạm phần dẫn điện đang có điện áp nguy hiểm. Có 2 phát biểu đúng.
}
\end{ex}

% ===== Câu 294 =====
% ID: L12C3B18-71-TN
% Mức: NB
\dangbt{Giải thích nguyên lí bếp từ, phanh điện từ và gia nhiệt cảm ứng}
\begin{ex}
% ID: L12C3B18-71-TN
% Mức: NB
Phát biểu nào sau đây đúng về Giải thích nguyên lí bếp từ, phanh điện từ và gia nhiệt cảm ứng?
\choice
{Bếp từ làm nóng trực tiếp mọi vật liệu bằng điện trường tĩnh.}
{Có thể bỏ qua phương, chiều và đơn vị trong mọi trường hợp.}
{Kết quả không phụ thuộc dữ kiện và điều kiện áp dụng.}
{\True Bếp từ làm nóng đáy nồi dẫn điện nhờ dòng điện Foucault và sự tỏa nhiệt Joule.}
\loigiai{
Bếp từ làm nóng đáy nồi dẫn điện nhờ dòng điện Foucault và sự tỏa nhiệt Joule. Vì vậy phương án $D$ đúng; các phương án còn lại sai.
}
\end{ex}

% ===== Câu 295 =====
% ID: L12C3B18-72-TL
% Mức: VDC
\dangbt{Bài toán tổng hợp về máy biến áp, truyền tải điện và an toàn điện}
\begin{bt}
% ID: L12C3B18-72-TL
% Mức: VDC
So sánh nhận định đúng “An toàn điện yêu cầu cách điện, nối đất phù hợp và không chạm phần dẫn điện đang có điện áp nguy hiểm.” với nhận định sai “Có thể bỏ qua nối đất và thiết bị bảo vệ khi dùng điện áp cao.”; chỉ rõ căn cứ Vật lí.
\loigiai{
Cần nêu được: An toàn điện yêu cầu cách điện, nối đất phù hợp và không chạm phần dẫn điện đang có điện áp nguy hiểm. Khi vận dụng phải xác định đúng dữ kiện, phương chiều nếu có, điều kiện áp dụng, hệ đơn vị và kiểm tra tính hợp lí. Nhận định “Có thể bỏ qua nối đất và thiết bị bảo vệ khi dùng điện áp cao.” sai.
}
\end{bt}

% ===== Câu 296 =====
% ID: L12C3B18-72-TLN
% Mức: VDC
\begin{ex}
% ID: L12C3B18-72-TLN
% Mức: VDC
Trong bốn phát biểu về Bài toán tổng hợp về máy biến áp, truyền tải điện và an toàn điện, có bao nhiêu phát biểu đúng? (Chỉ ghi một số nguyên.) (1) Khi giải cần xác định đúng dữ kiện, phương chiều, điều kiện áp dụng và đơn vị. (2) Có thể bỏ qua nối đất và thiết bị bảo vệ khi dùng điện áp cao. (3) Phải dùng đúng định luật cho tình huống. (4) An toàn điện yêu cầu cách điện, nối đất phù hợp và không chạm phần dẫn điện đang có điện áp nguy hiểm.
\shortans{3}
\loigiai{
Đối chiếu với kiến thức: An toàn điện yêu cầu cách điện, nối đất phù hợp và không chạm phần dẫn điện đang có điện áp nguy hiểm. Có 3 phát biểu đúng.
}
\end{ex}

% ===== Câu 297 =====
% ID: L12C3B18-72-TN
% Mức: NB
\dangbt{Giải thích nguyên lí bếp từ, phanh điện từ và gia nhiệt cảm ứng}
\begin{ex}
% ID: L12C3B18-72-TN
% Mức: NB
Trong nội dung Giải thích nguyên lí bếp từ, phanh điện từ và gia nhiệt cảm ứng?
\choice
{\True Bếp từ làm nóng đáy nồi dẫn điện nhờ dòng điện Foucault và sự tỏa nhiệt Joule.}
{Bếp từ làm nóng trực tiếp mọi vật liệu bằng điện trường tĩnh.}
{Có thể bỏ qua phương, chiều và đơn vị trong mọi trường hợp.}
{Kết quả không phụ thuộc dữ kiện và điều kiện áp dụng.}
\loigiai{
Bếp từ làm nóng đáy nồi dẫn điện nhờ dòng điện Foucault và sự tỏa nhiệt Joule. Vì vậy phương án $A$ đúng; các phương án còn lại sai.
}
\end{ex}

% ===== Câu 298 =====
% ID: L12C3B18-73-TN
% Mức: NB
\begin{ex}
% ID: L12C3B18-73-TN
% Mức: NB
Tình huống 3: chọn kết luận chính xác về Giải thích nguyên lí bếp từ, phanh điện từ và gia nhiệt cảm ứng?
\choice
{Kết quả không phụ thuộc dữ kiện và điều kiện áp dụng.}
{\True Bếp từ làm nóng đáy nồi dẫn điện nhờ dòng điện Foucault và sự tỏa nhiệt Joule.}
{Bếp từ làm nóng trực tiếp mọi vật liệu bằng điện trường tĩnh.}
{Có thể bỏ qua phương, chiều và đơn vị trong mọi trường hợp.}
\loigiai{
Bếp từ làm nóng đáy nồi dẫn điện nhờ dòng điện Foucault và sự tỏa nhiệt Joule. Vì vậy phương án $B$ đúng; các phương án còn lại sai.
}
\end{ex}

% ===== Câu 299 =====
% ID: L12C3B18-74-TN
% Mức: TH
\begin{ex}
% ID: L12C3B18-74-TN
% Mức: TH
Nội dung nào mô tả đúng nhất Giải thích nguyên lí bếp từ, phanh điện từ và gia nhiệt cảm ứng?
\choice
{Có thể bỏ qua phương, chiều và đơn vị trong mọi trường hợp.}
{Kết quả không phụ thuộc dữ kiện và điều kiện áp dụng.}
{\True Bếp từ làm nóng đáy nồi dẫn điện nhờ dòng điện Foucault và sự tỏa nhiệt Joule.}
{Bếp từ làm nóng trực tiếp mọi vật liệu bằng điện trường tĩnh.}
\loigiai{
Bếp từ làm nóng đáy nồi dẫn điện nhờ dòng điện Foucault và sự tỏa nhiệt Joule. Vì vậy phương án $C$ đúng; các phương án còn lại sai.
}
\end{ex}

% ===== Câu 300 =====
% ID: L12C3B18-75-TN
% Mức: TH
\begin{ex}
% ID: L12C3B18-75-TN
% Mức: TH
Khi phân tích Giải thích nguyên lí bếp từ, phanh điện từ và gia nhiệt cảm ứng?
\choice
{Bếp từ làm nóng trực tiếp mọi vật liệu bằng điện trường tĩnh.}
{Có thể bỏ qua phương, chiều và đơn vị trong mọi trường hợp.}
{Kết quả không phụ thuộc dữ kiện và điều kiện áp dụng.}
{\True Bếp từ làm nóng đáy nồi dẫn điện nhờ dòng điện Foucault và sự tỏa nhiệt Joule.}
\loigiai{
Bếp từ làm nóng đáy nồi dẫn điện nhờ dòng điện Foucault và sự tỏa nhiệt Joule. Vì vậy phương án $D$ đúng; các phương án còn lại sai.
}
\end{ex}

% ===== Câu 301 =====
% ID: L12C3B18-76-TN
% Mức: TH
\begin{ex}
% ID: L12C3B18-76-TN
% Mức: TH
Một học sinh ôn tập Giải thích nguyên lí bếp từ, phanh điện từ và gia nhiệt cảm ứng?
\choice
{\True Bếp từ làm nóng đáy nồi dẫn điện nhờ dòng điện Foucault và sự tỏa nhiệt Joule.}
{Bếp từ làm nóng trực tiếp mọi vật liệu bằng điện trường tĩnh.}
{Có thể bỏ qua phương, chiều và đơn vị trong mọi trường hợp.}
{Kết quả không phụ thuộc dữ kiện và điều kiện áp dụng.}
\loigiai{
Bếp từ làm nóng đáy nồi dẫn điện nhờ dòng điện Foucault và sự tỏa nhiệt Joule. Vì vậy phương án $A$ đúng; các phương án còn lại sai.
}
\end{ex}

% ===== Câu 302 =====
% ID: L12C3B18-77-TN
% Mức: TH
\begin{ex}
% ID: L12C3B18-77-TN
% Mức: TH
Chọn phát biểu không trái với kiến thức về Giải thích nguyên lí bếp từ, phanh điện từ và gia nhiệt cảm ứng?
\choice
{Kết quả không phụ thuộc dữ kiện và điều kiện áp dụng.}
{\True Bếp từ làm nóng đáy nồi dẫn điện nhờ dòng điện Foucault và sự tỏa nhiệt Joule.}
{Bếp từ làm nóng trực tiếp mọi vật liệu bằng điện trường tĩnh.}
{Có thể bỏ qua phương, chiều và đơn vị trong mọi trường hợp.}
\loigiai{
Bếp từ làm nóng đáy nồi dẫn điện nhờ dòng điện Foucault và sự tỏa nhiệt Joule. Vì vậy phương án $B$ đúng; các phương án còn lại sai.
}
\end{ex}

% ===== Câu 303 =====
% ID: L12C3B18-78-TN
% Mức: VD
\begin{ex}
% ID: L12C3B18-78-TN
% Mức: VD
Cơ sở Vật lí đúng dùng để giải Giải thích nguyên lí bếp từ, phanh điện từ và gia nhiệt cảm ứng?
\choice
{Có thể bỏ qua phương, chiều và đơn vị trong mọi trường hợp.}
{Kết quả không phụ thuộc dữ kiện và điều kiện áp dụng.}
{\True Bếp từ làm nóng đáy nồi dẫn điện nhờ dòng điện Foucault và sự tỏa nhiệt Joule.}
{Bếp từ làm nóng trực tiếp mọi vật liệu bằng điện trường tĩnh.}
\loigiai{
Bếp từ làm nóng đáy nồi dẫn điện nhờ dòng điện Foucault và sự tỏa nhiệt Joule. Vì vậy phương án $C$ đúng; các phương án còn lại sai.
}
\end{ex}

% ===== Câu 304 =====
% ID: L12C3B18-79-TN
% Mức: VD
\begin{ex}
% ID: L12C3B18-79-TN
% Mức: VD
Trong một bài thực hành về Giải thích nguyên lí bếp từ, phanh điện từ và gia nhiệt cảm ứng?
\choice
{Bếp từ làm nóng trực tiếp mọi vật liệu bằng điện trường tĩnh.}
{Có thể bỏ qua phương, chiều và đơn vị trong mọi trường hợp.}
{Kết quả không phụ thuộc dữ kiện và điều kiện áp dụng.}
{\True Bếp từ làm nóng đáy nồi dẫn điện nhờ dòng điện Foucault và sự tỏa nhiệt Joule.}
\loigiai{
Bếp từ làm nóng đáy nồi dẫn điện nhờ dòng điện Foucault và sự tỏa nhiệt Joule. Vì vậy phương án $D$ đúng; các phương án còn lại sai.
}
\end{ex}

% ===== Câu 305 =====
% ID: L12C3B18-80-TN
% Mức: VD
\begin{ex}
% ID: L12C3B18-80-TN
% Mức: VD
Kết luận đúng khi vận dụng Giải thích nguyên lí bếp từ, phanh điện từ và gia nhiệt cảm ứng?
\choice
{\True Bếp từ làm nóng đáy nồi dẫn điện nhờ dòng điện Foucault và sự tỏa nhiệt Joule.}
{Bếp từ làm nóng trực tiếp mọi vật liệu bằng điện trường tĩnh.}
{Có thể bỏ qua phương, chiều và đơn vị trong mọi trường hợp.}
{Kết quả không phụ thuộc dữ kiện và điều kiện áp dụng.}
\loigiai{
Bếp từ làm nóng đáy nồi dẫn điện nhờ dòng điện Foucault và sự tỏa nhiệt Joule. Vì vậy phương án $A$ đúng; các phương án còn lại sai.
}
\end{ex}

% ===== Câu 306 =====
% ID: L12C3B18-81-TN
% Mức: NB
\dangbt{Nhận biết cấu tạo, nguyên lí hoạt động của máy biến áp}
\begin{ex}
% ID: L12C3B18-81-TN
% Mức: NB
Phát biểu nào sau đây đúng về Nhận biết cấu tạo, nguyên lí hoạt động của máy biến áp?
\choice
{Máy biến áp lí tưởng hoạt động được với nguồn điện một chiều không đổi.}
{Có thể bỏ qua phương, chiều và đơn vị trong mọi trường hợp.}
{Kết quả không phụ thuộc dữ kiện và điều kiện áp dụng.}
{\True Máy biến áp gồm lõi sắt và hai cuộn dây, hoạt động dựa trên cảm ứng điện từ với dòng điện xoay chiều.}
\loigiai{
Máy biến áp gồm lõi sắt và hai cuộn dây, hoạt động dựa trên cảm ứng điện từ với dòng điện xoay chiều. Vì vậy phương án $D$ đúng; các phương án còn lại sai.
}
\end{ex}

% ===== Câu 307 =====
% ID: L12C3B18-82-TN
% Mức: NB
\begin{ex}
% ID: L12C3B18-82-TN
% Mức: NB
Trong nội dung Nhận biết cấu tạo, nguyên lí hoạt động của máy biến áp?
\choice
{\True Máy biến áp gồm lõi sắt và hai cuộn dây, hoạt động dựa trên cảm ứng điện từ với dòng điện xoay chiều.}
{Máy biến áp lí tưởng hoạt động được với nguồn điện một chiều không đổi.}
{Có thể bỏ qua phương, chiều và đơn vị trong mọi trường hợp.}
{Kết quả không phụ thuộc dữ kiện và điều kiện áp dụng.}
\loigiai{
Máy biến áp gồm lõi sắt và hai cuộn dây, hoạt động dựa trên cảm ứng điện từ với dòng điện xoay chiều. Vì vậy phương án $A$ đúng; các phương án còn lại sai.
}
\end{ex}

% ===== Câu 308 =====
% ID: L12C3B18-83-TN
% Mức: NB
\begin{ex}
% ID: L12C3B18-83-TN
% Mức: NB
Tình huống 3: chọn kết luận chính xác về Nhận biết cấu tạo, nguyên lí hoạt động của máy biến áp?
\choice
{Kết quả không phụ thuộc dữ kiện và điều kiện áp dụng.}
{\True Máy biến áp gồm lõi sắt và hai cuộn dây, hoạt động dựa trên cảm ứng điện từ với dòng điện xoay chiều.}
{Máy biến áp lí tưởng hoạt động được với nguồn điện một chiều không đổi.}
{Có thể bỏ qua phương, chiều và đơn vị trong mọi trường hợp.}
\loigiai{
Máy biến áp gồm lõi sắt và hai cuộn dây, hoạt động dựa trên cảm ứng điện từ với dòng điện xoay chiều. Vì vậy phương án $B$ đúng; các phương án còn lại sai.
}
\end{ex}

% ===== Câu 309 =====
% ID: L12C3B18-84-TN
% Mức: TH
\begin{ex}
% ID: L12C3B18-84-TN
% Mức: TH
Nội dung nào mô tả đúng nhất Nhận biết cấu tạo, nguyên lí hoạt động của máy biến áp?
\choice
{Có thể bỏ qua phương, chiều và đơn vị trong mọi trường hợp.}
{Kết quả không phụ thuộc dữ kiện và điều kiện áp dụng.}
{\True Máy biến áp gồm lõi sắt và hai cuộn dây, hoạt động dựa trên cảm ứng điện từ với dòng điện xoay chiều.}
{Máy biến áp lí tưởng hoạt động được với nguồn điện một chiều không đổi.}
\loigiai{
Máy biến áp gồm lõi sắt và hai cuộn dây, hoạt động dựa trên cảm ứng điện từ với dòng điện xoay chiều. Vì vậy phương án $C$ đúng; các phương án còn lại sai.
}
\end{ex}

% ===== Câu 310 =====
% ID: L12C3B18-85-TN
% Mức: TH
\begin{ex}
% ID: L12C3B18-85-TN
% Mức: TH
Khi phân tích Nhận biết cấu tạo, nguyên lí hoạt động của máy biến áp?
\choice
{Máy biến áp lí tưởng hoạt động được với nguồn điện một chiều không đổi.}
{Có thể bỏ qua phương, chiều và đơn vị trong mọi trường hợp.}
{Kết quả không phụ thuộc dữ kiện và điều kiện áp dụng.}
{\True Máy biến áp gồm lõi sắt và hai cuộn dây, hoạt động dựa trên cảm ứng điện từ với dòng điện xoay chiều.}
\loigiai{
Máy biến áp gồm lõi sắt và hai cuộn dây, hoạt động dựa trên cảm ứng điện từ với dòng điện xoay chiều. Vì vậy phương án $D$ đúng; các phương án còn lại sai.
}
\end{ex}

% ===== Câu 311 =====
% ID: L12C3B18-86-TN
% Mức: TH
\begin{ex}
% ID: L12C3B18-86-TN
% Mức: TH
Một học sinh ôn tập Nhận biết cấu tạo, nguyên lí hoạt động của máy biến áp?
\choice
{\True Máy biến áp gồm lõi sắt và hai cuộn dây, hoạt động dựa trên cảm ứng điện từ với dòng điện xoay chiều.}
{Máy biến áp lí tưởng hoạt động được với nguồn điện một chiều không đổi.}
{Có thể bỏ qua phương, chiều và đơn vị trong mọi trường hợp.}
{Kết quả không phụ thuộc dữ kiện và điều kiện áp dụng.}
\loigiai{
Máy biến áp gồm lõi sắt và hai cuộn dây, hoạt động dựa trên cảm ứng điện từ với dòng điện xoay chiều. Vì vậy phương án $A$ đúng; các phương án còn lại sai.
}
\end{ex}

% ===== Câu 312 =====
% ID: L12C3B18-87-TN
% Mức: TH
\begin{ex}
% ID: L12C3B18-87-TN
% Mức: TH
Chọn phát biểu không trái với kiến thức về Nhận biết cấu tạo, nguyên lí hoạt động của máy biến áp?
\choice
{Kết quả không phụ thuộc dữ kiện và điều kiện áp dụng.}
{\True Máy biến áp gồm lõi sắt và hai cuộn dây, hoạt động dựa trên cảm ứng điện từ với dòng điện xoay chiều.}
{Máy biến áp lí tưởng hoạt động được với nguồn điện một chiều không đổi.}
{Có thể bỏ qua phương, chiều và đơn vị trong mọi trường hợp.}
\loigiai{
Máy biến áp gồm lõi sắt và hai cuộn dây, hoạt động dựa trên cảm ứng điện từ với dòng điện xoay chiều. Vì vậy phương án $B$ đúng; các phương án còn lại sai.
}
\end{ex}

% ===== Câu 313 =====
% ID: L12C3B18-88-TN
% Mức: VD
\begin{ex}
% ID: L12C3B18-88-TN
% Mức: VD
Cơ sở Vật lí đúng dùng để giải Nhận biết cấu tạo, nguyên lí hoạt động của máy biến áp?
\choice
{Có thể bỏ qua phương, chiều và đơn vị trong mọi trường hợp.}
{Kết quả không phụ thuộc dữ kiện và điều kiện áp dụng.}
{\True Máy biến áp gồm lõi sắt và hai cuộn dây, hoạt động dựa trên cảm ứng điện từ với dòng điện xoay chiều.}
{Máy biến áp lí tưởng hoạt động được với nguồn điện một chiều không đổi.}
\loigiai{
Máy biến áp gồm lõi sắt và hai cuộn dây, hoạt động dựa trên cảm ứng điện từ với dòng điện xoay chiều. Vì vậy phương án $C$ đúng; các phương án còn lại sai.
}
\end{ex}

% ===== Câu 314 =====
% ID: L12C3B18-89-TN
% Mức: VD
\begin{ex}
% ID: L12C3B18-89-TN
% Mức: VD
Trong một bài thực hành về Nhận biết cấu tạo, nguyên lí hoạt động của máy biến áp?
\choice
{Máy biến áp lí tưởng hoạt động được với nguồn điện một chiều không đổi.}
{Có thể bỏ qua phương, chiều và đơn vị trong mọi trường hợp.}
{Kết quả không phụ thuộc dữ kiện và điều kiện áp dụng.}
{\True Máy biến áp gồm lõi sắt và hai cuộn dây, hoạt động dựa trên cảm ứng điện từ với dòng điện xoay chiều.}
\loigiai{
Máy biến áp gồm lõi sắt và hai cuộn dây, hoạt động dựa trên cảm ứng điện từ với dòng điện xoay chiều. Vì vậy phương án $D$ đúng; các phương án còn lại sai.
}
\end{ex}

% ===== Câu 315 =====
% ID: L12C3B18-90-TN
% Mức: VD
\begin{ex}
% ID: L12C3B18-90-TN
% Mức: VD
Kết luận đúng khi vận dụng Nhận biết cấu tạo, nguyên lí hoạt động của máy biến áp?
\choice
{\True Máy biến áp gồm lõi sắt và hai cuộn dây, hoạt động dựa trên cảm ứng điện từ với dòng điện xoay chiều.}
{Máy biến áp lí tưởng hoạt động được với nguồn điện một chiều không đổi.}
{Có thể bỏ qua phương, chiều và đơn vị trong mọi trường hợp.}
{Kết quả không phụ thuộc dữ kiện và điều kiện áp dụng.}
\loigiai{
Máy biến áp gồm lõi sắt và hai cuộn dây, hoạt động dựa trên cảm ứng điện từ với dòng điện xoay chiều. Vì vậy phương án $A$ đúng; các phương án còn lại sai.
}
\end{ex}

% ===== Câu 316 =====
% ID: L12C3B18-91-TN
% Mức: NB
\dangbt{Tính điện áp, số vòng dây và cường độ dòng điện của máy biến áp lí tưởng}
\begin{ex}
% ID: L12C3B18-91-TN
% Mức: NB
Phát biểu nào sau đây đúng về Tính điện áp, số vòng dây và cường độ dòng điện của máy biến áp lí tưởng?
\choice
{Máy biến áp lí tưởng luôn có U2 = U1 bất kể số vòng dây.}
{Có thể bỏ qua phương, chiều và đơn vị trong mọi trường hợp.}
{Kết quả không phụ thuộc dữ kiện và điều kiện áp dụng.}
{\True Máy biến áp lí tưởng thỏa U2/U1 = N2/N1 và U1I1 = U2I2.}
\loigiai{
Máy biến áp lí tưởng thỏa U2/U1 = N2/N1 và U1I1 = U2I2. Vì vậy phương án $D$ đúng; các phương án còn lại sai.
}
\end{ex}

% ===== Câu 317 =====
% ID: L12C3B18-92-TN
% Mức: NB
\begin{ex}
% ID: L12C3B18-92-TN
% Mức: NB
Trong nội dung Tính điện áp, số vòng dây và cường độ dòng điện của máy biến áp lí tưởng?
\choice
{\True Máy biến áp lí tưởng thỏa U2/U1 = N2/N1 và U1I1 = U2I2.}
{Máy biến áp lí tưởng luôn có U2 = U1 bất kể số vòng dây.}
{Có thể bỏ qua phương, chiều và đơn vị trong mọi trường hợp.}
{Kết quả không phụ thuộc dữ kiện và điều kiện áp dụng.}
\loigiai{
Máy biến áp lí tưởng thỏa U2/U1 = N2/N1 và U1I1 = U2I2. Vì vậy phương án $A$ đúng; các phương án còn lại sai.
}
\end{ex}

% ===== Câu 318 =====
% ID: L12C3B18-93-TN
% Mức: NB
\begin{ex}
% ID: L12C3B18-93-TN
% Mức: NB
Tình huống 3: chọn kết luận chính xác về Tính điện áp, số vòng dây và cường độ dòng điện của máy biến áp lí tưởng?
\choice
{Kết quả không phụ thuộc dữ kiện và điều kiện áp dụng.}
{\True Máy biến áp lí tưởng thỏa U2/U1 = N2/N1 và U1I1 = U2I2.}
{Máy biến áp lí tưởng luôn có U2 = U1 bất kể số vòng dây.}
{Có thể bỏ qua phương, chiều và đơn vị trong mọi trường hợp.}
\loigiai{
Máy biến áp lí tưởng thỏa U2/U1 = N2/N1 và U1I1 = U2I2. Vì vậy phương án $B$ đúng; các phương án còn lại sai.
}
\end{ex}

% ===== Câu 319 =====
% ID: L12C3B18-94-TN
% Mức: TH
\begin{ex}
% ID: L12C3B18-94-TN
% Mức: TH
Nội dung nào mô tả đúng nhất Tính điện áp, số vòng dây và cường độ dòng điện của máy biến áp lí tưởng?
\choice
{Có thể bỏ qua phương, chiều và đơn vị trong mọi trường hợp.}
{Kết quả không phụ thuộc dữ kiện và điều kiện áp dụng.}
{\True Máy biến áp lí tưởng thỏa U2/U1 = N2/N1 và U1I1 = U2I2.}
{Máy biến áp lí tưởng luôn có U2 = U1 bất kể số vòng dây.}
\loigiai{
Máy biến áp lí tưởng thỏa U2/U1 = N2/N1 và U1I1 = U2I2. Vì vậy phương án $C$ đúng; các phương án còn lại sai.
}
\end{ex}

% ===== Câu 320 =====
% ID: L12C3B18-95-TN
% Mức: TH
\begin{ex}
% ID: L12C3B18-95-TN
% Mức: TH
Khi phân tích Tính điện áp, số vòng dây và cường độ dòng điện của máy biến áp lí tưởng?
\choice
{Máy biến áp lí tưởng luôn có U2 = U1 bất kể số vòng dây.}
{Có thể bỏ qua phương, chiều và đơn vị trong mọi trường hợp.}
{Kết quả không phụ thuộc dữ kiện và điều kiện áp dụng.}
{\True Máy biến áp lí tưởng thỏa U2/U1 = N2/N1 và U1I1 = U2I2.}
\loigiai{
Máy biến áp lí tưởng thỏa U2/U1 = N2/N1 và U1I1 = U2I2. Vì vậy phương án $D$ đúng; các phương án còn lại sai.
}
\end{ex}

% ===== Câu 321 =====
% ID: L12C3B18-96-TN
% Mức: TH
\begin{ex}
% ID: L12C3B18-96-TN
% Mức: TH
Một học sinh ôn tập Tính điện áp, số vòng dây và cường độ dòng điện của máy biến áp lí tưởng?
\choice
{\True Máy biến áp lí tưởng thỏa U2/U1 = N2/N1 và U1I1 = U2I2.}
{Máy biến áp lí tưởng luôn có U2 = U1 bất kể số vòng dây.}
{Có thể bỏ qua phương, chiều và đơn vị trong mọi trường hợp.}
{Kết quả không phụ thuộc dữ kiện và điều kiện áp dụng.}
\loigiai{
Máy biến áp lí tưởng thỏa U2/U1 = N2/N1 và U1I1 = U2I2. Vì vậy phương án $A$ đúng; các phương án còn lại sai.
}
\end{ex}

% ===== Câu 322 =====
% ID: L12C3B18-97-TN
% Mức: TH
\begin{ex}
% ID: L12C3B18-97-TN
% Mức: TH
Chọn phát biểu không trái với kiến thức về Tính điện áp, số vòng dây và cường độ dòng điện của máy biến áp lí tưởng?
\choice
{Kết quả không phụ thuộc dữ kiện và điều kiện áp dụng.}
{\True Máy biến áp lí tưởng thỏa U2/U1 = N2/N1 và U1I1 = U2I2.}
{Máy biến áp lí tưởng luôn có U2 = U1 bất kể số vòng dây.}
{Có thể bỏ qua phương, chiều và đơn vị trong mọi trường hợp.}
\loigiai{
Máy biến áp lí tưởng thỏa U2/U1 = N2/N1 và U1I1 = U2I2. Vì vậy phương án $B$ đúng; các phương án còn lại sai.
}
\end{ex}

% ===== Câu 323 =====
% ID: L12C3B18-98-TN
% Mức: VD
\begin{ex}
% ID: L12C3B18-98-TN
% Mức: VD
Cơ sở Vật lí đúng dùng để giải Tính điện áp, số vòng dây và cường độ dòng điện của máy biến áp lí tưởng?
\choice
{Có thể bỏ qua phương, chiều và đơn vị trong mọi trường hợp.}
{Kết quả không phụ thuộc dữ kiện và điều kiện áp dụng.}
{\True Máy biến áp lí tưởng thỏa U2/U1 = N2/N1 và U1I1 = U2I2.}
{Máy biến áp lí tưởng luôn có U2 = U1 bất kể số vòng dây.}
\loigiai{
Máy biến áp lí tưởng thỏa U2/U1 = N2/N1 và U1I1 = U2I2. Vì vậy phương án $C$ đúng; các phương án còn lại sai.
}
\end{ex}

% ===== Câu 324 =====
% ID: L12C3B18-99-TN
% Mức: VD
\begin{ex}
% ID: L12C3B18-99-TN
% Mức: VD
Trong một bài thực hành về Tính điện áp, số vòng dây và cường độ dòng điện của máy biến áp lí tưởng?
\choice
{Máy biến áp lí tưởng luôn có U2 = U1 bất kể số vòng dây.}
{Có thể bỏ qua phương, chiều và đơn vị trong mọi trường hợp.}
{Kết quả không phụ thuộc dữ kiện và điều kiện áp dụng.}
{\True Máy biến áp lí tưởng thỏa U2/U1 = N2/N1 và U1I1 = U2I2.}
\loigiai{
Máy biến áp lí tưởng thỏa U2/U1 = N2/N1 và U1I1 = U2I2. Vì vậy phương án $D$ đúng; các phương án còn lại sai.
}
\end{ex}
